\begin{proof}
\noindent\textbf{Original source, complete new filling and directed physical transfer.}\par\smallskip
The source, full Miller and theta coefficients, point/O modules and
period dictionary are the retained SI2, I2P and I2M maps below.  Supporting
proof 20 constructs a complete auxiliary atlas directly from their
original convex period cells.  Its 1,926-cone common fan and all three
endpoint atlases retain the entire original source, physical circle,
period functions, selected native star, scalar root and outer marking.

RB retains the whole fixed-source scalar reference.  NS and NC give
the actual protected native maps and the directed Hopf cap with its
full normal derivative.  The retained MN reference and IL source,
Cartier, cube-frame and cap maps are used with the new moving and
minimal lower stars.  RFB carries the entire reference fan change;
ASI2 retains compactness, the original finite-cap square, the exterior
cylinder, all four parameter edges and the physical inverse clutch.

The complete new compact pair is $(895,139)$, the twelve lower doubles
sum to $(12,-24)$, and the actual moving-minus-own-reference pair is
$(-906,-111)$.  Their sum is $(1,4)$.  The entire positive reference
change is $(72,0)$, and the explicit Jin interval reversal gives
$(71,-4)$.  The unchanged physical line and held quaternion give
$(\mu,\nu)=(-1,72)$ and the original oriented sphere class $12$.

The earlier auxiliary execution is retained separately in
\cref{iv:open:K111:original_auxiliary_arrays}, with its missing ordered
arrays and replay requirement stated there.  The construction and
numerical parts used here are exactly the new complete atlas and the
retained original geometric maps.

\par\bigskip\noindent\textbf{P3K.111 supporting proof 1: Original source, M2 Miller/point periods, full cubic cap and parameter square}\par\nobreak
\par

\medskip\noindent\textbf{\(M_{2}\) selected-I3: complete source, point-line, good-root and period maps}

\par\nobreak

Working exact construction, 18 November 2026. The original source is
\[
 y^2=x^3+rx^2+tx+t^2,\quad G=(0,t),\quad s=1/4-r.
\]
We use the actual point line \(M_2=\mathcal O([2G]-[O])\), the
chronological translation by \(-3G\), the ordered selected node
\(t_c(s)\to0\), cap \((m,n)=(3,1)\), its prescribed finite character,
and a fixed original scalar \(\lambda_0=r_0^3\) in the common strict
contraction range. The positive original \(s\)-circle and the
unit-\(+1\) Tate marking specify the initial fan. This note constructs
the source/line/root/annular maps and exact period parents; it assigns
no characteristic pair or return.

\par

\medskip\noindent\textbf{SI2.1. Original selected height and the full node morphism}

\par\nobreak

Use the original marked-node cover near \(w=\pm1\):
\[
 v=w^2-1,\quad D=v-2,\quad r=(1-2v)/D^2,\quad
 (X,Y,T)=(D^2x,D^3y,D^3t).
\]
The inverse divides by the same nonzero functions. Its exact equation,
selected zero and source parameter are
\[
 Y^2=X^3+(1-2v)X^2+DTX+T^2,\quad q=T-v=D^3q_{\rm old},
\]
\[
 q_{\rm old}=t-t_c(s),\quad
 t_c(s(v))=v/D^3,\quad s=v(v+4)/(4D^2),\quad s'(0)=1/4.
 \tag{SI2.1}
\]
Thus each of these marked parameter discs carries one positive
original source meridian.

Put \(s_0=v+v^2/4\),
\(U=Y-T-DX/2,\ V_s=Y+T+DX/2\). The literal source identity is
\(UV_s=X^2(X-s_0)\). On its ordered source-resolution chart
\[
 U=Xz,\quad V_s=X\omega,\quad X=s_0+z\omega,\quad
 Y=X(z+\omega)/2,\quad T=X(\omega-z-D)/2.
 \tag{SI2.2}
\]
Its selected node is \(z=-v/2,\omega=v/2\). Set
\(a=z+v/2,\ b=\omega-v/2\) and retain the analytic square root
\[
 R_b=\sqrt{(b+2)^2+4v},\quad R_b(0)=2,\quad K_b=b+v+2+R_b.
\]
The full source map is
\[
 x_n=a-vb/K_b,\quad
 y_n=[bK_b-(2b+v)a]/4,\qquad x_ny_n=q.
 \tag{SI2.3}
\]
Its inverse is the unique small solution
\[
 b=\frac{4y_n+vx_n}{2(R_b(b,v)-x_n)},\qquad
 a=x_n+vb/K_b.
 \tag{SI2.4}
\]
The denominator is a unit. On a small fixed ball the first right
side has derivative of norm less than \(1/2\) in \(b\); successive
substitution gives its holomorphic inverse and the displayed
substitution verifies both compositions. The normal Jacobian is
\(\sqrt{1+v}\) with value one. This keeps the full selected height
and source parameter, including both source branches.

The original inverse Weierstrass/Tate map E81, with its unit-\(+1\)
root, supplies the actual functions
\[
 Q=P^2R,\quad P=t u_P,\quad R-P=s u_s,\quad
 u_P(0)=2,\quad u_s(0)=1,\quad t=P u_t.
\]
The complete discriminant comparison gives
\[
 Q=t^2q_{\rm old}u_Q,\quad u_Q(0)=8,\quad
 q_{\rm old}=R c_q,\quad c_q=u_P^2/u_Q,\quad c_q(0)=1/2.
 \tag{SI2.5}
\]
These are the original analytic source functions, not just their
orders. In particular \(R=0\) is exactly the selected-height divisor.

\par

\medskip\noindent\textbf{SI2.2. Six Miller factors and the genuine \(M_{2}\) point frame}

\par\nobreak

Let \(T_G(p)=p-G\). The literal point-to-tensor coefficient is
\[
 f_2=\frac{y-t-x/2}{x-s},\qquad r_2=f_2^{-1}.
\]
The full chronological suspension on the genuine point line is
\[
 B_{2,3}(p)=f_2(p)f_2(T_G^3p)
                 f_2(T_Gp)^2f_2(T_G^2p)^2,\qquad
 B_{\rm sel}^{(2,3)}=(q_{\rm old}/t^3)B_{2,3}.
 \tag{SI2.6}
\]
The source scaling in SI2.1 gives
\[
 f_{2,\rm old}=D^{-1}f_{2,\rm new},\quad
 r_{2,\rm old}=D r_{2,\rm new},\quad
 B_{2,3,\rm old}=D^{-6}B_{2,3,\rm new}.
\]
Therefore \(B_{\rm sel}^{(2,3)}=q B_{2,3,\rm new}/T^3\)
exactly. The canonical meromorphic point-section coefficients are
identified by the source isomorphism. Its balanced tensor coefficient
has the additional \(D\) factor written below; that factor is retained.

Write \(f_P(Z)=\Theta_Q(Z/P)/\Theta_Q(Z)\) and retain the full
degree-one coefficient
\[
 \eta=a_T\,\Theta'_Q(1)\,\Theta_Q(P^2)/\Theta_Q(P)^2.
\]
Here \(a_T\) is the actual Weierstrass differential scale, with \(a_T(0)=1\); the separate derivative has \(\Theta'_Q(1)(0)=-1\), hence \(\eta(0)=-1\). The derivative factor remains in every point/tensor arrow and its inverse. In particular the full sixth power is \(\eta^6=a_T^6\Theta'_Q(1)^6\Theta_Q(P^2)^6/\Theta_Q(P)^{12}\). The formulas below use the actual \(\eta\) fixed by the displayed Miller identity; this derivative factor is present in their full analytic coefficients, as well as in the first-power point-frame comparison. The original chord identity and exact telescope are
\[
 \eta=f_2(p)\frac{f_P(P^{-1}Z)}{f_P(Z)},\qquad
 f_P(Z)f_P(P^{-1}Z)=f_{P^2}(Z).
 \tag{SI2.7}
\]
They give the full module arrows
\[
 V_{2,T}=Z^{-1}f_P(Z)^2r_2(p)v_2,\quad
 V_{\rm reg}=ZV_{2,T}=\eta^{-1}V_{{\rm pt},2},\quad
 V_{{\rm pt},2}=f_{P^2}(Z)v_2.
 \tag{SI2.8}
\]
Under SI2.1, \(\eta_{\rm old}=D^{-1}\eta_{\rm new}\),
\(V_{{\rm pt},2,\rm old}=V_{{\rm pt},2,\rm new}\), and
\(V_{\rm reg,old}=D V_{\rm reg,new}\). This explicitly types the
point/tensor change and its inverse.

The six factors in SI2.6 telescope, with their order retained, to
\[
 B_{2,3}=\eta^6
       \frac{f_{P^2}(Z)}{f_{P^2}(P^{-3}Z)}.
 \tag{SI2.9}
\]
Consequently the complete source/line periods are
\[
 h:(Z,V_{2,T})\mapsto(P^2RZ,R^{-1}V_{2,T}),\quad
 B:(Z,V_{2,T})\mapsto(P^{-3}Z,\lambda q_{\rm old}K V_{2,T}),
\]
\[
 K=(P/t)^3\eta^6,\quad K(0)=8,
 \tag{SI2.10}
\]
or, in the actual point-related coefficient,
\[
 h:(V_{\rm reg},Z)\mapsto(P^2V_{\rm reg},P^2RZ),\quad
 B:(V_{\rm reg},Z)\mapsto
   (\lambda q_{\rm old}\eta^6t^{-3}V_{\rm reg},P^{-3}Z).
 \tag{SI2.11}
\]
All \(q_{\rm old}=Rc_q\), \(t=Pu_t\), and \(\eta\) factors remain.

At \(s=0\), set \(A_*=y-t-x/2,\ B_*=y+t+x/2\).
Then \(A_*B_*=x^3\) and the three chronological coefficients are
\(A_*/x,-B_*/x,2t/x\). Their product is \(-2t\), hence
\(B_{2,3}=4t^2\). The selected height is now exactly \(t\);
the genuine point Hom map is \(B_{\rm sel}^{(2,3)}=4\), covering
the identity \(T_G^3=1\). In the retained Tate lifts \(Q=P^3\),
so the actual central period is
\[
 hB:(Z,V_{\rm reg})\longmapsto(Z,4\lambda V_{\rm reg}).
 \tag{SI2.12}
\]
This is the full central coefficient, derived from the Miller
maps, not a chosen scalar normalization.

On the protected selected \(q_{\rm old}\)-disc over a nonzero
\(s\)-annulus, put
\[
 Q=q_{\rm old}\kappa,\quad \kappa=P^2/c_q,\quad
 \alpha=P^2,\quad g=P^{-3},\quad
 \eta_*=\lambda_{\rm base}\eta^6t^{-3}.
 \tag{SI2.13}
\]
Their exact parameter windings are \((2,2,-3,-3)\). The full
fixed-scalar winding matrix in the source/point coefficient marking is
\[
 \begin{pmatrix}2&-3\\2&-3\end{pmatrix}.
 \tag{SI2.14}
\]
No return or free coordinate is inferred from this matrix.

\par

\medskip\noindent\textbf{SI2.3. Point/O modules and both ordered markings}

\par\nobreak

For a retained source band \(j\), set
\[
 x=Q^jZ,\quad y=Q^{1-j}/Z,\quad xy=P^2R.
\]
The actual point graph is the blowup of \((x,P^2)\):
\[
 x=P^2a,\ ay=R;\qquad P^2=xc,\ y=Rc.
 \tag{SI2.15}
\]
The second chart is resolved by
\[
 P=uv,\ x=u,\ c=uv^2;\qquad P=uv,\ x=u^2v,\ c=v,
 \tag{SI2.16}
\]
with the literal overlap \(u_2=v_1^{-1},v_2=u_1v_1^2\).
The source functions \(s=(R-P)/u_s,\ t=Pu_t\) travel through these
proper graph maps. In the first chart \(ds\) has a nonzero \(dP\)
coefficient near the collision; in the two latter charts it has
a nonzero \(dR\) coefficient. Thus the actual source parameter is
preserved and the projection is a submersion on these small charts.

Put \(b=P^2,e=R\). The full theta product has
\[
 b f_b(x)=(b-x)E_{b,e}(x,y),
\quad
 E_{b,e}=
 \frac{\prod_{r\ge1}(1-Q^{r-1}ex)(1-Q^{r-1}by)}
 {(1-x)\prod_{r\ge1}(1-Q^rx)(1-Q^{r-1}y)}.
 \tag{SI2.17}
\]
On each chosen small node chart this product and its reciprocal
converge normally to units. In the actual point generators
\(e_a=(a-1)^{-1},e_c=(c-1)^{-1}\), writing
\(v_2=e_a A=e_c C\), the full arrows are
\[
 V_{\rm reg}=-\eta^{-1}P^{-2j}E_{b,e}A,\qquad
 V_{\rm reg}=\eta^{-1}P^{-2j}c^{-1}E_{b,e}C,\quad C=-cA.
 \tag{SI2.18}
\]
Their inverses retain every displayed factor. Near O use
\(v_2=(x-1)C_O\); then
\[
 V_{\rm reg}=\eta^{-1}P^{-2j-2}(x-P^2)E_OC_O,
 \tag{SI2.19}
\]
where
\[
 E_O=
 \frac{\prod_{r\ge1}(1-Q^rx/P^2)(1-Q^rP^2/x)}
 {\prod_{r\ge1}(1-Q^rx)(1-Q^r/x)} .
\]
In particular the additional \(P^{-2}\) at O remains in the
coefficient map. Under \(h\) use \(j'=j-1\); the local source and
module coefficients are unchanged. Under \(B\), for the actual
target band \(j'\),
\[
 x'=Q^{j'-j}P^{-3}x,\qquad
 v_2'=\lambda q_{\rm old}\eta^6t^{-3}P^{2(j'-j)}
                 \frac{f_b(x)}{f_b(x')}v_2.
 \tag{SI2.20}
\]
The reciprocal is evaluated at the inverse source point. This is
the full six-factor Hom by SI2.9.

The original unit-\(+1\) Tate source has
\(Z_T=Z_w\) near \(w=-1\) and \(Z_T=Z_w^{-1}\) near \(w=+1\).
The full theta inversion gives the point-line arrow
\[
 (Z,V_{{\rm pt},2})\mapsto(Z^{-1},P^2V_{{\rm pt},2}).
\]
The target point representative is \(P^{-1}\); its actual source
period is \((Q,P^{-2})\). Composing gives
\[
 (Z,V_{{\rm pt},2})\mapsto(Q/Z,V_{{\rm pt},2}).
 \tag{SI2.21}
\]
This is the literal reflected-fan character substitution, with both
point-module maps SI2.18--20 retained. The Weierstrass differential
changes sign with \(Z\)-inversion; the corresponding sign in
\(\Theta(P^{-2})/\Theta(P^{-1})^2\) leaves the full \(\eta\)
unchanged. The source/cap family itself is the original one.

\par

\medskip\noindent\textbf{SI2.4. The cubic good root and full annular reference}

\par\nobreak

Use \(q_\infty=z^3\) and the literal original-source arrow
\[
 t=z^{-3},\quad x=z^{-2}X,\quad y=z^{-3}Y,\qquad
 Y^2=X^3+r z^2X^2+zX+1 .
 \tag{SI2.22}
\]
The deck is \(D_3(z,X,Y)=(\zeta_3z,\zeta_3^{-4}X,Y)\).
The marked sections are \(G'=(0,1)\) and
\(2G'=(s z^2,-1-(s/2)z^3)\). The root source is a smooth proper
elliptic family near \(s=z=0\).
All six Miller factors give
\[
 f_{2,\rm old}=z^{-1}f'_2,\quad
 B_{2,3,\rm old}=z^{-6}B'_{2,3},\quad
 B'_{\rm sel}=h_\infty B'_{2,3},\quad
 h_\infty=1-t_c(s)z^3 .
 \tag{SI2.23}
\]
On \(z=0\), the three factors are
\((Y-1)/X,-(Y+1)/X,2/X\), of product \(-2\).
On \(s=0\), the selected height has \(t_c(0)=0\), so the complete
central Hom is \(B'_{\rm sel}=4\), on both parameter axes.

The original point-line correction is retained on the quotient
graph of weights \((w,z)=(2,1)\):
\[
 \beta^*\mathcal O(2G'-O)=\pi^*M_2+E_P-E_O .
 \tag{SI2.24}
\]
The quotient rays are \((2/3,1/3),(1/3,2/3)\); the correction is
the literal minimum \(\min(v_w,2v_z)\), with its monomial generator
at both marked sections. These are OG.6\textquotesingle s actual module maps.

Let \(R_0'\) be the unique small division \(3R_0'=3G'\). Its formal
coordinate is obtained from
\[
 \tau(3G')=z^2
 \frac{-2s(2+sz^3)}{8+6sz^3+(s^2-s^3)z^6};
\quad
 \tau(R_0')=s z^2\ell_R,\quad \ell_R(0,0)=-1/6.
 \tag{SI2.25}
\]
It and its inverse preserve the actual refinement ideals
\((w,z),(w,z^2),(z,w^2)\). Their canonical divisor sections carry
the complete correction SI2.24.
For \(P'=2G'\), the normalized chord is
\[
 F=\frac{\tau(R_0')\ell_{P',R_0'}}{X-X_{P'+R_0'}} .
 \tag{SI2.26}
\]
Writing \(X_{R_0'}=A/\tau^2,Y_{R_0'}=-A/\tau^3\), its numerator is
\[
 \tau(Y-Y_{P'})+
 \frac{A+\tau^3Y_{P'}}{A-\tau^2X_{P'}}(X-X_{P'}),
\]
whose denominator is a unit and whose exact Hom value is one at
\(\tau=0\). The horizontal divisor and the zero vertical orders
extend this Hom and its inverse jointly across both parameter axes.
The chord and \(\tau\) deck weights cancel, so \(F\) commutes with
\(D_3\).

Set \(\beta=B'_{\rm sel}/F^3\), with the full chronological product
understood. This is a nonzero base function of central value four.
Its prescribed invariant analytic cube root gives
\[
 \widetilde F=\beta^{1/3}F,\quad
 A_0=\zeta_3^d\widetilde F D_3,\quad
 A_\lambda=r_\lambda A_0,\quad
 B_\lambda=\lambda B'_{\rm sel},\quad r_\lambda^3=\lambda.
 \tag{SI2.27}
\]
Thus \(A_\lambda^3=B_\lambda\) and they commute, including
the full finite character and both point corrections.

On each original infinity sector let
\(\ell_\alpha=\log_\alpha q_\infty/(2\pi i)\),
\(z_\alpha=e^{2\pi i\ell_\alpha/3}\). The full original annular map is
\[
 \Phi_\alpha=
 \operatorname{Exp}\{\ell_\alpha\operatorname{Log}_\alpha\widetilde F+
 \ell_\alpha(\log r_\lambda/(2\pi i)+d/3)\delta+h_\alpha\}\,I_{z_\alpha}.
 \tag{SI2.28}
\]
Here \(I_{z_\alpha}\) includes SI2.24. The complete OG.12--13 sector
construction, with the displayed \(t_c(s)\) and \(M_{2}\) modules substituted,
gives
\(\Phi_\beta=A_\lambda^{j_0}B_\lambda^{j_{\alpha\beta}}\Phi_\alpha\)
and \(\Phi_\alpha B_\lambda=B_\lambda\Phi_\alpha\).
The reference-zero Cauchy operators use fixed annular contours;
the retained cochain contains the full source period and line lattice.
Its inverse reverses the sector, root and Lie arrows. Changing from
the distant selected node changes the normalized root by the actual
invariant cube root of
\((1-t_c(s)z^3)/(1-t_{\rm distant}(s)z^3)\).
This full unit and the original integer reference cochain remain
in SI2.28; neither is reset.

\par

\medskip\noindent\textbf{SI2.5. The admitted twelfth-power scalar and its complete phase square}

\par\nobreak

Choose
\[
 \lambda=\lambda_0s^{12},\quad r_\lambda=r_0s^4,\quad r_0^3=\lambda_0.
 \tag{SI2.29}
\]
For \(A=-\log|P|,\ C=-\log|R|,\ D_B=-\log|\lambda q_{\rm old}K|\),
the full period determinant is
\[
 (2A+C)D_B-3AC.
 \tag{SI2.30}
\]
The actual equation \(R-P=su_s\) gives
\(D_B\ge12\min(A,C)+C-C_0\), with bounded full-unit logarithms.
Since \(3AC/(2A+C)\le3\min(A,C)\), shrinking the bidisc gives a
strict positive determinant gap with every unit retained.

Use the independent root parameter \(r\), with
\(B=r^3B_0(s),A=rA_0(s)\), and the entire annular reference
\(F=F_0(s)+\rho\,\ell_\infty\delta\),
\(\rho=\log r/(2\pi i)\). Along SI2.29 the varying part is
\(4\tau\ell_\infty\delta\), \(\tau=\log s/(2\pi i)\).
The ordinary and cap displacements are
\[
 o_{\rm ord}=4\tau\frac{\log q_{\rm old}}{2\pi i}\delta,\qquad
 o_{\rm cap}=4\tau\frac{\log h_\infty}{2\pi i}\delta.
 \tag{SI2.31}
\]
The exact equality \(q_\infty q_{\rm old}=h_\infty\) gives
\(F+o_{\rm ord}-o_{\rm cap}=F_{\rm ref}\). The full forward square is
\[
 \Phi_{\rho+4}q_{\rm old}^{\,4}=h_\infty^4\Phi_\rho .
 \tag{SI2.32}
\]
Its endpoint-to-initial arrows are their full inverses.
With \(\beta_2=\eta^6t^{-3}\), the selected ordinary endpoint has
\[
 D_{\rm ord}^4=B^4R_4,\qquad
 R_4=(P^{12}Z,\lambda^{-4}\beta_2^{-4}V_{{\rm pt},2}).
 \tag{SI2.33}
\]
This retains the selected source translation and line factor at
the selected-height zero.

At a reference source point \(s_*\), the actual root basepoint is
\(r_{\rm base}=r_0s_*^4\). On the source cut put
\[
 G=\operatorname{Log}(P^{-3}),\quad
 J=\operatorname{Log}(\lambda_{\rm base}\beta_2),\qquad
 C_\varrho=(e^{\varrho G},
   e^{\varrho J+3\pi i\varrho(\varrho-1)}).
 \tag{SI2.34}
\]
The complete source-cut increments of both \(G\) and \(J\) are
\(-6\pi i\). With
\(\eta_\varrho=e^{6\pi i\varrho}\lambda_{\rm base}\beta_2\),
\[
 C_{\varrho+1}(g^{-1},\eta_\varrho^{-1})=C_\varrho,\quad
 C_0=1,\quad C_4=(g^4,\eta_*^4)=R_4^{-1}.
 \tag{SI2.35}
\]
For the native Cox lift
\((a_0,a_1,c_0,c_1)\mapsto(a_0,Aa_1,(AD)^{-1}c_0,Dc_1)\),
the source-cut mismatch is \((A,D)=(e^{-6\pi i\varrho},e^{-6\pi i\varrho})\).
Its literal quaternion map is
\[
 q_H\longmapsto e^{3\pi i\varrho}q_H e^{9\pi i\varrho}.
 \tag{SI2.36}
\]
At every integer seam its two half-phase signs are equal.
The base and orbit sphere rotate by the displayed left adjoint
map, which extends over the actual three-ball and carries the
principal physical line. The source-cut faces and their Spin
endpoints are therefore fully specified.

In the original DZ origin, substitution of the four root turns
in its actual quadratic product formula gives the retained factor
\[
 e^{3\pi i(16t^2+4t)}
   =e^{48\pi it(t-1)}e^{60\pi it}.
 \tag{SI2.37}
\]
The first logarithm has zero endpoint values; the second is the
physical loop of winding \(30\). It travels through both compared
boundary products, with their entire Cox and cap maps.
Thus the source/root parameter words are \((1,4)\), \((0,4)\)
and \((1,0)\), with the actual edge morphisms SI2.31--37 retained.
The open scalar path from \(r_{\rm base}\) to \(r_0\) is
\(r(a)=r_0\exp(4(1-a)\operatorname{Log}s_*)\), with its full lifted
log in SI2.28. The normal and compact central attachment has yet
to be evaluated, so these word identities assign no return.

\par

\medskip\noindent\textbf{SI2.6. Exact three endpoint period parents}

\par\nobreak

On \(R=pP\), put
\[
 c_{12}=\lambda_0c_q\eta^6/(u_s^{12}u_t^3),\quad
 c_{12}(0)=4\lambda_0,\quad f_0=(p-1)^{12}c_{12}.
 \tag{SI2.38}
\]
The original periods and the words \(K=hB,\ J=h\) are
\[
 h=(P^2V_{\rm reg},P^3pZ),\quad
 B=(P^{10}p f_0V_{\rm reg},P^{-3}Z),
\]
\[
 K=(P^{12}p f_0V_{\rm reg},pZ),\qquad
 J=(P^2V_{\rm reg},P^3pZ).
 \tag{SI2.39}
\]
In common characters \(X_c=V_{\rm reg}/Z,\ Z_c=Z\), take
\(L=K^2J,T=KJ\). Their entire arrows are
\[
 L=(P^{23}p^{-1}f_0^2X_c,P^3p^3Z_c),\qquad
 T=(P^{11}p^{-1}f_0X_c,P^3p^2Z_c).
 \tag{SI2.40}
\]
The generic period matrix is
\(\left(\begin{smallmatrix}23&11\\3&3\end{smallmatrix}\right)\)
with area \(36\). Both character changes and period words retain
their determinant and inverse; they do not replace a coefficient
by its valuation.

At zero use \(X_0=V_{\rm reg}^{-1},Y_0=Z\) and
\(L_0=B^{-1},T_0=hB^{-1}\):
\[
 L_0=(P^{10}p f_0X_0,P^3Y_0),\qquad
 T_0=(P^8p f_0X_0,P^6pY_0).
 \tag{SI2.41}
\]
Here \(f_0\) is a unit near \(p=0\). For \(r=\operatorname{ord}P\),
\(z=\operatorname{ord}p\), the full matrix and determinant are
\[
 M_0=\begin{pmatrix}10r+z&8r+z\\3r&6r+z\end{pmatrix},\qquad
 \det M_0=36r^2+13rz+z^2.
 \tag{SI2.42}
\]
The exact inspection map \(S=\left(\begin{smallmatrix}3&-3\\0&2\end{smallmatrix}\right)\)
identifies the original lattice with its index-six image lattice.
It gives
\[
 SM_0=\begin{pmatrix}21r+3z&6r\\6r&12r+2z\end{pmatrix}.
 \tag{SI2.43}
\]
The whole periodic support is the negative-square interpolation of
this symmetric integral form on that specified lattice. Its transverse
jumps are \(6r,15r+3z,6r+2z\). Pulling the exact dual cell back gives
\[
 (0,0),\ (-5r-z,0),\ (-10r-z,-3r),\
 (-13r-2z,-6r-z),\ (-8r-z,-6r-z),\ (-3r-z,-3r-z).
 \tag{SI2.44}
\]
This polygon spans two actual Tate source bands, separated by
\(y+3r+z=0\). Their maps are
\[
 (x_1,y_1)=(QZ,Z^{-1}),\qquad
 (x_2,y_2)=(Q^2Z,Q^{-1}Z^{-1}),
 \tag{SI2.45}
\]
with \(Q=P^3p\); all exponents are nonnegative in their respective
half-polygons. The coefficient arrow is SI2.18--20 with the actual
band index. The source-compatible regular four-fan is being
constructed from this periodic parent and these source cuts.

At one put \(v_1=p-1,\ X_1=Z^{-1},Y_1=V_{\rm reg}/Z\).
For \(L_1=h^{-1},T_1=hB\), the exact maps are
\[
 L_1=(P^3pX_1,PpY_1),\qquad
 T_1=(p^{-1}X_1,P^{12}v_1^{12}c_{12}Y_1).
 \tag{SI2.46}
\]
Their matrix is
\(\left(\begin{smallmatrix}3r&0\\r&12r+12z\end{smallmatrix}\right)\),
determinant \(36r^2+36rz\), with \(z=\operatorname{ord}v_1\).
The source and line \(p\)-factors are still present.

At infinity use \(u=P/R\) and
\[
 f_\infty=(1-u)^{12}c_{12}(Ru,R),\qquad f_0=u^{-12}f_\infty.
\]
The same common words are
\[
 L=(R^{23}f_\infty^2X_c,R^3Z_c),\quad
 T=(R^{11}f_\infty X_c,R^3uZ_c).
 \tag{SI2.47}
\]
Their matrix is
\(\left(\begin{smallmatrix}23r&11r\\3r&3r+z\end{smallmatrix}\right)\),
determinant \(36r^2+23rz\). The substitution \(P=Ru,p=u^{-1}\),
all reciprocal point generators and the selected-height factor remain
in this equality.

The actual original four-form, in the zero characters, is
\[
 \Omega=(2J_{ts}/a_T)\,P\,dP\wedge dp\wedge
                        d\log X_0\wedge d\log Y_0.
 \tag{SI2.48}
\]
Here \(J_{ts}=\det\partial(t,s)/\partial(P,R)\), with \(J_{ts}(0)=1/2\),
and \(a_T(0)=1\). Its full unit and exceptional \(P\) factor are
retained. Thus a primitive zero-end ray of base heights \((r,z)\)
has original canonical order \(2r+z-1\), with the actual chart
Jacobian sign still to be recorded on each regular cone.

The earlier sixth-power period attempt is preserved separately
with its exact constant-symmetrizer obstruction. It supplies no
completed boundary and no numerical input here.


% End expanded verified cp36 input


% Expanded verified cp36 input: notes/part3_e6_selected_i3_m2_return_support_02.tex
\par\bigskip\noindent\textbf{P3K.111 supporting proof 2: Original marked point/O graph and inverse}\par\nobreak
\par

\medskip\noindent\textbf{Selected-I3 \(M_{2}\) N12: the original point graph and complete Cartier coefficient}

\par\nobreak

This note retains SI2.1-\/-6, the original source parameter functions of
E81 and the complete new common/end atlases of supporting proof 20. It establishes
local source and line-module arrows. The finite marked boundary
attachment and compact fixed-normal ledger are separate calculations;
no characteristic pair is assigned here.

\par

\medskip\noindent\textbf{I2P.1. The literal point and source-node graph}

\par\nobreak

Write \(Q=P^{2} R\), \(b=P^{2}\) and \(e=R\), with every E81 analytic unit retained.
The actual Tate representative of 2G is b, and \(be=Q\). In a retained
source band j put \(x=Q^j Z\), \(y=Q^{1-j}/Z\), so \(xy=P^{2} R\).
The source period is always its written arrow; changing j is accompanied
by its \(P^{2j}\) line factor below.

The closure of the graph of \([x:P^{2}]\) is the blowup of the ideal
\((x,P^{2})\), with its strict-transform equations
\[
 x=P^2 a,\quad ay=R;\qquad
 P^2=x c,\quad y=Rc.                                    \tag{I2P.1}
\]
Both project to the original source by the displayed x,y,P,R, and
therefore preserve the original s,t through E81. On the common open
set \(c=a^{-1}\). The point section is \(a=1\) or \(c=1\). The first chart is
smooth with coordinates (P,a,y), \(R=ay\). The second is the product
of \(xc=P^{2}\) with the R coordinate. Its actual two-chart resolution is
\[
 P=uv,\quad x=u,\quad c=uv^2,\quad y=Ruv^2;
 \qquad
 P=uv,\quad x=u^2v,\quad c=v,\quad y=Rv.                 \tag{I2P.2}
\]
The two charts have the literal overlap \(u_2=v_1^{-1}\),
\(v_2=u_1v_1^{2}\). Their maps give the ordered P,P,R resolution of
\(xy=P^{2}R\). The projection is the composition of the graph blowup and
this explicit \(A_{1}\) resolution. These are proper morphisms over the
original parameter bidisc. The original zero section is retained
in its separate source chart x near 1 (or the adjacent-period chart).
On Q nonzero their Tate points cannot coincide: equality would give
\(P^{2}=Q^k\); the valuations for \(|P|,|R|<1\) exclude
this for every integer k. Indeed \(k\le 0\) would make the right side
have modulus at least 1, and \(k\ge 1\) gives \(1=P^{2k-2}R^k\), whose modulus
is strictly less than 1. Over the parameter axes the strict point
remains in \(a=1\) or \(c=1\), with \(x=P^{2}\) small, whereas the zero section is
in x near 1. These are disjoint source neighbourhoods on the graph.

\par

\medskip\noindent\textbf{I2P.2. Full theta products and the actual module generators}

\par\nobreak

The complete theta identity on a source-node neighbourhood is
\[
 b f_b(x)=(b-x)E_{b,e}(x,y),
\]
\[
 E_{b,e}=
 \frac{\prod_{r\ge1}(1-Q^{r-1}ex)(1-Q^{r-1}by)}
 {(1-x)\prod_{r\ge1}(1-Q^r x)(1-Q^{r-1}y)}.              \tag{I2P.3}
\]
This is obtained by factoring the \(r=0\) factor from
\(\Theta_Q(x/b)/\Theta_Q(x)\) and substituting \(Q=be=xy\). On the
chosen closed small node chart all factors in the displayed ratio
are units; the products converge normally and their reciprocals do
also. Its central value is 1. The actual Miller unit is
\(\eta=a_T \Theta'_Q(1) \Theta_Q(P^{2})/\Theta_Q(P)^{2}, \eta(0)=-1\),
with the original invariant differential \(a_T(0)=1\). This entire
function and its reciprocal remain in every following arrow.

Let \(v_{2}\) be the coefficient in the original canonical meromorphic
point-line frame. The exact SI2.2/band arrow is
\[
 V_{reg}=\eta^{-1}P^{-2j} f_b(x)v_2,
 \qquad v_2=\eta P^{2j}V_{reg}/f_b(x).                    \tag{I2P.4}
\]
This equality denotes a morphism of the specified Cartier modules.
In I2P.1 use \(e_a=1/(a-1)\), \(e_c=1/(c-1)\) as the local meromorphic
point-divisor generators, and write \(v_{2}=e_a A=e_c C\). Direct substitution
in I2P.3 gives the full frame arrows and their inverses:
\[
 V_{reg}=-\eta^{-1}P^{-2j}E_{b,e}A,
 \qquad A=-\eta P^{2j}E_{b,e}^{-1}V_{reg},               \tag{I2P.5}
\]
\[
 V_{reg}=\eta^{-1}P^{-2j}c^{-1}E_{b,e}C,
 \qquad C=\eta P^{2j}cE_{b,e}^{-1}V_{reg}.                \tag{I2P.6}
\]
The apparent c pole in I2P.6 has its stated Cartier type. It is
retained when pulling back to the two I2P.2 charts. On their point
neighbourhood c is a unit. On the two point charts \(e_c=-a e_a\),
so \(C=-c A\). Substitution verifies exactly the equality of I2P.5 and
I2P.6, including the minus sign, \(\eta\) and every P power.

Near O, let x be near 1 and use its actual generator \(e_O=x-1\) for
the negative point divisor. The full smooth-chart theta unit \(E_O\) is
\[
 f_b(x)=\frac{x-b}{b(x-1)}E_O,
\quad E_O=
 \frac{\prod_{r\ge1}(1-Q^r x/b)(1-Q^r b/x)}
 {\prod_{r\ge1}(1-Q^r x)(1-Q^r/x)}.
\]
Here \(Q/b=R\), so \(E_O\) and its reciprocal extend through the whole
small parameter bidisc. With \(v_{2}=e_O C_O\) the actual arrow is
\[
 V_{reg}=\eta^{-1}P^{-2j-2}(x-P^2)E_O C_O,
 \quad C_O=\eta P^{2j+2}(x-P^2)^{-1}E_O^{-1}V_{reg}.
                                                               \tag{I2P.7}
\]
The extra factor \(P^{2}\) at O is retained. It has not been folded into
\(\eta\) or an unspecified unit. The relation with the actual original
\(M_{2}\) module is SI2.2\textquotesingle s \(r_{2}=f_{2}^{-1}\) and SI2.2\textquotesingle s telescoping theta identity;
I2P.4 is exactly that relation before the graph substitutions. Thus
I2P.5-\/-7 carry both original marked divisors and the original
point-to-tensor coefficient, rather than choosing a new line from
its Chern class.

\par

\medskip\noindent\textbf{I2P.3. Directed source periods and suspension on every frame}

\par\nobreak

For the actual source period H, set \(j'=j-1\) after applying
\((Z,V_{\rm reg})\longmapsto (QZ,P^{2}V_{\rm reg})\). Then x,y and each coefficient \(A,C,C_O\)
in I2P.5-\/-7 are literally unchanged. This proves the full module
square on both source-node charts and on O; no unrecorded reindex
is available in this square.

Here \(q=t-t_c(s)=R c_q\) is the actual selected height,
\(R-P=s u_s\) and \(t=P u_t\). For the actual suspension B, write
\[
 d_B=\lambda_0s^{12}q\eta^6t^{-3},\quad
 (Z,V_{reg})\longmapsto(P^{-3}Z,d_B V_{reg}).             \tag{I2P.8}
\]
If the target retained band is j\textquotesingle, its full source coordinate is
\(x'=Q^{j'-j}P^{-3}x\). The meromorphic point coefficient transforms by
\[
 v_2'=d_B P^{2(j'-j)}\frac{f_b(x)}{f_b(x')}v_2.          \tag{I2P.9}
\]
The inverse has the reciprocal coefficient evaluated at the actual
inverse source point. Equations I2P.5-\/-7, on source and target, give
its local module coefficients by substitution, retaining both theta
units and their actual marked sections. At \(j'=j\) this is exactly
SI2.2\textquotesingle s six-factor chronological Hom: the original f2 theta identity
telescopes through all three source translations. For general j\textquotesingle{}
the equality \(f_b(Q^r z)=b^r f_b(z)\) supplies precisely the displayed
\(P^{2(j'-j)}\), so no new Hom scalar occurs.

On each localized exterior chart and on the thirty-six lower-source
charts used below, x,y,P,R are the full monomials and units of SI2.6
and I2M.  The exact cone inequalities and chart substitutions in
supporting proof 20 C give morphisms to I2P.1 and its resolution,
together with the complete line coefficients above.  The other
common/end interior charts form the independently constructed
regular toric filling; no additional graph blowup is inserted there.
The full I2P.8--9 maps carry every translated graph owner on the
regions where the exterior and lower-source comparisons use that
graph.  The complete boundary assembly is supplied in supporting
proof 20 F.

The source translation, full coefficient, point graph and inverse
Cartier maps above are fixed before any characteristic calculation.
The full \(\eta\) sixth power retains \(\Theta'_Q(1)^{6}\); the forward
point frame and its inverse retain \(\eta\) itself.


% End expanded verified cp36 input


% Expanded verified cp36 input: notes/part3_e6_selected_i3_m2_return_support_03.tex
\par\bigskip\noindent\textbf{P3K.111 supporting proof 3: Complete common period dictionary}\par\nobreak
\par

\medskip\noindent\textbf{Selected-I3 \(M_{2}\) N12: the complete common period dictionary}

\par\nobreak

The original SI2 source, point/O modules, finite character,
full annular reference and fixed original scalar are retained.
Use \(Q=P^{2}R\), \(R-P=s u_s\), \(t=P u_t\), \(q=R c_q\), \(p=R/P\) and

\[
\begin{gathered}
c_{12}=\lambda_0 c_q \eta^{6}/(u_s^{12} u_t^{3}), \\
f_{0}=(p-1)^{12} c_{12}.
\end{gathered}
\]

\par

\medskip\noindent\textbf{I2M.1. Common characters and full periods}

\par\nobreak

The genuine point coefficient has \(V_{\rm reg}=\eta^{-1} V_{{\rm pt},2}\). In the
ordered genuine characters \((V_{\rm reg},Z)\), the original periods are

\[
\begin{gathered}
K=hB=(P^{12} p f_{0} V_{\rm reg},pZ), \\
J=h=(P^{2} V_{\rm reg},P^{3}pZ).
\end{gathered}
\]

Set \((X_c,Z_c)=(V_{\rm reg}/Z,Z)\), \(L_c=K^{2}J\) and \(T_c=KJ\). Their exact maps are

\[
\begin{gathered}
L_c:(X_c,Z_c)\longmapsto (P^{23} p^{-1} f_{0}^{2} X_c,P^{3}p^{3} Z_c), \\
T_c:(X_c,Z_c)\longmapsto (P^{11} p^{-1} f_{0} X_c,P^{3}p^{2} Z_c).
\end{gathered}
\]

The inverse character is \(V_{\rm reg}=X_c Z_c\). The inverse period words
are \(K=L_c T_c^{-1}\) and \(J=T_c^{2} L_c^{-1}\). Both matrices have determinant 1.
Every entire coefficient is fixed by both period actions, so
the displayed compositions commute exactly, including their
units. The generic period matrix is \(\begin{pmatrix}23&11\\3&3\end{pmatrix}\), with
original area 36.

\par

\medskip\noindent\textbf{I2M.2. The zero and one parent maps in both directions}

\par\nobreak

At zero use \((X_{0},Y_{0})=(V_{\rm reg}^{-1},Z)\), \(L_{0}=K^{-1}J\) and \(T_{0}=K^{-1}J^{2}\).
The exact maps are

\[
\begin{gathered}
L_{0}:(X_{0},Y_{0})\longmapsto (P^{10} p f_{0} X_{0},P^{3}Y_{0}), \\
T_{0}:(X_{0},Y_{0})\longmapsto (P^{8} p f_{0} X_{0},P^{6} pY_{0}).
\end{gathered}
\]

The common character and period matrices are

\[
\begin{gathered}
A_{c0}=\begin{pmatrix}-1&-1\\0&1\end{pmatrix}, C_{c0}=\begin{pmatrix}-5&-3\\3&2\end{pmatrix}.
\end{gathered}
\]

They act by \((X_c,Z_c)=((X_{0}Y_{0})^{-1},Y_{0})\), \(L_c=L_{0}^{-5} T_{0}^{3}\) and
\(T_c=L_{0}^{-3} T_{0}^{2}\). Their exact inverse maps are
\((X_{0},Y_{0})=((X_cZ_c)^{-1},Z_c)\), \(L_{0}=L_c^{-2} T_c^{3}\), \(T_{0}=L_c^{-3} T_c^{5}\).
Both matrices have determinant −1. Direct composition of the
full functions in these displayed arrows gives I2M.1.

At one use \((X_{1},Y_{1})=(Z^{-1},V_{\rm reg}/Z)\), \(L_{1}=J^{-1}\) and \(T_{1}=K\).
Writing \(v=p-1\), the entire period functions are

\[
\begin{gathered}
L_{1}:(X_{1},Y_{1})\longmapsto (P^{3}p X_{1},PpY_{1}), \\
T_{1}:(X_{1},Y_{1})\longmapsto (p^{-1} X_{1},P^{12} v^{12} c_{12} Y_{1}).
\end{gathered}
\]

The exact common matrices are

\[
\begin{gathered}
A_{c1}=\begin{pmatrix}0&1\\-1&0\end{pmatrix}, C_{c1}=\begin{pmatrix}-1&-1\\2&1\end{pmatrix}.
\end{gathered}
\]

Thus \((X_c,Z_c)=(Y_{1},X_{1}^{-1})\), \(L_c=L_{1}^{-1} T_{1}^{2}\), \(T_c=L_{1}^{-1}T_{1}\).
Their inverses are \((X_{1},Y_{1})=(Z_c^{-1},X_c)\), \(L_{1}=L_c T_c^{-2}\),
\(T_{1}=L_c T_c^{-1}\). Both matrices have determinant 1. Substitution
of \(f_{0}=v^{12} c_{12}\) gives exactly the two full I2M.1 maps.
The actual generic matrix identities are

\[
\begin{gathered}
A_{c0} \begin{pmatrix}10&8\\3&6\end{pmatrix} C_{c0} \\
= A_{c1} \begin{pmatrix}3&0\\1&12\end{pmatrix} C_{c1} \\
= \begin{pmatrix}23&11\\3&3\end{pmatrix}.
\end{gathered}
\]

\par

\medskip\noindent\textbf{I2M.3. Infinity and every full translating coefficient}

\par\nobreak

At infinity the characters already equal \((X_c,Z_c)\). The actual
height is \(R=pP\), \(u=p^{-1}\) and \(f_{\rm infinity}=p^{-12} f_{0}\). Consequently

\[
\begin{gathered}
R^{23} f_{\rm infinity}^{2}=P^{23} p^{-1} f_{0}^{2}, \\
R^{11} f_{\rm infinity}=P^{11} p^{-1} f_{0}, \\
R^{3}=P^{3}p^{3}, R^{3}u=P^{3}p^{2}.
\end{gathered}
\]

These are exact full analytic equalities and the literal
arrows of SI2.47. For the translate \(L_c^i T_c^j\), the full unit
vectors in character order \((X_c,Z_c,P)\) are

\[
\begin{gathered}
u_p(i,j)=(-i-j,3i+2j,0), \\
u_f(i,j)=(2i+j,0,0).
\end{gathered}
\]

Their original valuation shift is \(((23i+11j)h,(3i+3j)h,0)\).
At zero p remains the second original parameter, so its full
orders remain coordinate orders. At one the exact substitution
\(f_{0}=v^{12} c_{12}\) adds \(12(2i+j)\) to the first character\textquotesingle s v order.
At infinity use both \(P=Ru\) and \(f_{0}=u^{-12} f_{\rm infinity}\). All inverse
arrows use the same full functions and opposite word. These
identities fix every period, source and unit map used by the
common toric overlay.


% End expanded verified cp36 input


% Expanded verified cp36 input: notes/part3_e6_selected_i3_m2_return_support_04.tex
\par\bigskip\noindent\textbf{P3K.111 supporting proof 12: Actual whole scalar reference body}\par\nobreak
\par

\medskip\noindent\textbf{Selected-I3 \(M_{2}\): the complete fixed-source reference body and its boundary line}

\par\nobreak

Working proof, 18 November 2026. This construction retains the original
source and point line of SI2.1--4, its actual selected height
\(q=t-t_c(s)\), the cubic good cap, and the four-root-turn square
SI2.31--37. The source parameter is \(s\). The independent scalar-disc
coordinate below is \(\sigma\); its root is \(r_0\sigma^4\), so its
suspension scalar is \(\lambda_0\sigma^{12}\). We construct the
reference over one fixed nonzero source value \(s_*\), with the
original small \(t\)-disc containing both the \(I_2\) fibre at zero
and the selected \(I_1\) fibre at \(t_c(s_*)\).

The calculation identifies the entire source pieces on the scalar
central face, their original line, all three retained vanishing
discs, and their prescribed outer boundary. It supplies the relative
reference normal square used when the twelve one-end pieces are
doubled. The collision normal data, the moving selected fixed surface,
its Hopf cap, and the total outer characteristic transfer are retained
as their own geometric pieces. No return is assigned here.

\par

\medskip\noindent\textbf{RB.1. The original source cover and an exact frame at the unselected I2 fibre}

\par\nobreak

Fix the actual functions of SI2 and E81:
\[
 Q=P^2R,\qquad R-P=s u_s,\qquad t=P u_t,
 \qquad q=R c_q,\qquad Q=\kappa q,\quad\kappa=P^2/c_q.
 \tag{RB.1}
\]
The Miller coefficient is the full
\(\eta=a_T\Theta'_Q(1)\Theta_Q(P^2)/\Theta_Q(P)^2\), with
\(\eta(0)=-1\). Put
\[
 V_{\rm pt}=f_{P^2}(Z)v_2,\quad
 \alpha=P^2,\quad g=P^{-3},\quad
 \Gamma=\lambda_0\eta^6t^{-3}.
\]
On the selected neighbourhood of the fixed source, \(P,t,\kappa,
\alpha,g,\Gamma\) are nonzero analytic functions. Its complete
periods are
\[
 h:(Z,V_{\rm pt})\longmapsto(QZ,\alpha V_{\rm pt}),\qquad
 B_\sigma:(Z,V_{\rm pt})\longmapsto
                 (gZ,\sigma^{12}q\Gamma V_{\rm pt}).
 \tag{RB.2}
\]
Every coefficient is evaluated at the same fixed-source base point;
the source translation is the actual chronological \(-3G\).

Near \(t=0\), \(R\) is a unit because \(s_*\ne0\). The full
theta identity \(\Theta_Q(Q^{-1}z)=-Q^{-1}z\Theta_Q(z)\) gives
\[
 f_{P^2}(Z)=-\frac{Z}{P^2}f_{R^{-1}}(Z),\qquad
 V_I=-P^2Z^{-1}V_{\rm pt}=f_{R^{-1}}(Z)v_2.
 \tag{RB.3}
\]
This is an equality of the original Cartier modules, with inverse
\(V_{\rm pt}=-ZP^{-2}V_I\). On their common smooth source annulus
the factors are nonzero. At the marked sections the two sides use
their actual point or zero generators; hence the identity extends
as a module morphism there as well. Its two period squares are
\[
 h:(Z,V_I)\longmapsto(QZ,R^{-1}V_I),\qquad
 B_\sigma:(Z,V_I)\longmapsto
       (P^{-3}Z,\sigma^{12}\Gamma_I V_I),
 \quad \Gamma_I=\lambda_0q\eta^6/u_t^3.
 \tag{RB.4}
\]
In particular \(\Gamma_I\) is an actual unit on the fixed-source
\(I_2\) neighbourhood. The cancellation of the apparent \(P^{-3}\)
in the line coefficient follows from the displayed \(Z^{-1}\)
factor in RB.3; it is not a replacement of the Hom by its leading term.

The two-node source is covered by
\[
 x_i=Z/P^i,\qquad y_i=P^{i+1}/Z,\qquad x_i y_i=P.
 \tag{RB.5}
\]
On consecutive charts \(x_{i+1}=y_i^{-1}\),
\(y_{i+1}=x_i y_i^2\). The full periods are
\[
 h:(i,x_i,y_i)\longmapsto(i+2,Rx_i,R^{-1}y_i),\qquad
 B_\sigma:(i,x_i,y_i)\longmapsto(i-3,x_i,y_i).
 \tag{RB.6}
\]
These maps retain the unit \(R\), and RB.4 gives their line lifts.
The original base coordinate is still \(t=Pu_t\); all maps keep
the original source parameter. The source graph I2P.1--2 restricts
here to these two ordered \(P\)-factor charts because its third
factor \(R\) is a unit. The strict point and zero sections are those
of I2P.5--7. Thus RB.3--6 identify the two-node source and its point
line with the original resolved family, including its marked sections.

\par

\medskip\noindent\textbf{RB.2. Fourteen ordered factors at the selected node}

\par\nobreak

On a selected Tate source band put
\[
 x_i=Z/Q^i,\qquad \bar y_i=Q^{i+1}/(\kappa Z),
                  \qquad x_i\bar y_i=q.
 \tag{RB.7}
\]
This retains the original height rather than substituting \(Q\)
for it. Its adjacent-band arrow is
\[
 x_{i+1}=(\kappa\bar y_i)^{-1},\qquad
 \bar y_{i+1}=Q\bar y_i.
 \tag{RB.8}
\]
The source period \(h\), accompanied by \(i\mapsto i+1\), fixes
\(x_i,\bar y_i\); its line coefficient is \(\alpha\). The source
part of \(B_\sigma\) sends \((x_i,\bar y_i)\) to
\((gx_i,g^{-1}\bar y_i)\). This is the actual translation on the
fixed selected node chart, including the full base-dependent units.

Use the associated line-node charts
\[
 U_j=V_{\rm pt}/(\sigma^{12}q)^j,\qquad
 W_j=(\sigma^{12}q)^{j+1}/V_{\rm pt},\qquad
               U_jW_j=\sigma^{12}x_i\bar y_i.
 \tag{RB.9}
\]
The chronological factor order is
\[
 \underbrace{\sigma,\ldots,\sigma}_{12\ {
m factors}},
                       x_i,\bar y_i.
 \tag{RB.10}
\]
The fourteen charts of its ordered graph have the following literal
maps. For \(1\le r\le12\), use coordinates
\((u_r,v_r,x_i,\bar y_i)\) and
\[
 \sigma=u_rv_r,\quad
 U_j=\sigma^{r-1}u_r,\quad
 W_j=\sigma^{12-r}x_i\bar y_i v_r.
 \tag{RB.11}
\]
For factor thirteen use
\[
 x_i=u_xv_x,\quad U_j=\sigma^{12}u_x,\quad
 W_j=\bar y_i v_x,
 \qquad (u_x,v_x,\sigma,\bar y_i).
 \tag{RB.12}
\]
For factor fourteen use
\[
 \bar y_i=u_yv_y,\quad U_j=\sigma^{12}x_i u_y,\quad
 W_j=v_y,
 \qquad (u_y,v_y,\sigma,x_i).
 \tag{RB.13}
\]
These are the standard affine charts of the successive graph closures
of the written ordered factor ideals; their displayed products give
the original hypersurface exactly. Each chart is an affine smooth
fourfold. Consecutive charts overlap by the actual reciprocal arrow
\[
 u_{r+1}=v_r^{-1},\qquad v_{r+1}=\sigma v_r
                         \quad(1\le r<12),
 \tag{RB.14}
\]
and the last two transitions are
\[
 u_x=v_{12}^{-1},\quad v_x=x_i v_{12},\quad
 \sigma=u_{12}v_{12};\qquad
 u_y=v_x^{-1},\quad v_y=\bar y_i v_x,\quad x_i=u_xv_x.
 \tag{RB.15}
\]
Their inverses are obtained from these same equations, on the stated
nonzero overlap coordinate. This proves both the graph maps and their
directions, without identifying two different factor orders.

The physical action has weights \((+1,-1,0,0)\) on each displayed
chart. In RB.11 its fixed surface is exactly
\[
 u_r=v_r=0,
\]
with the entire source chart \((x_i,\bar y_i)\) left over. It contains
the selected source node \(x_i=\bar y_i=0\). In RB.12--13 the fixed
surfaces are instead \(u_x=v_x=0\) and \(u_y=v_y=0\), with
\((\sigma,\bar y_i)\) or \((\sigma,x_i)\) as coordinates.
The first RB.15 overlap requires \(v_{12}\ne0\); neither fixed
surface lies in that overlap. The same statement for every earlier
overlap follows from RB.14. Thus the scalar-source surfaces and the
moving selected fixed surface are disjoint in the resolved graph.
This conclusion follows from their fixed equations and exact overlap,
not from the names assigned to the pieces.

\par

\medskip\noindent\textbf{RB.3. All source, line and period squares on the reference pieces}

\par\nobreak

The full \(B_\sigma\) arrow on RB.9 has target label \(j+1\):
\[
 (x_i,\bar y_i,U_j,W_j)\longmapsto
        (gx_i,g^{-1}\bar y_i,\Gamma U_j,\Gamma^{-1}W_j).
 \tag{RB.16}
\]
On the first twelve factor charts it is
\((u_r,v_r)\mapsto(\Gamma u_r,\Gamma^{-1}v_r)\).
On the last two charts it is respectively
\[
 (u_x,v_x)\longmapsto(\Gamma u_x,g\Gamma^{-1}v_x),
 \qquad
 (u_y,v_y)\longmapsto(\Gamma g^{-1}u_y,\Gamma^{-1}v_y).
 \tag{RB.17}
\]
The inverse uses the reciprocal coefficients at the inverse source
point. Substitution in RB.11--13 verifies the source, line, height and
every product. The \(h\) arrow has \(i\mapsto i+1\), fixes the
source coordinates, and sends each normal pair to
\((\alpha u,\alpha^{-1}v)\). These maps carry RB.14--15 to themselves.
They do not change the factor index \(r\).

The moving fixed charts also have an explicit full source-band
overlap. On \(\bar y_i\ne0\), the old x-factor chart and the
next band\textquotesingle s y-factor chart are related by
\[
 x_{i+1}=(\kappa\bar y_i)^{-1},\qquad
 u_y'=\kappa\bar y_i u_x,\qquad
 v_y'=\bar y_i v_x,\qquad \sigma'=\sigma.
 \tag{RB.17a}
\]
Indeed \(u_y'v_y'=\kappa x_i\bar y_i^2=Q\bar y_i\),
and the two expressions for \(U_j,W_j\) in RB.12--13 agree
literally. The inverse has
\(\bar y_i=(\kappa x_{i+1})^{-1}\),
\(u_x=x_{i+1}u_y'\), and
\(v_x=\kappa x_{i+1}v_y'\), with the same original base \(q\).
The scalar-factor normal coordinates remain unchanged under
this source-band transition. Substitution of RB.16--17 shows
that RB.17a commutes with the full B and h arrows, including
\(g,\Gamma,\alpha,\kappa\).

On the overlap \(q\ne0\) with an ordinary scalar chain
\(V_{\rm pt}=\sigma^k A_k\), \(A_kB_k=\sigma\), the exact map is
\[
 k=12j+r-1,\qquad A_k=q^j u_r,\qquad B_k=q^{-j}v_r.
 \tag{RB.18}
\]
For example RB.16 becomes \(k\mapsto k+12\) with multiplier
\(q\Gamma\) on \(A_k\), exactly RB.2. Both integer powers of
\(q\) remain in the arrow. At \(q=0\) the full factor charts
RB.11--13 supply its continuation; RB.18 is used only on its
specified nonzero overlap.

On the unselected neighbourhood use the ordinary chain in \(V_I\)
from RB.3. The period \(h\) has source label \(i+2\), normal
factor \(R^{-1}\); \(B_\sigma\) has source label \(i-3\), scalar
label \(k+12\), and normal factor \(\Gamma_I\). On the common
smooth source region the normal arrow is precisely
\[
 (A_k^I,B_k^I)=(-P^2Z^{-1}A_k,-ZP^{-2}B_k).
 \tag{RB.19}
\]
Its product is unchanged and the source coordinate is the identity.
Equations RB.3--4 show its two full period squares. Thus both
unselected source nodes and the selected source node belong to the
same retained reference source pieces.

The genuine normal maps are obtained before differentiating by
the original point/O modules. On a source band \(\ell\), I2P gives
\[
 V_{\rm pt}=-P^{-2\ell}E_{P^2,R}A
   =P^{-2\ell}c^{-1}E_{P^2,R}C,\quad C=-cA,
 \tag{RB.20}
\]
and at the zero section
\[
 V_{\rm pt}=P^{-2\ell-2}(x-P^2)E_OC_O.
 \tag{RB.21}
\]
The full normally convergent units and their inverses occur here.
The two module generators are \((a-1)^{-1},(c-1)^{-1}\), and the
zero generator is \(x-1\). These equations carry RB.18--19 into
the actual coefficient of \(M_2\), not merely into a line with
matching degree. In a genuine local vector frame \(e_i\) write
the coefficient as \(v_i\). For \(v_b=g_{ba}v_a\), the scalar
normal coordinates satisfy
\[
 A_b=g_{ba}A_a,\qquad B_b=g_{ba}^{-1}B_a.
 \tag{RB.22}
\]
The source graph and all negative normal maps are therefore
reciprocal to the positive ones. They identify the two eigenlines,
in physical order \((-1,+1)\), with \((M_2^{-1},M_2)\).

For each \(r=1,\ldots,12\), the \(h\)-glued scalar-source piece
is the entire compact resolved fixed-source body \(S_{s_*}\).
The \(B_\sigma\) action changes \(j\) by one, so a representative
\(j=0\) has no further source identification by a nonzero B word.
Equivalently, at the unselected fibre the joint labels change by
\((i,k)\mapsto(i+2,k)\) and \((i-3,k+12)\), leaving one full
two-node source for each \(k\bmod12\). These statements concern
the pieces inside the retained \(t\)-disc. The cubic finite cap
has its own label shift four, written next; we do not identify these
twelve pieces with twelve closed components of the fully capped space.

\par

\medskip\noindent\textbf{RB.4. The actual finite cap, parameter closing and selected-tube boundary}

\par\nobreak

On the original infinity annulus retain \(q_\infty=t^{-1}\), the
root \(q_\infty=z^3\), and
\[
 h_\infty=1-t_c(s)z^3=q_\infty q.
\]
The root source, the correction \(E_P-E_O\), the normalized chord,
finite character, and sector maps are exactly SI2.22--28. In particular
\(A_0^3=B_0\). Write its full genuine line map as
\(A_0:(p,v)\mapsto(a_0p,a_{\rm line}(p)v)\). On the scalar chain
\(v=\sigma^kA_k\), \(A_kB_k=\sigma\), the actual primitive cap map is
\[
 (p,k,A_k,B_k)\longmapsto
 (a_0p,k+4,r_0a_{\rm line}(p)A_k,
                    r_0^{-1}a_{\rm line}(p)^{-1}B_k).
 \tag{RB.23}
\]
Its chronological cube is the full twelve-step B map; the coefficient
is the product of \(a_{\rm line}\) at its three successive source
points. The inverse uses \(a_0^{-1}p\). Every sector overlap keeps
the original word \(A^{j_0}B^{j_{\alpha\beta}}\) of SI2.28.

On the outer scalar circle the prescribed forward square is
\[
 \Phi_{\rho+4}q^4=h_\infty^4\Phi_\rho.
 \tag{RB.24}
\]
It has the reciprocal endpoint-to-initial square. At the selected
collar \(D_{\rm ord}^4=B^4R_4\), with
\(R_4=(g^{-4},\eta_*^{-4})\),
\(\eta_*=\lambda_{\rm base}\eta^6t^{-3}\). The exact native
comparison is
\[
 C_\rho=(e^{\rho G},e^{\rho J+3\pi i\rho(\rho-1)}),
 \quad G=\operatorname{Log}g,\quad J=\operatorname{Log}\eta_*,
 \quad C_4=R_4^{-1}.
 \tag{RB.25}
\]
The full source-cut increments of \(G,J\) are \(-6\pi i\).
Their Cox map, including both odd seam signs, is the actual
\(q_H\mapsto e^{3\pi i\rho}q_He^{9\pi i\rho}\) of PPH.2.
Its left adjoint extends over the specified three-balls and its
written quaternion map carries the principal Hopf line. The complete
DZ factor is the based exponential \(e^{48\pi it(t-1)}\) followed
by the retained physical loop \(e^{60\pi it}\). Thus the boundary
map contains the winding thirty, its two cap lifts, and the original
sector integers. The q-dependent maps and their derivatives use the
protected-tube velocity and inverse flow of PPH.2.

These are boundary and collar maps at nonzero scalar radius. We use
RB.11--13 at the scalar-zero selected corner and RB.23 at the actual
scalar-zero finite cap. In particular we do not apply \(q^4\) as
a unit at the selected zero or claim that it fixes every central
factor chart. Its entire height shift is retained by RB.24--25 and
the source-period factorization. The full module map, not an omitted
height factor, joins these boundary products.

The corresponding actual coupled-source endpoint has the complete
coordinate form I2C.1--3. With
\(f_0=(p-1)^{12}\lambda_0c_q\eta^6/(u_s^{12}u_t^3)\) and
\(A=(c_q/f_0)^4\), it is
\[
 B^{-4}q^4:(Z,V_{\rm pt})\longmapsto
                   (P^{12}Z,P^{-36}A V_{\rm pt}).
 \tag{RB.25a}
\]
Indeed \(\lambda^{-4}(\eta^6t^{-3})^{-4}
=P^{-36}u_s^{48}u_t^{12}/
[\lambda_0^4\eta^{24}(p-1)^{48}]\). In the actual selected
native chart \(q=Za_Nc_N\), \(c_N=V_{\rm pt}\), this sends
\((Z,a_N,c_N)\) to
\((P^{12}Z,P^{24}A^{-1}a_N,P^{-36}Ac_N)\).
The reciprocal chart carries the same map with the two normal
powers \(36,-24\). The positive spinor \((\bar a_N,c_N)\)
therefore has determinant phase \((P/|P|)^{-60}(A/|A|)^2\).
These are exact boundary arrows; they do not change RB.33\textquotesingle s
intrinsic source cube. The complete cut field, its inverse, the
Hermitian-ball map and the held-Cox normal-sphere extension are
I2C.5--11 in
\path{notes/part3_e6_selected_i3_m2_moving_cap_working.md}.
Their cap square is the very same RB.24, with the original
reference integer on each sector. The resolved moving normal
matrix and its Chern classes remain to be obtained from the
new atlas.

The moving fixed surface meets the outer scalar circle. The scalar
source pieces lie on \(\sigma=0\). Their only possible common
source-node overlap was RB.15, where the actual nonzero coordinate
excludes both fixed sets. On the rest of the reference source cover
the fixed equations and reciprocal chain overlaps give the same
separation. In the compact retained source body these are disjoint
closed fixed subsets. Choose an invariant metric in the written
charts and an invariant small tubular neighbourhood of the moving
fixed subset disjoint from all twelve source pieces; a finite cover
and the positive separation of disjoint compact subsets give a
positive admissible radius. Its outer-radius part is transported
through RB.25 and the actual Hopf cap. Hence excising this moving
selected tube leaves the whole source body, including its selected
vanishing disc, in each scalar-source piece.

The finite cap lies outside the retained small \(t\)-disc. Its
common ordinary-annulus product is the prescribed SI2.28 map;
RB.23--24 give its central and boundary arrows. Connections and
products on this cap are transported from the same ordinary annulus.
The source-period and scalar-label faces are paired by RB.2, RB.4,
RB.16 and RB.23 with their inverses. Thus no artificial boundary
around the selected source node is introduced into \(S_{s_*}\).
The moving surface and its cap remain separate terms in the later
fixed-set and outer-boundary ledger.

\par

\medskip\noindent\textbf{RB.5. Relative source and original-line transport, with all three discs}

\par\nobreak

Choose the original source boundary \(|t|=\epsilon_t\), with
\(|t_c(s)|<\epsilon_t/3\) on the source path from zero to \(s_*\).
Retain the proper I2P graph over the resulting compact source domain.
Its central source is the minimal \(I_3\) source. Its three local
charts have equations
\[
 R=ay,\quad s=(ay-P)/u_s(P,ay);
 \qquad P=uv,\quad s=(R-uv)/u_s(uv,R)
 \tag{RB.26}
\]
in the first graph chart and the two \(A_{1}\) charts, respectively.
The partial derivatives \(\partial_Ps\), respectively
\(\partial_Rs\), are nonzero units with central values
\(-1,+1,+1\). The local complex source lifts are explicitly
\[
 X_0=(\partial_Ps)^{-1}\partial_P,
 \qquad X_1=(\partial_Rs)^{-1}\partial_R,
 \qquad X_2=(\partial_Rs)^{-1}\partial_R.
 \tag{RB.27}
\]
On the retained boundary use the actual E81 marked source charts
and the prescribed ordinary boundary product, holding \(t\)
fixed. Their maps have inverses and are smooth over the source
path, so differentiating this product gives a specified boundary
lift \(X_\partial\) with \(ds(X_\partial)=1\).

For clarity the source-period part of that boundary map is explicit.
On its marked Tate cover write \(Z=\exp(2\pi i(u+\tau b))\),
\(\tau=\operatorname{Log}Q/(2\pi i)\). The quotient
\(Q(s,t)/Q(0,t)\) has a logarithm on this boundary source path:
\[
 \frac{Q(s,t)}{Q(0,t)}=
 \frac{u_Q(s,t)}{u_Q(0,t)}\left(1-\frac{t_c(s)}t\right).
 \tag{RB.28}
\]
Use the value zero at \(s=0\). Keeping the marked real coordinates
\((u,b)\) sends \(Z_0\) to
\(\exp(b\operatorname{Log}(Q_s/Q_0))Z_0\). Under the source
deck \(b\mapsto b+1,Z_0\mapsto Q_0Z_0\), its target is multiplied
by precisely \(Q_s\). Before the genuine point/O module
compositions, the line coefficient on this same cover is
\[
 V_{{\rm pt},s}=\exp\!\left(
      b\operatorname{Log}\frac{P_s^2}{P_0^2}\right)V_{{\rm pt},0}.
 \tag{RB.28a}
\]
The logarithm is the full one with value zero at \(s=0\);
\(P_s/P_0=u_P(s,t)/u_P(0,t)\) is a unit on the entire boundary
source path. Under \(b\mapsto b+1\), the old coefficient acquires
\(P_0^2\), and RB.28a acquires exactly the target factor
\(P_s^2\). The inverse uses both reciprocal exponentials at the
recovered marked coordinates \((u,b)\). The full point/O coefficient on each local
chart is then carried by RB.20--21 and its inverse. At a prescribed
boundary-product chart this construction is composed with that
written chart map; its derivative is the boundary lift just fixed.
Thus the integral source-period lift is retained rather than reset.

Choose a finite partition \(\chi_i\) subordinate to the actual
source charts, equal to the prescribed boundary partition on a
collar. Since each local lift projects to one,
\(X=\sum\chi_iX_i\) also projects to one. Along a source path
\(s(a)\), use the real vector field
\(\dot s(a)X+\overline{\dot s(a)X}\). It equals the prescribed
boundary velocity and is tangent to \(|t|=\epsilon_t\).
Properness and compactness make its flow defined on the whole
interval. The inverse is its backward flow. This is an explicit
relative diffeomorphism from the complete central source to
\(S_{s_*}\); no node chart or zero-section disc has been removed.

Use actual vector frames with \(e_j=g_{ij}e_i\). Define connection
coefficients for the lift by
\[
 a_i=-\sum_k\chi_k\,Xg_{ik}/g_{ik},\qquad
 a_j=a_i+Xg_{ij}/g_{ij}.
 \tag{RB.29}
\]
The horizontal vector is \(X-a_iv_i\partial_{v_i}\), with the
corresponding real conjugate. All \(g_{ij}\) are the full original
Cartier transitions RB.20--21; derivatives of the units are retained.
A relative cutoff changes the coefficients to the prescribed full
boundary derivative, which obeys the same transition identity.
Its linear ODE supplies the line transport and its exact inverse.
The reciprocal ODE transports the dual normal. This construction
takes place in the genuine point/O frames, so no meromorphic
coefficient is treated as a nonsingular vector frame.

On a period translate define the transported chart by the full
RB.16 or RB.4 map and its inverse; differentiating the same square
includes the source-dependent period derivative. On the finite
cap use the SI2.28 arrow on its ordinary annulus, extend the
relative lift through its actual proper cap charts, and retain
the central map RB.23. This is the same prescribed-annulus
construction as OG.4, now with \(h_\infty=1-t_c(s)z^3\) and
the complete selected cochain. It preserves RB.24 and its inverse.
All changes used to define these transports are outside the
already separated selected fixed tube or are its specified
quaternion/Hopf transport. Consequently they do not delete any
part of the scalar-source surfaces.

The three central node charts therefore give three actual
vanishing discs in every transported reference source body.
The two-node and one-node descriptions RB.5--8 locate their
fixed-source continuations. In the central plumbing each disc
has the same primitive boundary cycle \(a\); the complementary
elliptic cycle \(b\) survives. The retained zero-section disc
kills the base meridian \(\alpha_t\). Equivalently the three
literal plumbing arrows give
\[
 [a,b]=1,\quad\alpha_t a\alpha_t^{-1}=a,\quad
 \alpha_t b\alpha_t^{-1}=a^3b.
\]
Attaching the three discs and the zero-section disc gives the
actual inclusion
\[
 H_1(\partial S)=\mathbb Z\langle b,\alpha_t\rangle
       \oplus\mathbb Z/3\langle a\rangle
 \longrightarrow H_1(S)=\mathbb Z\langle b\rangle,
 \qquad (a,\alpha_t,b)\longmapsto(0,0,b).
 \tag{RB.30}
\]
Its use here is justified by RB.7--29: the selected source disc
has been explicitly retained and the outer boundary is the
prescribed one throughout the source and line transport.

\par

\medskip\noindent\textbf{RB.6. The twelve one-end normal squares and the intrinsic cube frame}

\par\nobreak

In the actual one-end chart use \(v=p-1\),
\(X=Z^{-1}\), \(Y=V_{\rm reg}/Z\), and
\[
 \alpha_l=Yv^{-l},\qquad\beta_l=v^{l+1}/Y,
                \quad0\le l\le11.
\]
The full directed lower-to-reference map on the retained source
boundary is
\[
 A_l=H_{+,l}\alpha_l,\quad B_l=H_{-,l}\beta_l,
\quad H_{+,l}=\eta u_s^lP^{-l}X^{-1},\quad
 H_{-,l}=\eta^{-1}XP^{l+1}u_s^{-l-1}.
 \tag{RB.31}
\]
It satisfies \(H_{+,l}H_{-,l}=P/u_s\), taking
\(\alpha_l\beta_l=v\) to \(A_lB_l=s\). The inverse keeps
the reciprocal functions at the same source point. The source
and reference parameters in this boundary comparison are the
ones in \(s=Pv/u_s\); \(\sigma\) was the independent parameter
used to construct its fixed-source reference filling.

Here the final fixed-source target includes the actual transport
of RB.5. Denote its source map from \(S_0\) to \(S_s\) by
\(T_s\), and put \(K_s=T_{s_*}T_s^{-1}\). With its transported
line and reciprocal dual maps, the complete boundary morphism is
\[
 (p,\alpha_l,\beta_l)\longmapsto
 \bigl(K_s p,\ K_s^{\rm line}(H_{+,l}(p)\alpha_l),
                 K_s^{\rm dual}(H_{-,l}(p)\beta_l)\bigr).
 \tag{RB.31a}
\]
Its inverse first applies \(K_s^{-1}\) and then the reciprocal
coefficients at that recovered source point. The normal product
remains \(s\), since the two transported line maps are dual.
Thus RB.31 is the full coefficient arrow before this specified
source transport, and RB.31a names its fixed-source target.

Under \(L=(P^3pX,PpY)\), the complete positive and negative
reference factors are \(P^{-2},P^2\). Under
\(T=(p^{-1}X,P^{12}v^{12}c_{12}Y)\) use \(l\mapsto l+12\).
The positive factor is
\[
 p u_s^{12}c_{12}
   =\lambda_0p c_q\eta^6/u_t^3
   =P^2\bigl(\lambda_0q\eta^6/t^3\bigr),
 \tag{RB.32}
\]
and the negative one is its reciprocal. This is exactly the
original \(hB\) point-line coefficient after dividing by the
additional \(s^{12}\). At \(p=1\), the actual central Miller
identity gives \(B_{2,3}=\eta^6P^2=4t^2\), hence
\(\eta^6=4u_t^2\), and RB.32 is \(4\lambda_0\).
All units and derivatives remain in the equality before this
central restriction.

The actual minimal-source point section is on \(C_1\), and
the zero section is on \(C_0\), by the complete I2P/MR source
and module arrows. Thus the intrinsic line degrees are
\((-1,1,0)\). Its source-period cube coefficient is
\(X^2V_{\rm pt}^3\). With
\(\tau_L=\log|X|/\log|Q|\), the full boundary frame is
\[
 W_p=e^{-2\tau_L\operatorname{Log}p}X^2V_{\rm pt}^3.
 \tag{RB.33}
\]
Here \(\operatorname{Log}p\) is the actual small logarithm on
the retained outer source boundary. Its two exact multipliers
under \(L\) cancel. The inverse vector frame at \(p=1\) has
the full divisor \(6F-2C_1-C_2\); multiplication by \(P^{-6}\)
gives coefficient zero on \(C_0\), retaining the entire fibre
clutch. This is the original cube frame of \(M_2\). The
four-turn endpoint factor \(P^{-12}\) in \(R_4\) is a
different, explicitly retained boundary map in RB.24--25; it
does not replace this intrinsic cube coefficient.

The integral matrix on \((C_1,C_2)\) is
\(I=\left(\begin{smallmatrix}-2&1\\1&-2\end{smallmatrix}\right)\).
The vector \((1,0)^t\) has order three in its cokernel and
\(I^{-1}(1,0)^t=(-2/3,-1/3)^t\). Consequently the relative
lift is \(D_+=-(2C_1+C_2)/3\) and \(D_+^2=-2/3\).
For the actual frame RB.33 it is \(2F+D_+\), with the same
square because \(F^2=FC_i=0\).

Finally RB.30 gives
\(H^1(\partial S;\mathbb Z)/\operatorname{im}H^1(S;\mathbb Z)
=\mathbb Z\langle\alpha_t^*\rangle\). Its connecting class
is the primitive fibre, realized by the literal vector-frame
clutch \(e_R=(P/\epsilon)e_L\). The full \(u_t\) carries
the same statement for the original \(t\)-meridian. The
\(b\)-character extends; the zero-period remainder has a
real logarithm extended by a collar cutoff. Thus every actual
boundary ratio from RB.31, including the full Miller unit,
point/O coefficients and the prescribed phase maps, is retained.
It can change the relative lift by a fibre multiple only.

Taking actual third roots of RB.33 gives the order-three flat
boundary connection; its dual is reciprocal. The interpolation
to the transported boundary connection on the other half of
the double is the identical full interpolation with reversed
collar coordinate. Its transgressions cancel pointwise.
The reference normal pair on each retained source piece is
therefore, in physical order,
\[
 (N_-,N_+)=(M_2^{-1},M_2),\qquad
 (c_-^2,c_-c_+,c_+^2)=(-2/3,+2/3,-2/3).
 \tag{RB.34}
\]
This conclusion uses the actual fourteen-factor source inclusion,
the two-node frame RB.3, the full original cap/parameter square,
and the relative source and line transport above. It leaves the
moving selected fixed surface and its cap in the subsequent
ledger, and assigns no total characteristic pair.


% End expanded verified cp36 input


% Expanded verified cp36 input: notes/part3_e6_selected_i3_m2_return_support_13.tex
\par\bigskip\noindent\textbf{P3K.111 supporting proof 13: Collision selected native star and true fan}\par\nobreak
\par

\medskip\noindent\textbf{Selected-I3 \(M_{2}\) N12: the protected native star and its full phase square}

\par\nobreak

This proof uses the new zero endpoint and its actual common lift in supporting proof 20,
the original selected height, and the complete point-line coefficient.
It identifies the native tube on the outer source annulus by explicit
chart maps, their inverses and their source/period squares. It retains
the old I2C and RB sources as records of the preceding point-coordinate
calculation. The correction to their moving-tube interpretation is
specified in NS.6 below. No moving Chern number or return is assigned.

\par

\medskip\noindent\textbf{NS.1. The two protected facets and the original base}

\par\nobreak

Keep the SI2 source and line data
\[
 X=V_{\rm reg}^{-1},\quad Y=Z_{\rm old},\quad
 V_{\rm pt}=\eta/X,\quad R=Pp,\quad Q=P^3p,
\]
\[
 q=q_{\rm old}=Pp c_q,\qquad
 s=P(p-1)/u_s,\qquad t=Pu_t,
\]
\[
 f_0=(p-1)^{12}\frac{\lambda _0c_q\eta^6}{u_s^{12}u_t^3},
 \qquad
 \eta=a_T\Theta'_Q(1)\frac{\Theta_Q(P^2)}{\Theta_Q(P)^2}.
 \tag{NS.1}
\]
All functions are evaluated at the original source point. In particular
the theta derivative, the value \(\eta(0)=-1\), and the complete sixth
power remain present. On the selected outer collar \(P\ne0\), the
functions \(c_q,u_s,u_t,\eta,f_0\) are nonzero. The actual source map
\((P,R)\mapsto(t,s)\) has determinant \(J_{ts}\), with
\(J_{ts}(0)=1/2\). Since \(q=t-t_c(s)\),
\[
 \det\frac{\partial(s,q)}{\partial(P,p)}=-P J_{ts}.
 \tag{NS.2}
\]
Thus the original inverse source map gives the holomorphic inverse
\((s,q)\mapsto(P,p)\) on this collar. The formulas below always use
that inverse, including all its source-dependent units.

The original zero lattice is ordered as
\(\operatorname{ord}(X,Y,P,p)\). Write
\[
 A=(0,0,0,1),\quad B_1=(-1,0,0,1),\quad
 B_2=(-2,-1,0,1),\quad B_3=(-1,-1,0,1).
 \tag{NS.3}
\]
The complete common lift has exactly the two lower facets
\[
 \langle A,B_1,B_2\rangle,\qquad
 \langle A,B_2,B_3\rangle.
 \tag{NS.4}
\]
They are exactly the pure-lower facets of the new 54,121-chart zero
endpoint.  Its entire P-invertible subfan is unchanged from the original
raw model.  The complete physical face extraction has thirty-one charts
in the single orbit represented by these two normal faces.  Supporting
proof 20 and its same-atlas certificate retain every original ray,
source band, normal character, canonical order and full period unit.

The actual physical vector is \(\delta=(-1,0,0,0)\). On the first
face \(\delta=B_1-A\); on the second \(\delta=B_2-B_3\).
The actual zero period
\[
 T_0=hB^{-1}:(X,Y,P,p)\longmapsto
        (P^8p f_0 X,P^6pY,P,p)
 \tag{NS.5}
\]
sends \(B_2\) to \(B_1\) and \(B_3\) to \(A\), with the same
physical vector. This is the specific period identification used below.
No source or fibre character is selected from a fan label.

\par

\medskip\noindent\textbf{NS.2. Both regular native charts and their literal overlap}

\par\nobreak

On the first lower facet use its actual primitive dual characters
\[
 u_-=pX/Y,\qquad u_+=Y^2/X,\qquad z=Y^{-1}.
 \tag{NS.6}
\]
Their inverse is
\[
 X=u_+^{-1}z^{-2},\quad Y=z^{-1},\quad p=u_-u_+z.
 \tag{NS.7}
\]
These inverse torus formulas specify the affine chart by its forward
semigroup generators; negative powers in NS.7 are used only on the
corresponding torus localization. With \(P\) invertible, NS.6 are
the three regular coordinates of this lower chart. The physical
weights of \((u_-,u_+,z)\) are \((-1,+1,0)\). Define
\[
 a_\infty=\frac{Pc_q}{\eta}u_-,\qquad
 c_\infty=\eta u_+,\qquad z=z.
 \tag{NS.8}
\]
Then \(q=z a_\infty c_\infty\). Given \((s,z,a_\infty,c_\infty)\),
first compute that literal \(q\), recover \((P,p)\) by NS.2, and set
\[
 u_-=\frac{\eta}{Pc_q}a_\infty,\qquad
 u_+=\eta^{-1}c_\infty.
 \tag{NS.9}
\]
Both compositions are identities, including their dependence on
\(q\). In particular the coefficients in NS.9 are evaluated at the
recovered source point, not at a substituted new normal product.

On the second lower facet use its primitive dual coordinates
\[
 u_A=pY,\qquad u_2=Y/X,\qquad u_3=X/Y^2,
 \tag{NS.10}
\]
with inverse \(X=u_2^{-2}u_3^{-1}\),
\(Y=u_2^{-1}u_3^{-1}\), \(p=u_Au_2u_3\).
Their physical weights are \((0,+1,-1)\).
Apply the full period NS.5 before comparing with the first chart.
The resulting native coordinates are
\[
 Z=P^6u_A,\qquad
 a_0=P^{-3}c_q f_0\eta^{-1}u_3,\qquad
 c_0=\eta P^{-2}f_0^{-1}u_2.
 \tag{NS.11}
\]
Again \(q=Za_0c_0\). Its full inverse, after the same actual base
recovery from \((s,q)\), is
\[
 u_A=P^{-6}Z,\quad
 u_2=P^2 f_0\eta^{-1}c_0,\quad
 u_3=P^3\eta(c_qf_0)^{-1}a_0.
 \tag{NS.12}
\]
These are holomorphic inverse maps on the protected outer collar.

On their common localization the complete period NS.5 gives
\[
 u_-^{\infty}=P^2f_0u_Au_3,\qquad
 u_+^{\infty}=P^4f_0^{-1}u_Au_2,\qquad
 z=P^{-6}u_A^{-1}.
 \tag{NS.13}
\]
Its inverse is
\[
 u_A=P^{-6}z^{-1},\quad
 u_2=P^2f_0u_+^{\infty}z,\quad
 u_3=P^4f_0^{-1}u_-^{\infty}z.
 \tag{NS.14}
\]
Substitution into NS.8 and NS.11 yields the exact two-chart map
\[
 z=Z^{-1},\qquad a_\infty=Za_0,\qquad c_\infty=Zc_0,
 \qquad Za_0c_0=z a_\infty c_\infty=q.
 \tag{NS.15}
\]
This proves the regular native tube and its fixed sphere directly
from the two protected lower facets and their specified period.
The original selected height, the physical circle and every analytic
coefficient are carried by the displayed maps.

On the aligned dense torus the same map reads
\[
                    c_0=ZV_{\rm pt},\qquad
                    a_0=\frac{q}{Z^2V_{\rm pt}}.
 \tag{NS.16}
\]
NS.8--15 prove its extension to both native vertex charts. They also
specify its inverse there, rather than using NS.16 alone at a zero
or a pole of \(Z\).

\par

\medskip\noindent\textbf{NS.3. Full point modules, source periods and original determinant}

\par\nobreak

For source band \(j\), retain \(x=Q^jZ\),
\(y=Q^{1-j}/Z\) and the actual point graph
\(\operatorname{Bl}_{(x,P^2)}(xy=P^3p)\). In the I2P point
generators and the zero generator the coefficient is literally
\[
 V_{\rm pt}=-P^{-2j}E A_{\rm pt}
 =P^{-2j}c^{-1}E C_{\rm pt}
 =P^{-2j-2}(x-P^2)E_OC_O.
 \tag{NS.17}
\]
Here \(C_{\rm pt}=-cA_{\rm pt}\), and \(E,E_O\) are the entire
I2P theta products. Multiplication of each expression by \(Z\)
gives \(c_0\); its reciprocal module map, with \(q/Z\), gives
\(a_0\). In the other native chart the coefficients are those of
NS.15. Thus every point/O graph uses the actual I2P Cartier module
and its inverse. The factor \(P^{-2}\) at the zero section is
retained. The use of \(\eta\) in NS.8--12 comes from
\(V_{\rm pt}=\eta V_{\rm reg}\); it cancels in NS.17 only by that
exact stated module equality.

The original point periods are
\[
 h:(Z,V_{\rm pt})\mapsto(P^3pZ,P^2V_{\rm pt}),\qquad
 B:(Z,V_{\rm pt})\mapsto(P^{-3}Z,P^{10}pf_0V_{\rm pt}).
 \tag{NS.18}
\]
Under the map NS.16 their complete native forms are
\[
 \begin{aligned}
 h:&(Z,a_0,c_0)\mapsto
 (P^3pZ,P^{-8}p^{-2}a_0,P^5p c_0),\\
 B:&(Z,a_0,c_0)\mapsto
 (P^{-3}Z,P^{-4}p^{-1}f_0^{-1}a_0,P^7pf_0c_0).
 \end{aligned}\tag{NS.19}
\]
Their products are \(q\), and their inverses divide by the same
full factors at the unchanged \((P,p)\). In particular
\[
 \begin{aligned}
 B^{-1}:&(P^3Z,P^4pf_0a_0,P^{-7}p^{-1}f_0^{-1}c_0),\\
 hB^{-1}:&(P^6pZ,P^{-4}p^{-1}f_0a_0,P^{-2}f_0^{-1}c_0).
 \end{aligned}\tag{NS.20}
\]
These squares include both native charts through NS.15. The
source-band change and point-module factors are exactly I2P.8--9;
the formulas do not erase the selected-height powers at \(p=0\).
Their continuation there is the named period map between the
appropriate affine charts, including NS.13--14.

The original four-form is
\[
 \Omega_{\rm orig}=\frac{2J_{ts}}{a_T}P\,
 dP\wedge dp\wedge d\log X\wedge d\log Y.
 \tag{NS.21}
\]
Using \(X=\eta Z/c_0\), \(Y=Z\), the derivatives of \(\eta\)
wedge to zero against the two original base differentials. NS.2 and
\(q=Za_0c_0\) then give, with the displayed orders of coordinates,
\[
 \Omega_{\rm orig}=\frac2{a_T}\,ds\wedge dZ\wedge da_0\wedge dc_0
 =-\frac2{a_T}\,ds\wedge dz\wedge da_\infty\wedge dc_\infty.
 \tag{NS.22}
\]
The last sign is the actual determinant of NS.15, namely \(-1\).
The full unit \(a_T\) and these signs are retained. At the interior
\(P=0\) divisors the original orders remain the exported
\(2r+z-1\); NS.22 describes the protected outer collar and does
not discard those divisors.

For each of the thirty-one new resolved moving-star charts, the
complete certificate exports the actual matrix and its inverse. Align
the second face by NS.5. In the aligned lattice the normal characters
for the zero native chart are
\[
 m_-=(1,-2,1,1),\qquad m_+=(-1,1,0,0),
 \tag{NS.23}
\]
with coefficient units \(c_q\eta^{-1}f_0^n\) and
\(\eta f_0^{-n}\), respectively, where \(n=0\) on the first
face and \(n=1\) after NS.5. The matrix NS.5 carries all P powers;
these are not included a second time in that unit column. For the
infinity native chart add \((0,1,0,0)\) to both rows. Evaluating
these rows on each executed cone gives the full tangent monomial
coefficient of its two normal derivatives. On the two normal columns
the rows are exactly \((1,0)\) and \((0,1)\), in physical order
\((-1,+1)\). Both P and the native fixed-curve coordinate have
zero normal order; p contains the product of the two normals.
Consequently derivatives of the full base units have zero first
normal term on the fixed surface. This proves that the exported
coefficient matrix is its actual normal derivative.

More explicitly, let \((v_-,v_+)\) be the two actual local normal
coordinates of one executed chart, in physical order \((-1,+1)\),
and let \(e_-,e_+\) be the full tangent monomials and units obtained
from NS.23. In the positive complex coefficient convention the
resolved-to-held-Cox derivative is
\[
 N_\theta^+=G_0(Z)\operatorname{diag}(\overline{e_-},e_+),
 \qquad G_0(Z)=\begin{pmatrix}1&-Z\\\bar Z&1\end{pmatrix}
 \tag{NS.23a}
\]
on the zero native chart. On the other chart replace \(G_0\) by
\(G'_0(z)=\left(\begin{smallmatrix}\bar z&-1\\1&z\end{smallmatrix}\right)\)
and use the two infinity rows. The identity
\(G'_0(1/Z)\operatorname{diag}(\bar Z,Z)=G_0(Z)\)
is the exact normal overlap. The full transported Hermitian norm is
used on its sphere. These maps define \(N_\theta\) on the actual
outer boundary; a tangent monomial with nonzero central divisor
remains that divisor in the interior atlas.

The new quotient source fan and all thirty-one native normal charts
are the complete ones in supporting proof 20 E.  Its thirty compact
rays and two horizontal rays cover the marked half-plane once.  The
twenty-nine recorded ordinary sum-of-neighbor blowdowns, with their
reciprocal source and normal maps, end at the same original three-ray
fan of the native sphere times its P disc.

\par

\medskip\noindent\textbf{NS.4. The coupled endpoint in the actual native tube}

\par\nobreak

Put \(\mathcal A=(c_q/f_0)^4\). The exact point-coordinate endpoint
from SI2 and I2C is
\[
 E:(Z,V_{\rm pt})\mapsto(P^{12}Z,P^{-36}\mathcal A V_{\rm pt}).
 \tag{NS.24}
\]
Conjugation by the two actual native maps proved above gives
\[
 E_0:(Z,a_0,c_0)\mapsto
 (P^{12}Z,P^{12}\mathcal A^{-1}a_0,P^{-24}\mathcal A c_0),
 \tag{NS.25}
\]
\[
 E_\infty:(z,a_\infty,c_\infty)\mapsto
 (P^{-12}z,P^{24}\mathcal A^{-1}a_\infty,
                         P^{-12}\mathcal A c_\infty).
 \tag{NS.26}
\]
Both preserve the literal height, and NS.15 conjugates one to the
other. All inverses use the reciprocal functions at the same source
point. For the positive spinor \((\overline{a_0},c_0)\) the full
normal matrix is
\(\operatorname{diag}(\overline{P}^{12}\overline{\mathcal A}^{-1},
P^{-24}\mathcal A)\). Its determinant phase is therefore
\[
             (P/|P|)^{-36}(\mathcal A/|\mathcal A|)^2.
 \tag{NS.27}
\]
The full constant and based-unit phases remain in this equality.
On the selected fixed locus \(s=-P/u_s(P,0)\), the original positive
source meridian has positive degree one in P. The entire unit
\(\mathcal A(P,0)\) has its chosen logarithm relative to its nonzero
constant. Thus NS.27 has winding \(-36\) in these actual native
normal frames. Its use in a cap requires the exact Cox product below.

\par

\medskip\noindent\textbf{NS.5. Every phase-square edge and the full DZ factor}

\par\nobreak

Retain the relative root coordinate \(\rho\), the actual
\(r_{\rm base}=r_0s_*^4\), and the full SI2 logarithms
\[
 G=\operatorname{Log}(P^{-3}),\qquad
 J=\operatorname{Log}(\lambda_{\rm base}\eta^6t^{-3}).
\]
The point-coordinate product is
\[
 a_\rho=e^{\rho G},\quad
 d_\rho=e^{\rho J+3\pi i\rho(\rho-1)}.
 \tag{NS.28}
\]
Its actual native conjugate, using NS.16, is
\[
 C^0_\rho:(Z,a_0,c_0)\mapsto
 (a_\rho Z,(a_\rho^2d_\rho)^{-1}a_0,a_\rho d_\rho c_0).
 \tag{NS.29}
\]
In the standard native notation its second factor is
\(\beta_\rho=a_\rho d_\rho\), whose full logarithm is
\(\rho(G+J)+3\pi i\rho(\rho-1)\). These formulas preserve q
and have their reciprocal inverses at the same q and source parameter.

Write \(g=e^G\), \(\eta_*=e^J\),
\(\eta_\rho=e^{6\pi i\rho}\eta_*\). Subtracting the full
logarithms proves
\[
 C^0_{\rho+1}\,S(g^{-1},\eta_\rho^{-1})=C^0_\rho,
 \qquad C^0_0=1,\qquad C^0_4=S(g^4,\eta_*^4),
 \tag{NS.30}
\]
where the exact conjugation of a point map \((a,d)\) is
\[
 S(a,d):(Z,a_0,c_0)\mapsto(aZ,(a^2d)^{-1}a_0,ad c_0).
 \tag{NS.31}
\]
The map at four is the inverse of the same endpoint NS.25 at the
identical root endpoint. The map at the initial source cut uses the
full values of G and J there. Both source-cut logarithm increments
are \(-6\pi i\); hence its native mismatch is
\[
 (a,\beta)=(e^{-6\pi i\rho},e^{-12\pi i\rho}).
 \tag{NS.32}
\]
This identifies all four edges and every intermediate integer seam.

Use the right-complex held Cox coefficient
\[
 T_Z=\begin{pmatrix}1&-\overline Z\\Z&1\end{pmatrix},\qquad
 q_H=T_Z(a_0,\overline{c_0})^t.
 \tag{NS.33}
\]
The second chart uses
\(T'_z=\left(\begin{smallmatrix}z&-1\\1&\bar z\end{smallmatrix}\right)\);
the literal identity
\(T'_{1/Z}\operatorname{diag}(Z,\bar Z)=T_Z\)
proves that this is the same held quaternion. Under NS.32 put
\(w=e^{\pi i\rho}\). Direct matrix multiplication gives
\[
 T_{w^{-6}Z}\operatorname{diag}(w^{18},w^{12})
       =\operatorname{diag}(w^{18},w^{12})T_Z.
\]
Left and right multiplication on the right-complex coefficient pair
therefore give exactly
\[
                  q_H\longmapsto w^3q_Hw^{15}.
 \tag{NS.34}
\]
At each integer seam both factors have sign \((-1)^\rho\).
Their full Spin signs are retained. The left adjoint is the actual
base and normal-orbit rotation and extends linearly over the specified
three-ball; the same quaternion arrow supplies its principal Hopf
line map. The right factor is part of that map.

For arbitrary full functions \(a_\rho,d_\rho\), the normal sphere
map is given, with the transported Hermitian norm, by the invertible
matrix
\[
 T_{a_\rho Z}
 \operatorname{diag}((a_\rho^2d_\rho)^{-1},
                    \overline{a_\rho d_\rho})T_Z^{-1}.
 \tag{NS.35}
\]
Its inverse is the inverse product at the recovered point. The full
q-dependent native velocity preserves q. Keeping the prescribed
q-preserving source lift, multiply only that vertical correction by
the protected invariant cutoff. Its flow and backward flow provide
the collar comparison on both vertex charts. This carries every
source derivative and all endpoint jets. The finite cap is outside
the cutoff support, so its unchanged full square is
\[
              \Phi_{\rho+4}q^4=h_\infty^4\Phi_\rho.
 \tag{NS.36}
\]
It retains the root chart \(k\mapsto k+4\), its three chronological
steps, every point correction, sector integer and finite character.

On the diagonal \(\rho=4t\), the native second-factor logarithm
from NS.29 is
\[
       4t(G+J)+3\pi i(16t^2-4t).
\]
The original DZ product has the same source logarithm \(4tG\)
and native second-factor logarithm
\[
                         4t(G+J)+96\pi it^2.
\]
This follows by adding the unchanged source logarithm to the actual
point-line DZ logarithm \(4tJ+96\pi it^2\). Thus the entire G and J,
including all analytic units and their constant values, cancel
pointwise in the comparison. The remaining directed map is precisely
the physical scalar
\[
 F(t)=e^{3\pi i(16t^2+4t)}
     =e^{48\pi it(t-1)}e^{60\pi it},\qquad
 (a_0,c_0)\mapsto(F(t)^{-1}a_0,F(t)c_0).
 \tag{NS.37}
\]
The first logarithm has zero values at both endpoints. The second is
the retained positive physical winding thirty. In the held Cox
quaternion it is the inverse right phase. The explicit intertwining
\(S\delta_F=\delta_F S\) also follows directly from NS.16, so this
loop is carried on the same physical line and both finite-cap
interfaces. No contraction of its nonzero winding is used.

\par

\medskip\noindent\textbf{NS.6. Exact scope bridge for I2C, PPH and RB}

\par\nobreak

Keep I2C at SHA256
\path{e9b6b65b1ef64c7c086fca00bde3c937bb34d6b07d0b2e7ce8b422b2c80cd072}
and RB at SHA256
\path{a8f69a012bc2ec528d2da2181ae6049803b41d74710e3d7e262b4974a90fdb79}.
Their point-coordinate endpoint NS.24 is exact. On its dense common
torus the old coefficients satisfy
\[
 c_N=V_{\rm pt},\quad a_N=q/(ZV_{\rm pt}),\qquad
 a_0=Z^{-1}a_N,\quad c_0=Zc_N.
 \tag{NS.38}
\]
This is the exact morphism relating the two coefficient systems on
that localization. At the regular native point \(Z=0,a_0=0,c_0=1\),
the old expression \(c_N=c_0/Z\) has a pole. Therefore that displayed
old point-coordinate map does not extend to this native vertex chart.
NS.8--15 supply the actual regular extension selected by the protected
lower atlas.

In positive normal coefficients the common-locus comparison is
\(S_+(Z)=\operatorname{diag}(\bar Z^{-1},Z)\). Its directed endpoint
square is
\[
 S_+(P^{12}Z)
 \operatorname{diag}(\bar P^{24}\bar{\mathcal A}^{-1},
                         P^{-36}\mathcal A)
 S_+(Z)^{-1}
 =\operatorname{diag}(\bar P^{12}\bar{\mathcal A}^{-1},
                         P^{-24}\mathcal A).
 \tag{NS.39}
\]
This proves the change from the old phase exponent \(-60\) to the
actual native exponent \(-36\) without declaring the intervening
meromorphic character an extending gauge. Likewise NS.29--35 are
the full native replacement for the old point-coordinate interpretation
of the SI2.36/PPH.2 Cox product; its right phase becomes fifteen.
NS.37 proves the complete DZ comparison through that exact map.

For RB.25a the point endpoint stays NS.24, while its native-tube
assertion uses NS.25--27 and its cut/Cox maps use NS.29--36. The disjoint
fourteen-factor source pieces, their source/line/cap transports, and
the intrinsic cube and relative pairings in RB.31--34 are retained.
Their comparison with the moving boundary is now made through this
actual native tube. The compact resolved normal bundle, its cap
clutch and the global outer characteristic ledger remain subsequent
calculations. In particular this proof supplies their coordinate and
orientation input and assigns none of their characteristic numbers.


% End expanded verified cp36 input


% Expanded verified cp36 input: notes/part3_e6_selected_i3_m2_return_support_14.tex
\par\bigskip\noindent\textbf{P3K.111 supporting proof 14: Directed native Hopf cap and five-dimensional bundle bordism}\par\nobreak
\par

\medskip\noindent\textbf{Selected-I3 \(M_{2}\) N12: the actual native cap and its directed bundle bordism}

\par\nobreak

This construction uses the thirty-one actual native-star charts of supporting proof 20 and the protected native maps NS.1-\/-6,
the original SI2 height and point/O modules, the cubic cap with norm
one, its fixed finite character, the complete annular reference, and
the coupled scalar \(\lambda=\lambda_0 s^{12}\) with its chosen root \(r_{0} s^{4}\).
The source and period maps in NS are retained literally. The following
cap attaches to that regular native tube. The older I2C source remains
preserved at its stated hash; NS.38-\/-39 give the exact localization map
and its pole at the native vertex. No Chern number is inferred from
the older coefficient system.

\par

\medskip\noindent\textbf{NC.1. The full native endpoint, its inverse, and the cut collar}

\par\nobreak

Write \(q=q_{\rm old}\) and keep the complete functions

\[
 Q=P^3p,\quad q=Pp c_q,\quad
 s=P(p-1)/u_s,\quad t=Pu_t,\quad
 f_0=(p-1)^{12}\frac{\lambda_0c_q\eta^6}{u_s^{12}u_t^3},
 \qquad {\cal A}=(c_q/f_0)^4.                       \tag{NC.1}
\]

Here \(\eta\) includes the actual theta derivative and \(\eta(0)=-1\). All
functions and their reciprocals are evaluated at the original source
point recovered by the inverse of NS.2. On the protected outer source
collar P is nonzero and each displayed unit is nonzero. The regular
native charts proved in NS satisfy

\[
 q=Za_0c_0=z a_\infty c_\infty,\qquad
 z=Z^{-1},\quad a_\infty=Za_0,\quad c_\infty=Zc_0.
                                                        \tag{NC.2}
\]

The actual endpoint and its inverse in these charts are

\[
 E_0=(P^{12}Z,P^{12}{\cal A}^{-1}a_0,
                         P^{-24}{\cal A}c_0),\qquad
 E_0^{-1}=(P^{-12}Z,P^{-12}{\cal A}a_0,
                         P^{24}{\cal A}^{-1}c_0),       \tag{NC.3}
\]
\[
 E_\infty=(P^{-12}z,P^{24}{\cal A}^{-1}a_\infty,
                            P^{-12}{\cal A}c_\infty).
                                                        \tag{NC.4}
\]

These formulas fix s and q. Multiplying their three coordinate
factors gives one; composing either direction with NC.2 proves the
whole overlap square, including the full analytic functions. Their
inverses use the reciprocal factors at the same recovered source
point. NS.8-\/-14 already prove their regular extension at both native
vertices, so the dense expressions \(c_{0}=Z V_{\rm pt}\) and \(a_{0}=q/(Z^{2} V_{\rm pt})\) are
used here with those extending chart maps.

Cut the positive original s circle. Choose its continuous logarithm
of P and the based logarithm of the unit A on this cut. For a smooth
\(\chi(\theta)\) constant zero on the initial collar and constant one on
the terminal collar, set

\[
 \alpha_\theta=\exp[-\chi(\theta)\operatorname{Log}(P^{12})],
 \qquad
 \beta_\theta=\exp[-\chi(\theta)
                          \operatorname{Log}(P^{-24}{\cal A})].
                                                        \tag{NC.5}
\]

The logarithms include their constant values and their entire source
dependence. Define

\[
 C_\theta^0(Z,a_0,c_0)=
 (\alpha_\theta Z,(\alpha_\theta\beta_\theta)^{-1}a_0,
                                             \beta_\theta c_0),
                                                        \tag{NC.6}
\]
\[
 C_\theta^\infty(z,a_\infty,c_\infty)=
 (\alpha_\theta^{-1}z,\beta_\theta^{-1}a_\infty,
                              \alpha_\theta\beta_\theta c_\infty).
                                                        \tag{NC.7}
\]

NC.2 intertwines these maps exactly. Both fix q; hence the inverse
recovers that same q, recovers P,p from s,q, and divides by precisely
the written factors. This proves both compositions as maps of the
original source and the actual line modules. The endpoint values are
\(C_{0}=1\) and \(C_{1}=E^{-1}\). On the following cut define \(C_{\theta+1} E=C_\theta\).
Because \(\chi\) is constant on collars, this is an equality of full
collar maps and all their endpoint derivatives.

The real vector field of NC.6 has holomorphic-coordinate part

\[
 (\partial_\theta\log\alpha)Z\partial_Z
 -(\partial_\theta\log(\alpha\beta))a_0\partial_{a_0}
 +(\partial_\theta\log\beta)c_0\partial_{c_0}.            \tag{NC.8}
\]

It annihilates q by direct differentiation. NC.7 gives the same field
on the other vertex chart. Add the prescribed parameter-horizontal
lift fixing q, and multiply only this vertical correction by a smooth
physical-circle-invariant cutoff equal to one on the protected tube
and zero off a slightly larger tube. Its support is disjoint from the
good cap, its annular interface and the outer source boundary. The
real field is smooth on the compact cut fibres; its ordinary flow and
backward flow exist for the full compact parameter interval. They
give the exact NC.6 map on the smaller tube and the prescribed
unchanged comparison on each protected interface. The field commutes
with the physical action because its coefficients depend on the
invariant s,q and its coordinate weights are diagonal.

Thus this is an actual interval identification of the family and its
line bundles. The finite cap retains the original full square

\[
                     \Phi_{\rho+4}q^4=h_\infty^4\Phi_\rho.
                                                        \tag{NC.9}
\]

The root chart shift by four, its three chronological comparisons,
the point correction at each step, every sector integer and the
finite character are all the original SI2/NS arrows, since the flow
is the identity on that interface. The physical scalar of winding
thirty remains the full NS.37 comparison. In particular this collar
construction preserves that scalar and does not remove its phase.

\par

\medskip\noindent\textbf{NC.2. The whole Cox map and its extension over the actual ball}

\par\nobreak

Use the right-complex coefficient and the other sphere chart

\[
 T_Z=\begin{pmatrix}1&-\overline Z\\ Z&1\end{pmatrix},\quad
 q_H=T_Z(a_0,\overline{c_0})^t,\qquad
 T'_z=\begin{pmatrix}z&-1\\1&\overline z\end{pmatrix}.
                                                        \tag{NC.10}
\]

The literal identity
\(T'_{1/Z} \operatorname{diag}(Z,\overline{Z})=T_Z\) proves that both expressions give the same
held quaternion. The physical circle acts by its inverse right
phase. In this coefficient the entire first-normal map of NC.6 is

\[
 G_\theta(Z)=T_{\alpha_\theta Z}
 \operatorname{diag}((\alpha_\theta\beta_\theta)^{-1},
                                      \overline{\beta_\theta})
 T_Z^{-1}.                                             \tag{NC.11}
\]

In NC.11 \(\alpha\) and \(\beta\) are their actual restrictions to \(q=0\). Their
normal derivatives vanish at the fixed surface: \(q=Za0c_{0}\) has zero
first normal differential there and P,p are functions of s,q.
Thus NC.11 is the full first-normal derivative, including the
original source units. Its determinant is nonzero and its inverse
is the reversed matrix product at the recovered point. On the other
sphere chart the NC.10 identity and NC.7 give exactly this same
matrix. It is smooth at infinity; the other chart gives there the
diagonal value \(\operatorname{diag}(\overline{\alpha \beta},\beta^{-1})\). Sphere normalization
uses the transported Hermitian norm. The inverse applies \(G^{-1}\) and
then normalizes, so both sphere compositions are identities.

Represent the closed three-ball as the trace-one positive
semidefinite Hermitian matrices Y. Its boundary point is

\[
 Y_Z=\frac1{1+|Z|^2}
       \begin{pmatrix}|Z|^2&Z\\\overline Z&1\end{pmatrix}.
                                                        \tag{NC.12}
\]

Put \(D_\theta=\operatorname{diag}(\alpha_\theta|_{q=0},1)\). The actual ball map is

\[
 \Psi_{D_\theta}(Y)=
 \frac{D_\theta YD_\theta^*}
      {\operatorname{tr}(D_\theta YD_\theta^*)}.           \tag{NC.13}
\]

The denominator is positive on the whole ball: D is invertible and
Y is nonzero positive semidefinite. Its inverse is \(\Psi_{D^{-1}}\), as
direct substitution shows. It preserves rank, and on rank-one \(Y_Z\)
it is \(Y_{\alpha Z}\). Composition is \(\Psi_D \Psi_F=\Psi_{DF}\), including
the normalization factors. Hence NC.13 has exactly the boundary
source map and its directed endpoint seam.

Write \(Y=I/2+r H\) with \(\operatorname{tr} H=0\) and H on the unit direction sphere. Let
\(\tau(r)\) be smooth, zero near \(r=0\) and one near the boundary. In NC.11
replace \(\alpha\),\(\beta\) by

\[
 \exp(\tau(r)\log\alpha),\qquad
 \exp(\tau(r)\log\beta),                                \tag{NC.14}
\]

using the fixed cut logarithms and the sphere direction H for Z.
This defines \(\widehat G_\theta(Y)\). It is independent of direction near
the centre because it is the identity there. Away from the centre,
the exact other-chart identity in NC.10 proves smoothness across
the infinity direction. Every factor is invertible; the inverse
product is smooth by the same statements. The actual full cap map is

\[
 \widehat C_\theta(Y,q_H)=
 \left(\Psi_{D_\theta}(Y),
 \frac{\widehat G_\theta(Y)q_H}
      {\|\widehat G_\theta(Y)q_H\|}\right).              \tag{NC.15}
\]

Its inverse first recovers Y by \(\Psi_{D^{-1}}\) and then uses the inverse
matrix at that recovered Y. This order proves the two compositions
even though \(\widehat G\) varies with Y. It restricts to NC.6-\/-11 on the
common sphere boundary. The forward cap seam is the actual map

\[
                    \widehat E=\widehat C_1^{-1}\widehat C_0.
                                                        \tag{NC.16}
\]

It restricts to NC.3-\/-4 with the full source and line functions.
The equality \(\widehat C_{1} \widehat E=\widehat C_{0}\) proves that \(\widehat C\) identifies this
specific ball/sphere mapping torus with the product cap. The inverse
is the inverse just constructed on the cut and agrees at the seam.
Thus the construction fixes the cap product, its source ball and
its normal sphere, rather than merely selecting a determinant degree.

\par

\medskip\noindent\textbf{NC.3. The executed normal derivative and the exact cap product}

\par\nobreak

On each of the thirty-one actual new resolved moving-star charts let
\((v_{\rm minus},v_{\rm plus})\) be the local normal coordinates in physical order
\((-1,+1)\). NS.23 evaluates the original native coefficient characters
on the actual executed cone and obtains their tangent monomials and
full units \(e_{\rm minus},e_{\rm plus}\). Their normal exponents are exactly the
two standard basis vectors, so the map is the actual derivative.
In particular the point/O formulas NS.17 and the full \(T_{0}\) period are
carried in those coefficients. No source band or zero-section factor
is removed. In positive complex coefficients the exact map is

\[
 N_\theta^+=G_0(Z)
          \operatorname{diag}(\overline{e_-},e_+),\qquad
 G_0(Z)=\begin{pmatrix}1&-Z\\\overline Z&1\end{pmatrix},
                                                        \tag{NC.17}
\]

and on the other native chart it uses

\[
 G_1(z)=\begin{pmatrix}\overline z&-1\\1&z\end{pmatrix},
 \qquad G_1(1/Z)\operatorname{diag}(\overline Z,Z)=G_0(Z).
                                                        \tag{NC.18}
\]

The two charts and the NS normal-coordinate inverses prove agreement
on every common localization. The meromorphic tangent monomials at
interior divisors keep their actual divisor orders; NC.17 is used as
an invertible boundary map where P is nonzero.

Let \(N_\theta\) be the corresponding right-complex map before conjugating
its coefficient pair. The map from the resolved normal sphere into
the actual cut family is \(B_\theta=C_\theta^{-1} N_\theta\). Composing the
specified cap map, in its actual order and at its actual source point,
gives

\[
                     \widehat C_\theta B_\theta=N_\theta.
                                                        \tag{NC.19}
\]

Its inverse equality is \(B_\theta^{-1} \widehat C_\theta^{-1}=N_\theta^{-1}\). The
base equality follows from NC.13, and the full normal equality
follows from NC.11 and NC.15. Conjugating the two right-complex
coefficient columns gives the same identity for the positive complex
rank-two bundle NC.17. Thus the completed bundle has precisely the
periodic full \(N_\theta\) as its clutch in the product cap of NC.16.
Every P power and unit of C, B and \(\widehat C\) has occurred in a displayed
map and its inverse before this exact cancellation.

\par

\medskip\noindent\textbf{NC.4. The actual oriented replacement bordism}

\par\nobreak

The fixed surface in the cap product is \(C_{0}=-D^3_Z\times S^1_P\). Take
\(C=\mathbb P^1_Z\times D^2_P\) with its product orientation. The oriented product
boundary is

\[
 \partial(D^3_Z\times D^2_P)=C+C_0.                    \tag{NC.20}
\]

Round the corner in its two radial coordinates. Attach its rounded
facet to the boundary collar of the resolved fixed surface times an
interval. Carry the exact map N of NC.19 constantly across that
interval, with its full inverse. Over \(D^3\times D^2\) the positive
rank-two bundle is trivial in the conjugated held Cox frame. This
constructs a compact smooth five-dimensional bundle bordism from
the actual cap-completed fixed surface and normal bundle to the
closed replacement \(F_{\rm star}-C\) and its attached bundle. Equality of
their \(c_2\) and \(c_1^2\) follows by extending a Hermitian connection
over this explicit bundle bordism and applying Stokes to its closed
characteristic forms.

On the C face the two ordered line subspaces in that frame are
\(C(1,\overline{Z})\) and \(C(-Z,1)\). NC.18 gives their coefficient transitions
\(\overline{Z}\) and \(Z\), so their first Chern numbers on the complex Z sphere
are respectively +1 and -1. Retain these exact subspaces on the
replacement face. The diagonal derivative coefficients NC.17
preserve them and therefore glue the two positive normal lines.

For a P-disc radius epsilon, the opposite face -C has the coordinate
\(\eta=\overline{\xi}/\epsilon^{2}\) on its abstract disc. This real map has negative
determinant and equals 1/P on the common boundary \(\xi=P\). Thus -C
has its \(\eta\)-complex orientation. The closed replacement is the
actual twenty-nine-fold blowup of \(\mathbb P^1_Z\times \mathbb P^1_{\rm parameter}\) given by the
new source fan and its recorded chronological coordinate blowup maps in supporting proof 20.
All twenty-nine centres lie in the original P disc. The meromorphic normal
dual sections and their chart inverses determine the divisor lines
on these blowups; they are retained for the separate intersection
calculation.

The sign of each parameter degree is fixed by this actual attachment.
For the boundary coefficient arrow \(z_{\rm cap}=P^n z_F\), the meromorphic
section with coefficient one on the original disc has coefficient
\(\eta^-n\) on the new cap chart. Its divisor has degree -n. For a
conjugated P coefficient, \(\overline{P}=P^{-1} |P|^{2}\) on the boundary; the
positive modulus has its real logarithm and is retained by the
explicit nonzero homotopy. This proves the sign convention directly
for the arrows in NC.17.

For completeness every unit \(u(P)\) used in these normal coefficients
has its actual nonzero value \(u(0)\) and a unique based logarithm \(L(P)\)
on the source disc: \(u(P)=u(0)exp L(P)\). The coefficient path
\(u(0)exp(t L(P))\) and its reciprocal are invertible throughout. On
the cap it extends by the radial cutoff of NC.14 in the global
projector onto the corresponding NC.18 line subspace. This is smooth
at the ball centre and has the reciprocal exponential as inverse.
The constant value remains in the attached frame. Therefore these
particular full units have zero parameter degree by an explicit
extending bundle map. No meromorphic P or Z character is included
among those units; their divisor orders stay in the thirty-one-chart
calculation of supporting proof 20.

\par

\medskip\noindent\textbf{NC.5. Relation to the preserved point-coordinate calculation}

\par\nobreak

The dense localization map from the preserved I2C coefficients is
\(a_{0}=Z^{-1} a_N\) and \(c_{0}=Z c_N\). Its normal positive matrix is
\(\operatorname{diag}(\overline{Z}^{-1},Z)\). At the regular point \(Z=0,a_{0}=0,c_{0}=1\) the old inverse
\(c_N=c_{0}/Z\) has a pole. The exact endpoint conjugation is NS.39, giving
the actual normal exponent pair \((-12,-24)\) in positive phases and
their sum \(-36\). NC.3-\/-19 use this actual extending native pair at
both vertices. The original source endpoint, full finite-cap square,
the NS phase-square maps and the physical winding thirty retain
their original arrows. This specifies precisely which local maps are
used for the new cap and which old formulas remain statements on
the common torus localization.


% End expanded verified cp36 input


% Expanded verified cp36 input: notes/part3_e6_selected_i3_m2_return_support_15.tex
\par\bigskip\noindent\textbf{P3K.111 supporting proof 15: The actual fixed-source reference and ordered native bases}\par\nobreak

\medskip\noindent\textbf{MN.3. The actual RB moving reference and its own cap}

\par\nobreak

Keep the full fixed-source selected functions of RB:
\[
 Q=\kappa q,\quad \alpha=P^2,\quad g=P^{-3},\quad
 \Gamma=\lambda_0\eta^6t^{-3}.
\]
Here the original source parameter equals \(s_*\ne0\), and the
independent scalar-disc coordinate is \(\sigma\), with suspension
\(\lambda_0\sigma^{12}\). The complete source and line maps
are exactly RB.2 and RB.16--17, with their genuine Cartier frames.

Take the y-factor chart of RB.13 with \(i=j=0\). Then
\[
 Z=x_0,\quad q=Zu_yv_y,\quad
 V_{\rm pt}=\sigma^{12}Z u_y.
\]
Its regular native coefficients are
\[
 c_R=Z^{-1}V_{\rm pt}=\sigma^{12}u_y,\qquad
 a_R=q/V_{\rm pt}=\sigma^{-12}v_y,
 \qquad q=Za_Rc_R.
 \tag{MN.9}
\]
The powers in MN.9 are boundary coefficient maps at nonzero scalar
radius. The resolved normal variables \(u_y,v_y\) themselves
remain regular at \(\sigma=0\).

On the x-factor chart of band \(i=-1,j=0\), put
\[
 z=\kappa\bar y_{-1}=Z^{-1},\qquad
 c_{R,\infty}=\sigma^{12}u_x,\qquad
 a_{R,\infty}=\kappa^{-1}\sigma^{-12}v_x.
 \tag{MN.10}
\]
The exact RB.17a overlap, including its inverse, gives
\[
 u_y'=Z^{-1}u_x,\qquad v_y'=(\kappa Z)^{-1}v_x,
\qquad
 (a_{R,\infty},c_{R,\infty})=Z(a_R,c_R).
 \tag{MN.11}
\]
Both native products are q. The inverse recovers q first, then
the same actual \(P(q),\kappa(q),\Gamma(q),\eta(q)\) and the
displayed reciprocal coefficients. Thus this is the actual regular
reference native sphere, derived before any numerical comparison.

Changing the source or chain representative uses the full RB maps.
In the genuine point coefficient the actual period is
\[
 B_\sigma:(Z,V_{\rm pt})\longmapsto
       (gZ,\sigma^{12}q\Gamma V_{\rm pt}).
\]
For the target label \(j+1\), division by
\((\sigma^{12}q)^{j+1}\) gives \(U_{j+1}'=\Gamma U_j\),
and the reciprocal coordinate gives \(W_{j+1}'=\Gamma^{-1}W_j\).
Thus the original q factor cancels only through this literal chain
arrow. At \(q=0\) its two resolved continuations are exactly RB.17.
The h map has \(i\mapsto i+1\) and line coefficient \(\alpha\).
The B map has \(j\mapsto j+1\), source coefficients
\((g,g^{-1})\), and the normals in RB.17. Their inverses carry
the same full \(\alpha,g,\Gamma,\kappa\) at the inverse source
point. On the moving fixed surface \(q=0\) these are nonzero constants
on the scalar disc. Thus their actual logarithms extend over that
disc, while the powers \(\sigma^{-12},\sigma^{12}\) in
MN.9--10 remain unchanged. The twelve other scalar-source pieces
are disjoint by the explicit \(v_{12}\ne0\) overlap of RB.15;
none of them is included in this moving surface.

For completeness its own cap is the NC.6--19 construction applied
to MN.9--11. The actual point endpoint is
\[
 (Z,V_{\rm pt})\longmapsto
        (g^{-4}Z,\sigma^{-48}\Gamma^{-4}V_{\rm pt}).
\]
Its native reference factors are
\[
 Z\mapsto g^{-4}Z,\quad
 a_R\mapsto\sigma^{48}\Gamma^4a_R,\quad
 c_R\mapsto g^4\sigma^{-48}\Gamma^{-4}c_R.
 \tag{MN.12}
\]
Their product is one. Use the full logarithms of
\(g^{-4}\) and \(g^4\sigma^{-48}\Gamma^{-4}\) in the two
NC cut factors. The two native chart equations, the q-preserving
velocity and its inverse, the positive Hermitian-ball map and the
whole held-Cox normal matrix are exactly those NC formulas with
these explicitly substituted functions. Their cap seam is
\(\widehat C_1^{-1}\widehat C_0\), and their resolved boundary
map obeys the same full identity \(\widehat C B=N_R\).
The actual original finite-cap square remains
\(\Phi_{\rho+4}q^4=h_\infty^4\Phi_\rho\), with all root
characters and sector integers, because the cut field is supported
in the selected tube. This constructs the reference\textquotesingle s own cap.

Its positive normal coefficient map is
\(G_0\operatorname{diag}(\overline{\sigma^{-12}},\sigma^{12})\)
in the zero native chart. The infinity chart has the additional
full nonzero coefficient \(\kappa(0)^{-1}\) on the negative line.
That constant and the full point/O units extend by the explicit NC
unit map. Applying the same opposite parameter-disc bundle bordism,
the closed reference replacement is
\(\mathbb P^1\times\mathbb P^1\), and its two ordered line
classes are exactly
\[
                    x_R=h-12u,\qquad y_R=-h-12u.
 \tag{MN.13}
\]
In particular
\[
 x_R^2=-24,\quad x_Ry_R=0,\quad y_R^2=24,\qquad
 c_1(E_R)^2=0,\quad c_2(E_R)=0,\quad p_1(TF_R)=0.
 \tag{MN.14}
\]
The same explicit bundle bordism transfers these three characteristic
numbers to the actual reference chart with its own specified cap. It has not been substituted for the collision
cap by a dense character map.

\par

\medskip\noindent\textbf{MN.4. The exact ordered bases and the remaining whole-pair map}

\par\nobreak

In the original \((\operatorname{ord}Z,\operatorname{ord}V_{\rm pt})\)
plane at selected height one, the physical edge is
\(0\to(0,1)\). After the literal collision \(T_0\) alignment,
its opposite triangle vertices are \((1,-1)\) and \((-1,2)\).
The reference charts MN.9--10 give vertices \((1,1)\) and
\((-1,0)\). Thus the full ordered matrices are
\[
 P_N=\begin{pmatrix}1&-1\\-1&2\end{pmatrix},\qquad
 P_R=\begin{pmatrix}1&-1\\1&0\end{pmatrix},\qquad
 P_N(1,1)^t=P_R(1,1)^t=(0,1)^t.
 \tag{MN.15}
\]
Both determinants are one. Their dense coefficient relation is
\[
 (Z,a_0,c_0)=(Z,Z^{-2}a_R,Z^2c_R),\qquad
 \begin{pmatrix}1&0\\2&1\end{pmatrix}P_N=P_R.
 \tag{MN.16}
\]
At the regular reference vertex \(Z=0,a_R=1,c_R=0\), the displayed
\(a_0\) has a pole. Therefore MN.16 does not itself supply the
regular vertex or cap identification. The original point/O arrows
cannot remove that pole on this selected node collar: their
numerator/denominator divisors already cancel in the actual frames
NS.17, and their remaining local factors are nonzero.

The full point product \((a_\rho,d_\rho)\) acts in the reference
native chart by
\[
 (Z,a_R,c_R)\mapsto
       (a_\rho Z,d_\rho^{-1}a_R,(d_\rho/a_\rho)c_R).
 \tag{MN.17}
\]
Its source-cut mismatch has \(a=e^{-6\pi i\rho}\) and second
native factor one. Its held quaternion is consequently
\(q_H\mapsto e^{3\pi i\rho}q_He^{3\pi i\rho}\), with both
odd signs retained. The collision right factor in NS.34 is fifteen;
the exact dense conjugation MN.16 accounts for their difference.
Both complete diagonal products retain the physical DZ factor
NS.37, by the same pointwise subtraction of the full logarithms.

The full RFB source, Cox, line, reference and cap continuation below
compares the actual complete references in their specified native fans.
The collision own-cap bundle used in the present assembly is the new
thirty-one-chart bundle computed in supporting proof 20 E; MN.14 is
its unchanged fixed-source own reference.

% End expanded verified cp36 input


% Expanded verified cp36 input: notes/part3_e6_selected_i3_m2_return_support_16.tex
\par\bigskip\noindent\textbf{P3K.111 supporting proof 16: Twelve full lower/reference doubles with original cube}\par\nobreak
\par

\medskip\noindent\textbf{Selected-I3 \(M_{2}\): all twelve lower source surfaces and their full reference doubles}

\par\nobreak

Working proof, 18 November 2026. The source is
\(y^2=x^3+r_{\rm src}x^2+tx+t^2\), with \(G=(0,t)\),
\(M_2=\mathcal O([2G]-[O])\), chronological translation \(-3G\),
the good cubic cap \((m,n)=(3,1)\), and its prescribed finite character.
The source parameter is \(s=1/4-r_{\rm src}\). We retain SI2\textquotesingle s coupled
scalar \(\lambda=\lambda_0s^{12}\), scalar root \(r_0s^4\), original
selected height \(q=t-t_c(s)\), and complete annular cochain.

The twelve surfaces used here are the complete minimal three-chart
lower source stars of the new one endpoint in supporting proof 20.
They are closed against the actual twelve RB scalar-source pieces by
the full source, line, cube-frame and cap maps retained below.  Their
normal pairings and physical pair are evaluated in supporting proof 20 E.

\par

\medskip\noindent\textbf{IL.1. Original source, point line, normal coordinates and every period}

\par\nobreak

Use the full functions from SI2 and I2O:
\[
 Q=P^2R=P^3p,\quad p=1+v,\quad v=su_s/P,\quad
 t=Pu_t,\quad q=R c_q,\quad
 c_{12}=\frac{\lambda_0c_q\eta_2^6}{u_s^{12}u_t^3}.
 \tag{IL.2}
\]
Here
\(\eta_2=a_T\Theta'_Q(1)\Theta_Q(P^2)/\Theta_Q(P)^2\).
All products and inverse coefficients retain these entire functions
on their specified source domains. With \(X=Z^{-1}\) and
\(Y=V_{\rm reg}/Z\), the full two periods are
\[
 L:(X,Y)\longmapsto(P^3pX,PpY),\qquad
 T:(X,Y)\longmapsto(p^{-1}X,P^{12}v^{12}c_{12}Y).
 \tag{IL.3}
\]
They fix \(P,v\); their inverses use the reciprocal complete factors.
The two actual original line coefficients are
\[
 V_{\rm pt}=\eta_2Y/X,\qquad
 v_2=\frac{\eta_2P^2Y}{Xf_{P^2}(pP^3/X)}.
 \tag{IL.4}
\]
The latter is a morphism in the original point/O Cartier modules.
Its reciprocal is \(Y=\eta_2^{-1}P^{-2}X
f_{P^2}(pP^3/X)v_2\), at the same recovered source point.

In the ordered lattice
\((\operatorname{ord}X,\operatorname{ord}Y,
\operatorname{ord}P,\operatorname{ord}v)\), the twelve minimal
physical fixed faces are
\[
 E_l=(0,l,0,1),\quad E_{l+1}=(0,l+1,0,1),
                  \qquad 0\leq l\leq11.
 \tag{IL.5}
\]
Their stabilizer is generated by the full \(L\) arrow. The independent
\(T\) arrow replaces \(l\) by \(l+12\), with its entire line multiplier;
the twelve residue labels therefore remain twelve separate surfaces.
The physical cocharacter is \((0,1,0,0)\). In a chart ordered by
\((E_l,E_{l+1},r_1,r_2)\), its weights are exactly \((-1,+1,0,0)\).
The source projection is the integral map
\((x,y,r,z)\mapsto(x,r)\). Both source functions \(X,P\) have zero
normal order.

Write the meromorphic parent normal coefficients as
\[
 \alpha_l=Yv^{-l},\qquad \beta_l=v^{l+1}/Y,\qquad
                     \alpha_l\beta_l=v.
 \tag{IL.6}
\]
The physical weights of \(\alpha_l,\beta_l\) are \(+1,-1\).
For a full ray \((x,y,r,z)\), their two source exponents are
\(y-lz,(l+1)z-y\). On every executed chart these are the two rows
of the full monomial map to the parent coefficients, with normal entries respectively
\((0,1)\) and \((1,0)\). These rows are retained on all thirty-six new lower charts.
In particular the normals being calculated are the original total-line
normal bundles restricted to the fixed surfaces.

The full normal period squares are
\[
 L:(\alpha_l,\beta_l)\longmapsto
       (Pp\,\alpha_l,(Pp)^{-1}\beta_l),
 \tag{IL.7}
\]
\[
 T:(\alpha_l,\beta_l,l)\longmapsto
   (P^{12}c_{12}\alpha_l,P^{-12}c_{12}^{-1}\beta_l,l+12).
 \tag{IL.8}
\]
The saved period-unit key \(f_1\) denotes the full \(c_{12}\)
of IL.2 in this one-end execution. Its character columns are
\((0,0,0,0)\) for \(L\) and \((0,1,0,0)\) for \(T\);
the \(p\) columns are \((1,1,0,0)\) and \((-1,0,0,0)\).
Every executed lower endpoint arrow carries those columns through
its full source and target inverse matrices, with its reciprocal
unit vector also retained.
The \(p\) and \(c_{12}\) factors are retained before restriction to
the fixed surface. A derivative of a source unit multiplying a
vanishing normal coefficient contributes zero additional normal
term at that surface. Also \(v=\alpha_l\beta_l\) has zero first
normal derivative there. Thus the actual normal derivatives are
precisely the displayed diagonal coefficients, together with the
full recorded source derivative.

\par

The new lower stars have the complete minimal source cycle, with
normal degree vectors \((-1,0,1)\) and \((-1,1,0)\).
On the marked two-component relative basis set
\[
 I=\begin{pmatrix}-2&1\\1&-2\end{pmatrix}.
\]
The actual relative lifts are \(I^{-1}(0,1)^t\) and
\(I^{-1}(1,0)^t\), giving the pairings
\((-2/3,-1/3,-2/3)\).  All three original source charts, their
full period closure and normal rows are in the new certificate.
The following source and frame morphisms are retained in full.

\medskip\noindent\textbf{IL.4. Exact point/O and source maps to the fixed reference}

\par\nobreak

The source bands \(x=Q^jZ,y=Q^{1-j}/Z\), with \(xy=P^2R\),
carry the proper graph of \((x,P^2)\):
\[
 x=P^2a,\quad ay=R;\qquad xc=P^2,\quad y=Rc.
 \tag{IL.20}
\]
The second chart is resolved by
\[
 P=uw,\ x=u,\ c=uw^2,\ y=Ruw^2;\qquad
 P=uw,\ x=u^2w,\ c=w,\ y=Rw.
 \tag{IL.21}
\]
These arrows preserve \(P,R\), and hence the original \(s,t\).
Their reciprocal overlaps are the actual blowup overlaps. At
\(s=0\), \(R=P\), the source coordinates
\[
       x_i=X/P^i,\qquad y_i=P^{i+1}/X,\qquad x_iy_i=P
 \tag{IL.22}
\]
identify their three node charts with the marked minimal
\(I_3\) source. The point section lies on \(C_1\): in IL.20
it is \(a=1\), so \(X=P\) after the actual \(L\) reindex.
The zero section lies on \(C_0\), at \(X=1\); the alternate
representative at \(X=P^3\) is carried back by the same complete
source and line period.

For \(b=P^2,e=R\), the full theta expression has the factorization
\[
 b f_b(x)=(b-x)E_{b,e}(x,y),
 \tag{IL.23}
\]
where \(E_{b,e}\) is I2P.3\textquotesingle s complete normally convergent product
on this node chart. In the original point modules
\[
 v_2=(a-1)^{-1}A=(c-1)^{-1}C,\quad C=-cA,
\]
\[
 V_{\rm reg}=-\eta_2^{-1}P^{-2j}E_{b,e}A
     =\eta_2^{-1}P^{-2j}c^{-1}E_{b,e}C .
 \tag{IL.24}
\]
At the zero section the actual generator is \(x-1\), and
\[
 V_{\rm reg}=\eta_2^{-1}P^{-2j-2}(x-P^2)E_OC_O.
 \tag{IL.25}
\]
The full \(E_O\), the additional \(P^{-2}\), every reciprocal,
and the original Miller arrow \(r_2=f_2^{-1}\) remain in these
formulas. Their point/O horizontal factors are interpreted in
their indicated line modules through coincident section points.

On the retained boundary of a fixed small original \(t\)-disc,
the directed normal map to the scalar-source chain is
\[
 A_l=H_{+,l}\alpha_l,\quad B_l=H_{-,l}\beta_l,\qquad
 H_{+,l}=\eta_2u_s^lP^{-l}X^{-1},\quad
 H_{-,l}=\eta_2^{-1}XP^{l+1}u_s^{-l-1}.
 \tag{IL.26}
\]
Their product is \(P/u_s\), so the map takes
\(\alpha_l\beta_l=v\) to \(A_lB_l=s\). Its inverse uses the
two reciprocal coefficients at the recovered source point.
Under \(L\) its two target factors are \(P^{-2},P^2\).
Under \(T,l\mapsto l+12\), they are
\[
 p u_s^{12}c_{12}
       =P^2\lambda_0q\eta_2^6/t^3
 \quad\hbox{and its reciprocal}.
 \tag{IL.27}
\]
Thus the whole selected height, source phase, Miller coefficient
and line period are retained. At \(p=1\) the exact central
identity is \(\eta_2^6=4u_t^2\), making IL.27 equal to
\(4\lambda_0\); the equality before this restriction remains
the actual boundary arrow. For an arbitrary source-cell representative
write the original chain label as \(12j+l\). The actual inverse period is
\[
 X_*=p^jX,\quad
 Y_*=P^{-12j}v^{-12j}c_{12}^{-j}Y,\quad
 V_{{\rm pt},*}=(P^{12}v^{12}pc_{12})^{-j}V_{\rm pt}.
 \tag{IL.27a}
\]
It takes the positive parent coefficient to
\(\alpha_*=(P^{12}c_{12})^{-j}\alpha_{12j+l}\).
Thus the two original reference-chain coefficients acquire,
in addition to IL.26 evaluated at \(X_*\), the respective
factors \(\gamma^j,\gamma^{-j}\), where
\[
                   \gamma=pu_s^{12}c_{12}
                         =P^2\lambda_0q\eta_2^6/t^3.
 \tag{IL.27b}
\]
This follows by substituting
\(P^{12}v^{12}pc_{12}=s^{12}\gamma\) in
IL.27a, and retains every chronological source and line power.

The final target of IL.26 is the actual fixed-source reference
at \(s_*\ne0\). We spell out its source transport. In the
first graph chart of IL.20,
\(s=(ay-P)/u_s(P,ay)\) has \(\partial_Ps\) a unit of
central value \(-1\). In either IL.21 chart,
\(s=(R-uw)/u_s(uw,R)\) has \(\partial_Rs\) a unit of
central value \(+1\). Take respectively
\((\partial_Ps)^{-1}\partial_P\) and
\((\partial_Rs)^{-1}\partial_R\) as local lifts.
On the original \(t\)-boundary the prescribed source-period map is
\[
 Z_s=\exp\!\left(b\,\operatorname{Log}(Q_s/Q_0)\right)Z_0,
 \quad
 V_{{\rm pt},s}=\exp\!\left(
 b\,\operatorname{Log}(P_s^2/P_0^2)\right)V_{{\rm pt},0}.
 \tag{IL.28}
\]
The marked real coordinate \(b\) increases by one under the
source deck. The two target factors are then exactly \(Q_s\)
and \(P_s^2\). The full ratios have their source-path logarithms
fixed to zero at \(s=0\), and the complete IL.24--25 module
arrows are composed with IL.28.

A relative partition with sum one joins these local source
lifts, agreeing with the derivative of IL.28 on a boundary
collar. The real field along a source path \(s(a)\) is
\(\dot s X+\overline{\dot s X}\). It is tangent to the
retained \(t\)-boundary and exists for the whole interval
because the source graph over the closed domain is proper.
Its backward flow is the inverse. Write the resulting source
map as \(T_s:S_0\to S_s\), and \(K_s=T_{s_*}T_s^{-1}\).

For original vector frames \(e_j=g_{ij}e_i\), set
\[
 a_i=-\sum_k\chi_k\,Xg_{ik}/g_{ik},\qquad
 a_j=a_i+Xg_{ij}/g_{ij}.
 \tag{IL.29}
\]
The horizontal vector in the coefficient \(w_i\) is
\(X-a_iw_i\partial_{w_i}\). A relative cutoff makes it agree
with the complete boundary derivative of IL.24--28. Its
linear ODE, and its inverse backward ODE, give the actual
line transport; the reciprocal ODE gives the dual.
Consequently the full boundary map is
\[
 (z,\alpha_l,\beta_l)\longmapsto
 \bigl(K_sz,\ K_s^{\rm line}(H_{+,l}\alpha_l),
                   K_s^{\rm dual}(H_{-,l}\beta_l)\bigr).
 \tag{IL.30}
\]
The inverse first applies \(K_s^{-1}\), then the reciprocal
IL.26 factors. On every period translate, the construction is
defined by conjugating with the full IL.3 map and its inverse;
this includes its source-dependent derivative.

The actual good cubic cap retains SI2.28\textquotesingle s ordinary-to-cap
module map, its deck character and all sector integers.
For the original height \(q\) its exact four-turn square is
\[
           \Phi_{\rho+4}\,q^4=h_\infty^4\Phi_\rho,
                 \qquad h_\infty=1-t_c(s)z^3.
 \tag{IL.31}
\]
The inverse closing uses \(q^{-4},h_\infty^{-4}\).
IL.30 is composed with that point-line cap map and its
reciprocal on the prescribed annulus. Its relative source
lift extends through the actual proper cap charts with that
fixed boundary derivative. Hence the full parameter square
is preserved.

All three vanishing discs remain in these source bodies.
The two-node unselected source and the selected node are the
three actual charts IL.22. The fourteen-factor RB resolution
has their scalar-source fixed components at
\(\alpha_l=\beta_l=0\); its overlap with the moving selected
fixed curve requires the nonzero coordinate \(v_{12}\).
It consequently retains each selected vanishing disc in the
scalar-source component. These explicit overlaps give the
source boundary inclusion
\[
 H_1(\partial S)=\mathbb Z\langle b,\alpha_t\rangle
                  \oplus\mathbb Z/3\langle a\rangle
 \longrightarrow H_1(S)=\mathbb Z\langle b\rangle,\qquad
                  (a,\alpha_t,b)\mapsto(0,0,b).
 \tag{IL.32}
\]
Indeed the three plumbing maps give
\(\alpha_t b\alpha_t^{-1}=a^3b\), and each vanishing
disc kills the same primitive \(a\); the retained O-disc
kills \(\alpha_t\). This also specifies the boundary on
which IL.30 is used.

\par

\medskip\noindent\textbf{IL.5. The actual cube frame and the twelve closed normal pairings}

\par\nobreak

The reference positive normal in the scalar chain is the
genuine point line \(M_2\), and its negative normal is
its reciprocal. IL.24--25 show the central point/O degree
vector \((-1,1,0)\). To retain its full boundary class,
write
\[
 \tau_L=\frac{\log|X|}{\log|Q|},\qquad
 W_p=e^{-2\tau_L\operatorname{Log}p}X^2V_{\rm pt}^3.
 \tag{IL.33}
\]
The logarithm of \(p\) is the actual small boundary logarithm.
Under \(L\), \(\tau_L\) increases by one and
\(X^2V_{\rm pt}^3\) acquires \(p^2\), so IL.33 is exactly
invariant. Its inverse is a vector frame of \(M_2^{\otimes3}\).
At \(p=1\), the complete theta identity in the actual
point/O modules gives this frame divisor
\[
                    6F-2C_1-C_2.
 \tag{IL.34}
\]
More explicitly it is
\[
 \frac{P^6X}{(X-P)^3}
       E_{P^2,P}(P^3/X,X)^{-3}\sigma_2^{\otimes3},
 \tag{IL.35}
\]
where \(\sigma_2\) is the original canonical meromorphic
section. The arguments of the entire product are exactly
the band-one source coordinates \((x,y)=(P^3/X,X)\).
Indeed IL.23 and the literal \(Q=P^3\) reindex give
\[
 f_{P^2}(1/X)
      =P^{-2}(1-P/X)E_{P^2,P}(P^3/X,X).
 \tag{IL.35a}
\]
Its cube, multiplied by \(X^2\), has IL.35 as its inverse
coefficient in the canonical section. The point/O horizontal factors cancel through
IL.24--25. Multiplication by \(P^{-6}\) gives the literal
coefficient-zero-on-\(C_0\) frame with divisor
\(-2C_1-C_2\), retaining that entire fibre clutch.

Thus the positive reference relative lift is
\[
 D_+=-(2C_1+C_2)/3,\qquad D_+^2=-2/3.
 \tag{IL.36}
\]
Its boundary character has exact order three, since
\(3(1,0)^t=I(-2,-1)^t\) while \(I^{-1}(1,0)^t\)
has denominators three. For the actual frame IL.33 the
lift is \(2F+D_+\), with the same square. The negative
reference line uses the reciprocal frame. Therefore
\[
 (c_-^2,c_-c_+,c_+^2)_{\rm ref}=(-2/3,+2/3,-2/3).
 \tag{IL.37}
\]

This computation retains the remaining full boundary gauges.
By IL.32,
\[
 H^1(\partial S;\mathbb Z)/\operatorname{im}H^1(S;\mathbb Z)
                    =\mathbb Z\langle\alpha_t^*\rangle.
 \tag{IL.38}
\]
Its connecting class is the primitive fibre, realized by
the vector-frame clutch \(e_R=(P/\epsilon)e_L\).
The original \(t\)-meridian has the same class through
\(t=Pu_t\), with its full \(u_t\) factor retained.
The \(b\)-character extends over \(S\), and a zero-period
boundary ratio has a logarithm extended by a collar cutoff.
Every actual gauge in IL.26--31 therefore changes a relative
lift only by the retained fibre multiple. It changes no
square or mixed pairing because \(F^2=F\cdot N_\pm=0\).

Choose the norm making the inverse IL.33 cube frame unit,
and take its actual third roots on a simply connected
boundary cover. Their transition functions are constant
elements of \(\mu_3\); declaring these frames parallel
gives the actual flat order-three boundary connection.
The reciprocal connection belongs to the negative line.
Transport both through IL.30. The interpolation to any
previous prescribed boundary connection is placed on the
opposite collar with its coordinate reversed. This is the
identical full interpolation pulled back by the reversal;
its complete transgression cancels pointwise in the glued
double.

The collision negative degree vector on \(C_1,C_2\)
differs from the reference negative vector by
\((1,1)^t=I(-1,-1)^t\). Thus their boundary torsion
characters agree by the actual coefficient arrow IL.26.
The positive characters agree as well. Gluing through IL.30, using the same cube connections and retained
fibre gauges, subtracts the reference pairings IL.37 from the minimal
collision pairings above.  The resulting closed normal pairings are
\((0,-1,0)\).  The entire minimal collision source is carried to
the reference body by the displayed source flow; its derivative and
inverse match the boundary tangent connections.  The identical
opposite-collar interpolation cancels, and the tangent Pontryagin
pairing of these new minimal doubles is zero.  No source blowup is
inserted into these lower stars.

On the fixed surface \(s=A_lB_l=0\), the first normal
derivative of \(s\) is zero. A circle-averaged source
partition has no linear term of nonzero physical weight.
The normal derivative of the transport is consequently
the complex diagonal map IL.26, in the original Cartier
modules, and adds no unrecorded tangent-normal mixing.

\par

\medskip\noindent\textbf{IL.6. Physical characteristic evaluation and exact provenance}

\par\nobreak

Let \(a=c_1(N_-)\), \(b=c_1(N_+)\), with physical weights
\((-1,+1)\). Use \(u=c_1(U)\), \(U=M^*\), so the
physical class is \(c_1(M)=-u\). The actual equivariant
normal Euler factor is \((-u+a)(u+b)\). Expanding its
inverse and the normal tangent Pontryagin factors gives
the surface contributions
\[
 k_M=\int(a^2-ab+b^2),\qquad
 \ell_M=\int\bigl(p_1(TS)+(a+b)^2\bigr).
 \tag{IL.41}
\]
In detail, the degree-four coefficient of
\((-u)^3/[(-u+a)(u+b)]\) is \(a^2-ab+b^2\).
The corresponding coefficient of
\[
 \frac{(-u)\bigl(p_1(TS)+(-u+a)^2+(u+b)^2\bigr)}
                {(-u+a)(u+b)}
\]
is \(p_1(TS)+(a+b)^2\). These are evaluated with
the original complex fixed-surface orientation and the complete
source, normal and cap maps just displayed.

For the new minimal doubles the complete pairings are
\((0,-1,0)\) and the tangent pairing is zero.  The displayed
physical Euler expansion therefore gives \((1,-2)\) for each of the
twelve actual labels, as verified in supporting proof 20.

The whole scalar-reference fan comparison RFB changes
the complete reference pair from \(T_1\) to \(T_{-1}\)
by \((-72,0)\). The new lower computation uses the actual RB \(T_1\)
reference pieces and their own cap maps. The whole RFB
correction is applied once in the complete outer
assembly, after these lower pieces and the separately
computed moving/reference cap have been included.


% End expanded verified cp36 input


% Expanded verified cp36 input: notes/part3_e6_selected_i3_m2_return_support_17.tex
\par\bigskip\noindent\textbf{P3K.111 supporting proof 17: Whole reference fan bridge with exact Jin type}\par\nobreak
\par

\medskip\noindent\textbf{Selected-I3 \(M_{2}\): the complete scalar-reference fan and cap comparison}

\par\nobreak

Working proof, 18 November 2026. This note retains the fixed original
source \(s=s_*\ne0\), \(M_2=\mathcal O(2G-O)\), the chronological
translation \(-3G\), the cubic good cap of norm one with its prescribed
finite character, and every SI2 annular integer. The independent
scalar-disc coordinate is \(\sigma\), its root is \(r_0\sigma^4\),
and the suspension scalar is \(\lambda_0\sigma^{12}\).

The reference constructed in RB has fan \(T_1\). The executed
collision native star NS has fan \(T_{-1}\). We construct the two
actual reverse \(C_1\) changes on the nonzero scalar boundary circle,
including their Cox correspondences, and compare their completed
circle quotients with the complete source, line, cap and outer maps.
The resulting \textbf{whole physical characteristic-pair correction} is

\[
 (k_{M,-1}-k_{M,1},\ell_{M,-1}-\ell_{M,1})=(-72,0).
 \tag{RFB.1}
\]

The pair in RFB.1 belongs to this independent scalar word \((0,4)\).
All reference fixed components and the reference moving cap are
included in its whole pair. No individual moving-cap value is obtained
by subtracting RFB.1 from a partial fixed-component ledger. The
comparison is a construction on the actual nonzero-parameter family
and its completed quotient. The analytic scalar-disc reference RB,
including its fourteen-factor charts, is kept unchanged.

\par

\medskip\noindent\textbf{RFB.1. The named periods, fan labels and the full endpoint}

\par\nobreak

Write \(q=t-t_c(s_*)\), and retain the complete RB functions

\[
 \begin{gathered}
 Q=\kappa q,\quad \kappa=P^2/c_q,\quad \alpha=P^2,\quad g=P^{-3},\\
 \eta_2=a_T\Theta'_Q(1)\frac{\Theta_Q(P^2)}{\Theta_Q(P)^2},\qquad
 \Gamma=\lambda_0\eta_2^6t^{-3},\qquad \eta=\sigma^{12}\Gamma .
 \end{gathered}
 \tag{RFB.2}
\]

On the selected neighbourhood of this fixed source, all functions
except the displayed \(q\) and \(\sigma\) factors are nonzero
holomorphic functions of the original \(q\). Their logarithms are
chosen on one common smaller \(q\)-disc. In particular \(\eta_2\)
is the full Miller/theta unit; none of its factors is replaced by
its central value. The original periods are

\[
 h:(Z,V)\mapsto(\kappa qZ,\alpha V),\qquad
 B:(Z,V)\mapsto(gZ,\eta qV),\qquad V=V_{\rm pt}.
 \tag{RFB.3}
\]

The positive scalar-circle coordinate is \(\sigma=\rho e^{2\pi i\theta}\),
with fixed \(\rho>0\). Hence the two line windings of RFB.3 are
\((0,12)\), while \(\kappa,g\) have winding zero.

For each of the three integers \(j=1,0,-1\), use the exact characters
and period generators

\[
 Z_j=Z,\qquad V_j=Z^{-j}V,\qquad h_j=hB^j,\quad B_j=B.
 \tag{RFB.4}
\]

Their inverses are \(V=Z_j^jV_j\) and \(h=h_jB_j^{-j}\). Substitution
in both original periods gives the full, untruncated functions

\[
 \kappa_j=\kappa g^j,\qquad
 \alpha_j=\alpha\eta^j\kappa^{-j}g^{-j^2},\qquad
 \eta_j=\eta g^{-j},
 \tag{RFB.5}
\]
\[
 h_j:(Z_j,V_j)\mapsto(\kappa_jqZ_j,\alpha_jV_j),\qquad
 B_j:(Z_j,V_j)\mapsto(gZ_j,\eta_jqV_j).
 \tag{RFB.6}
\]

The line windings are \((12j,12)\). All negative powers in RFB.5
are retained on the nonzero scalar circle. The actual affine
extensions are the primitive character charts on the indicated
fan, with full six-chart rows

\[
\begin{array}{c|ccc}
 &u&v&w\\\hline
 C_0&\kappa_j\alpha_jbc&\alpha_j^{-1}ab&\kappa_j^{-1}b^{-1}\\
 C_1&g\eta_jbc&\eta_j^{-1}c^{-1}&g^{-1}ac\\
 C_2&a^{-1}&ab&ac .
\end{array}
\qquad Z_j=b,\quad V_j=c,\quad q=abc=uvw.
 \tag{RFB.7}
\]

For each row recover the original \(q\) from the target triple
product and evaluate every inverse coefficient at that recovered
value. These maps and inverses retain the full normal jets.

Put \(u_j=(1,j)\), \(v_j=(-1,1-j)\), so \(u_j+v_j=\delta=(0,1)\).
The two positive-diagonal triangles are
\(\{0,u_j,\delta\}\), \(\{0,v_j,\delta\}\).
Translating the second by the actual \(h_j\) sends it to
\(\{u_j,\delta,u_j+\delta\}\). These are precisely the two
triangles of \(T_j=T(e_1+je_2,e_2)\); the translation carries the
complete RFB.6 line coefficient. Thus the RB matrix
\(\left(\begin{smallmatrix}1&-1\\1&0\end{smallmatrix}\right)\)
is \(T_1\), and the NS matrix
\(\left(\begin{smallmatrix}1&-1\\-1&2\end{smallmatrix}\right)\)
is \(T_{-1}\), with their exact point marking.

The actual ordinary forward endpoint is \(q^4\). After its specified
\(B^{-4}\) reindex the bounded endpoint is

\[
 E:(Z,V)\mapsto(AZ,DV),\qquad A=g^{-4},\quad D=\eta^{-4}.
 \tag{RFB.8}
\]

In the \(j\)-th native \(C_0\) chart set

\[
 a_j=q/(Z_jV_j),\quad c_j=V_j,\quad q=Z_ja_jc_j,\qquad
 D_j=D A^{-j}=\eta^{-4}g^{4j}.
 \tag{RFB.9}
\]

The full endpoint and its reciprocal inverse are

\[
 E_j:(Z_j,a_j,c_j)\mapsto
       (AZ_j,(AD_j)^{-1}a_j,D_jc_j).
 \tag{RFB.10}
\]

At \(j=1\), this is exactly MN.12. At \(j=-1\) the native
characters are exactly NS.16. The two normal factors in RFB.10
are, respectively,

\[
\begin{array}{c|cc}
j&a_j&c_j\\\hline
1&\eta^4&g^4\eta^{-4}\\
0&g^4\eta^4&\eta^{-4}\\
-1&g^8\eta^4&g^{-4}\eta^{-4}.
\end{array}
 \tag{RFB.11}
\]

In particular the full \(\sigma^{-48}\) factor remains in the
positive normal. This table is a direct substitution in the named
forward endpoint; it does not use the source-word return.

\par

\medskip\noindent\textbf{RFB.2. The two extending Cox changes and the unchanged exterior}

\par\nobreak

In the \(C_1\) row of RFB.7 the two native triples are

\[
 (z,n_0,n_1)=(c,a,b),\qquad
 (z',n'_0,n'_1)=(\eta_jv,gw,(g\eta_j)^{-1}u).
 \tag{RFB.12}
\]

They satisfy \(z'=z^{-1},n'_i=zn_i\) literally. Take Cox
coordinates \((a_0,a_1,c_0,c_1)\) with auxiliary weights
\((1,1,-1,-1)\), \(z=a_1/a_0\), \(n_i=a_0c_i\), and
\(q=a_0a_1c_0c_1\). On the opposite ruling\textquotesingle s \(c_1\ne0\)
chart the regular coordinates are

\[
 (a',b',c')=(c_0/c_1,a_0c_1,a_1c_1).
 \tag{RFB.13}
\]

On the common torus they are \((a/b,b,bc)\); their torus inverse
is \((a,b,c)=(a'b',b',c'/b')\). Thus the actual target characters
are \(Z_{j-1}=Z_j,V_{j-1}=Z_jV_j\), not a scalar-only substitution.
The opposite \(c_0\ne0\) chart and its reciprocal overlap give the
other target vertex. Its complete generators are \(h_{j-1}=h_jB^{-1}\),
and the full coefficients obey

\[
 \kappa_{j-1}=\kappa_jg^{-1},\quad
 \alpha_{j-1}=\alpha_j\eta_j^{-1}\kappa_jg^{-1},\quad
 \eta_{j-1}=\eta_jg .
 \tag{RFB.14}
\]

Here is also the common proper graph, with both actual morphisms.
On the total space of \(\mathcal O(-1,-1)\) use homogeneous vectors
\(A=(A_0,A_1),C=(C_0,C_1)\) and a coefficient \(\tau\), with

\[
 (A,C,\tau)\sim(lA,mC,(lm)^{-1}\tau).
\]

Its two projections are

\[
 p_+(A,C,\tau)=[A_0,A_1,\tau C_0,\tau C_1],\qquad
 p_-(A,C,\tau)=[\tau A_0,\tau A_1,C_0,C_1].
 \tag{RFB.15}
\]

Both pull the original height back to
\(\tau^2A_0A_1C_0C_1\). On the complement of the respective
zero section, the inverse to \(p_+\) takes
\([C]=[c_0:c_1]\) and the unique line coefficient \(\tau\)
with \(c=\tau C\); the inverse to \(p_-\) uses
\([A]=[a_0:a_1]\) and \(a=\tau A\). Their expressions in any
nonzero homogeneous coordinate are holomorphic reciprocal maps.
At \(\tau=0\) the displayed projections retain the two separate
projective rulings. No inverse across that exceptional divisor is
claimed.

Every bounded torus arrow \((Z_j,V_j)\mapsto(AZ_j,DV_j)\)
has the full \(C_1\) Cox lift

\[
 (a_0,a_1,c_0,c_1)\mapsto
 (a_0,Da_1,(AD)^{-1}c_0,Ac_1).
 \tag{RFB.16}
\]

It preserves \(q\), so its inverse divides by the same full
functions at that same \(q\). It commutes with the auxiliary and
physical circles. In the target chart RFB.13 its factors are
\(((A^2D)^{-1},A,AD)\), exactly the \(D\mapsto AD\) reset
of RFB.14. The same diagonal map on \(A,C\), with \(\tau\)
unchanged, intertwines both projections RFB.15. Thus every bounded
seam and every unit arrow has an actual lift on the common graph.
The period translates are defined by RFB.7 and RFB.14 on their
labelled affine charts; the written cancellations of their \(q\)
powers take place before the unit Cox maps are used.

Transport the actual Hermitian norms by RFB.16. The two common-link
Hopf inverses are \(v=u^{-1}q_H\), \(u=q_Hv^{-1}\); their auxiliary
circle ambiguity is exactly the written Cox action. The exterior
arrow and its inverse are

\[
 T(y,q_H)=(q_H^{-1}yq_H,q_H),\qquad
 T^{-1}(y,q_H)=(q_Hyq_H^{-1},q_H).
 \tag{RFB.17}
\]

These are the full IF3.17--19 common-link maps, with all four Cox
weights and paired half-phase endpoints retained. The explicit
normal-rescaling flow of IF3.2, using its \(q\)-height inverse,
carries the full nonlinear germ on a uniform smaller tube. Apply
RFB.12--17 first for \(j=1\) and then \(j=0\). On the original
outside source and finite-cap charts use the identity, with the
inverse character arrow RFB.4 on their torus overlap. This gives
the actual two fan families and their same original endpoint.

The analytic reference\textquotesingle s first reverse change can also be seen in
RB\textquotesingle s factor graph: interchange the final \(x_i,\bar y_i\) in
\((\sigma^{12},x_i,\bar y_i)\). On the nonzero scalar circle this
is exactly RFB.15. The subsequent change is the next periodic
\(C_1\) cell, with the complete \(B\) label and RFB.14; its scalar
factors are retained. The whole-pair calculation below is on this
completed nonzero-circle family and does not require extending a
dense \(Z^2\) multiplier through a scalar-zero vertex.

\par

\medskip\noindent\textbf{RFB.3. The exact cut, native sphere and finite-cap squares}

\par\nobreak

Choose an endpoint-flat cut function \(\chi\), and actual cut
logarithms of \(g\) and \(\eta\). Set

\[
 \mathcal A_\theta=\exp(4\chi\operatorname{Log}g),\quad
 \mathcal D_\theta=\exp(4\chi\operatorname{Log}\eta),\quad
 \mathcal D_{j,\theta}=\mathcal D_\theta\mathcal A_\theta^{-j}.
 \tag{RFB.18}
\]

The two native chart maps are

\[
 C^0_{j,\theta}(Z,a,c)=
 (\mathcal AZ,(\mathcal A\mathcal D_j)^{-1}a,\mathcal D_jc),
\]
\[
 C^\infty_{j,\theta}(z,a_\infty,c_\infty)=
 (\mathcal A^{-1}z,\mathcal D_j^{-1}a_\infty,
                         \mathcal A\mathcal D_jc_\infty).
 \tag{RFB.19}
\]

Both preserve the original \(q\); their inverse uses its unchanged
value before evaluating the reciprocal coefficients. They obey
\(C_{j,1}E_j=C_{j,0}\) on complete endpoint collars. Extend to the
next cut by \(C_{j,\theta+1}E_j=C_{j,\theta}\), retaining every
derivative. The angular part of \(\operatorname{Log}\eta\) remains
\(24\pi i\theta\), including its \(48\)-fold endpoint normal phase.

Write \(S_j(Z,a,c)=(Z,Z^ja,Z^{-j}c)\) on its stated character
localization, with reciprocal \(S_j^{-1}(Z,a,c)=(Z,Z^{-j}a,Z^jc)\).
There is the exact full cut square

\[
 C_{j,\theta}S_j=S_jC_{0,\theta},\qquad
 S_j^{-1}C_{j,\theta}=C_{0,\theta}S_j^{-1}.
 \tag{RFB.19a}
\]

For example the left side\textquotesingle s negative normal coefficient is
\(\mathcal A^{j-1}\mathcal D^{-1}Z^ja\), and its positive
coefficient is \(\mathcal D\mathcal A^{-j}Z^{-j}c\);
the right side has exactly these two coefficients. The identity
\(\mathcal D_j=\mathcal D\mathcal A^{-j}\) retains the complete
\(\eta\) phase. RFB.15--16 and the directed Cartier modules below
are its extending chart arrows. Its first-normal derivative and
the two cap charts therefore carry this same square, with its
reciprocal, through RFB.21--23. This is the shared inner gauge used
in RFB.42; no endpoint scalar phase is suppressed.

The vertical velocity is
\((\partial_\theta\log\mathcal A)Z\partial_Z
 -\partial_\theta\log(\mathcal A\mathcal D_j)a\partial_a
 +(\partial_\theta\log\mathcal D_j)c\partial_c\).
It annihilates \(q\). Add the prescribed horizontal parameter lift,
and cut off only this vertical difference on a larger selected tube.
Its flow and backward flow give the full relative collar map;
the actual cubic cap and its annular reference are fixed there.
In particular their square remains

\[
                 \Phi_{\rho+4}q^4=h_\infty^4\Phi_\rho,
 \qquad h_\infty=1-t_c(s_*)z^3,
 \tag{RFB.20}
\]

with the reciprocal square in the opposite direction. The actual
root action \(A_0^3=B_0\), the three chronological line factors,
the finite character and every sector integer are exactly RB.23--24.

For clarity, the selected Hopf cap is also fixed by a literal map.
Use NC\textquotesingle s matrices

\[
 T_Z=\begin{pmatrix}1&-\bar Z\\Z&1\end{pmatrix},\quad
 T'_z=\begin{pmatrix}z&-1\\1&\bar z\end{pmatrix},\quad
 T'_{1/Z}\operatorname{diag}(Z,\bar Z)=T_Z.
\]

The complete first-normal derivative of RFB.19 is

\[
 G_{j,\theta}(Z)=T_{\mathcal AZ}
 \operatorname{diag}((\mathcal A\mathcal D_j)^{-1},
                                  \overline{\mathcal D_j})T_Z^{-1}.
 \tag{RFB.21}
\]

Here the functions are their actual \(q=0\) restrictions; their
first normal derivatives vanish since \(dq\) vanishes on the
fixed sphere. The other chart equality above proves smoothness
at infinity, with the full inverse matrix. On the trace-one
positive Hermitian three-ball use

\[
 \Psi_{\mathcal A}(Y)=
 \frac{\operatorname{diag}(\mathcal A,1)Y
       \operatorname{diag}(\mathcal A,1)^*}
      {\operatorname{tr}(\operatorname{diag}(\mathcal A,1)Y
       \operatorname{diag}(\mathcal A,1)^*)}.
 \tag{RFB.22}
\]

Its inverse is \(\Psi_{\mathcal A^{-1}}\), and on the boundary
it is exactly \(Z\mapsto\mathcal AZ\). In RFB.21 replace the two
cut logarithms by their product with a radial function zero near
the ball centre and one near its boundary. The resulting
\(\widehat G_{j,\theta}\) is independent of the sphere direction
near the centre and has the exact other-chart expression. Hence

\[
 \widehat C_{j,\theta}(Y,q_H)=
 \left(\Psi_{\mathcal A}Y,
 \frac{\widehat G_{j,\theta}(Y)q_H}
      {\|\widehat G_{j,\theta}(Y)q_H\|}\right)
 \tag{RFB.23}
\]

is the actual invertible cap map. Its inverse first applies
\(\Psi_{\mathcal A^{-1}}\), then the inverse matrix at that recovered
point and the positive normalization. The cap forward seam is
\(\widehat E_j=\widehat C_{j,1}^{-1}\widehat C_{j,0}\).
The resolved boundary map is \(B_{j,\theta}=C_{j,\theta}^{-1}N_{j,\theta}\);
the exact square is \(\widehat C_{j,\theta}B_{j,\theta}=N_{j,\theta}\),
with reciprocal inverse. At \(j=1\) these are precisely MN.12--14
and NC.10--20. The \(j=0,-1\) caps are specified by this same
construction with their actual RFB.11 coefficients.

\par

\medskip\noindent\textbf{RFB.4. An actual completed orbit-piece and relative-line morphism}

\par\nobreak

We give the full named application of K72 SI.6--27, rather than
identifying the pieces from their winding matrix. For \(j\in\mathbb Z\)
use \(Q_j=\kappa qg^j\) and the actual contracting coefficient

\[
 \gamma_j=
 \frac{\log|Q_j|\log|V_j|-\log|\alpha_j|\log|Z_j|}
 {\log|Q_j|\log|q\eta_j|-\log|\alpha_j|\log|g|}.
 \tag{RFB.24}
\]

The denominator has leading term \((\log|q|)^2\) and is nonzero
on a uniform small disc. The displayed full positive norms are
independent of \(\theta\). Direct substitution gives
\(\gamma_j(h_jp)=\gamma_j(p)\) and
\(\gamma_j(Bp)=\gamma_j(p)+1\). The F9 quotient coordinate has

\[
 2b=r-|y|^2/r,\quad2b'=|x|^2r-r^{-1},\quad
 db'/db=(|x|^2+r^{-2})/(1+|y|^2r^{-2})>0.
 \tag{RFB.25}
\]

It is matched to the contracting coordinate by the proper decreasing
map \(1/2-\epsilon(r-1)/(r+1)\), with \(0<\epsilon<1/4\).
Both branch derivatives are \(-4\epsilon\). The second fixed-sphere
endpoint is the actual \(h_j^{-1}\) translate of the second
conifold endpoint and therefore has the same lifted \(\gamma_j\).
The endpoint heights stay distinct in that short interval; no
integer is inserted. Convex interpolation of the proper decreasing
maps, fixed on the outside collar, gives the full matching and its
inverse. Untwisting by \(\exp(\gamma\log g)\) and the complete
positive line norm gives the actual orbit piece
\(W_j=E_j\times S^1_\gamma\), with the marked outer coefficients
unchanged. The line phases remain in its principal circle bundle.

The explicit source comparison to \(j=0\) uses

\[
 Q_{a}=\kappa q\exp(aj\log g),\qquad
 X=\partial_a+j\log g(q)\,y\partial_y
          \quad\hbox{on }xy=Q_a,\quad0\le a\le1.
 \tag{RFB.26}
\]

It fixes the original \(q\). The inverse local base map recovers
\(q\) by the inverse of \(q\mapsto Q_a(q)\); its derivative at zero
is the nonzero \(\kappa(0)g(0)^{aj}\), so one smaller disc works
uniformly. On the outside collar prescribe instead the real marked
period map \(u+\tau_av\mapsto u+\tau_0v\). Both vector fields
project to \(\partial_a\); their difference is vertical. A finite
partition, fixed on the protected collars, glues them. The compact
boundary-tangent field has a complete flow, with the backward flow
as inverse. It is independent of the scalar phase.

The fixed normalization spheres have their lifted graphs in the
same short \(\gamma\) interval. Interpolate them with the distinct
endpoint heights fixed. The source normalization is injective away
from its two node preimages, so this interpolation is an embedding.
Extend its velocity over a tube and cut it off. This gives the actual
relative ambient isotopy, not a replacement of the nodal source by
real logarithmic coordinates.

The native normal arrow is

\[
 (Z,a_0,c_0)\mapsto(Z,Z^j a_0,Z^{-j}c_0),\qquad q=Za_0c_0.
 \tag{RFB.27}
\]

It is interpreted on the two \textbf{specified} fan stars. On their
normalization sphere its full positive-spinor map is the Cartier
module arrow

\[
 E_j^{\rm spin}=E_0^{\rm spin}\otimes
       \mathcal O(j(D_0-D_\infty))\longrightarrow E_0^{\rm spin},
                 \qquad w\longmapsto Z^j w.
 \tag{RFB.28}
\]

The defining Cartier generators cancel the written zeros and poles
at \(D_0,D_\infty\). RFB.27 is the character arrow from the \(0\)
star to the \(j\) star; its positive spinor coefficient matrix is
\(\operatorname{diag}(\overline Z^{\,j},Z^{-j})\). RFB.28 is its
directed inverse from the \(j\) spinor bundle to the \(0\) bundle,
with matrix
\(\operatorname{diag}(\overline Z^{-j},Z^j)\). After the transported
positive norms this inverse has the common \(Z^j\) phase. The
remaining positive factors are the explicit diagonal norm
rescalings, with their reciprocal inverse. This proves the
directed spinor map, including the physical circle. Its adjoint
carries the normal three-plane. The
source derivative and the graph velocity give their actual positive
triangular maps. Extend the sphere velocity over the ball by a
collar cutoff and use RFB.23 on its Hopf normal sphere. These maps
give the completed orbit-piece diffeomorphism \(I_j:C_j\to C_0\),
where

\[
 C_j=(W_j\setminus\operatorname{int}(F_j\times D^3))
                      \cup(D^3\times S^2).
 \tag{RFB.29}
\]

All inverses are the displayed reciprocal bundle maps and backward
flows. In the original cut coordinates the cap map is
\(\widehat C_{0,\theta}^{-1}I_j\widehat C_{j,\theta}\).
Its endpoint identity follows directly from
\(\widehat C_{j,1}\widehat E_j=\widehat C_{j,0}\) and the
periodic closed-coordinate \(I_j\). On the principal circle line
the additional tensor is retained next; it is trivialized on this
cap by making the graph map \(b\) equal to zero there.

Use the actual coefficient decomposition

\[
 (\log Z,\log V)/(2\pi i)=
 u(1,0)+v(\tau,\mu)+d(0,1)+\gamma(G,\zeta),
 \tag{RFB.30}
\]

where \(\tau=\log(\kappa q)/(2\pi i)\),
\(\mu=\log\alpha/(2\pi i)\), \(G=\log g/(2\pi i)\),
and \(\zeta=\log(q\eta)/(2\pi i)\). These are the full
functions in RFB.2. The exact inverse period matrix determines the
real coefficients and retains their complete derivatives. The
integer deck maps are

\[
 q\hbox{-turn}:(u,d)\mapsto(u-v,d-\gamma),\qquad
 \theta\hbox{-turn}:(u,d)\mapsto(u,d-12\gamma).
 \tag{RFB.31}
\]

RFB.4 gives exactly

\[
 u_j=u,\quad v_j=v,\quad\gamma_j=\gamma-jv,\quad d_j=d-ju.
 \tag{RFB.32}
\]

Write \(\alpha_q=\arg q/(2\pi)\). For the meridian-one line
\(U=M^*\), the actual coefficient connections are

\[
 A_0=dd+(\alpha_q+12\theta)d\gamma,\qquad
 A_j=dd_j+(\alpha_q+12\theta)d\gamma_j+12j\theta\,dv.
 \tag{RFB.33}
\]

They are invariant under both full deck maps, and their prescribed
outer difference is exactly

\[
                  \beta_j=-j(du+\alpha_qdv).
 \tag{RFB.34}
\]

Let \(b:C_0\to\mathbb R/\mathbb Z\) extend the source graph
coordinate \(v\). It extends across the node\textquotesingle s vanishing disc
because its attaching circle is the \(u\)-circle. Its restriction
to \(F_0\) has a real lift; a supported real homotopy makes it zero
on a tube containing the surgery and hence on the cap. Its exact
theta-cut gauge is \(\exp(2\pi i12j\theta h)\), where \(h\)
is this real change of graph map. Define the actual line
\(\mathcal T\) by the \textbf{endpoint-to-initial} coefficient clutch
\(w\mapsto e^{2\pi ib}w\). On its theta cut the coefficient
connection is \(A=dd+\theta\,db\); at the endpoint this is the
pullback of \(dd\) by that exact clutch. The reciprocal
initial-to-endpoint arrow multiplies by \(e^{-2\pi ib}\).
Then \(c_1(\mathcal T)=\Xi=d\theta\wedge db\), \(\Xi^2=0\).

The Thom and Mayer-\/-Vietoris maps of this displayed surgery give
\(H_1(C_0)=\mathbb Z\langle b,\gamma\rangle\) and
\(H_2(C_0)=\mathbb Z\langle\text{Hopf meridian},b\times\gamma\rangle\).
The primitive meridian has \(U\)-value one; the real-period torus
has value zero. The theta clutch fixes these generators and adds
the \(b\times\theta\) and \(\gamma\times\theta\) evaluations.
RFB.32 gives difference \(12j\) on the first and zero on the other
three. Thus the actual line morphism is

\[
          J_j:U_j\longrightarrow I_j^*U_0\otimes\mathcal T^{12j}.
 \tag{RFB.35}
\]

Its prescribed boundary arrow is also constructed, not inferred
from those evaluations. The boundary first cohomology has basis
\((\alpha_q,\theta,v,\gamma)\); the extending image is
\((\theta,b,\gamma)\). The only remaining obstruction is therefore
the \(q\)-angle. On one parameter arc choose a generic \(Z_*\)
with \(|Z_*|=1\), \(V=1\), avoiding the finite marked sections.
The literal free section disc is

\[
 (q/Z_*,Z_*,1)\mapsto(qZ_*^{j-1},Z_*,Z_*^{-j}).
 \tag{RFB.36}
\]

It misses the surgery and has \(v=\gamma=0\); both line factors
and their reciprocals extend over \(q=0\). The source flow can fix
this disc on that parameter arc. Thus the sole boundary obstruction
is zero, as an integer. On a good cover take logarithms of the
actual ratio-line transitions, including
\(j\log\Gamma-j\log\kappa-j^2\log g\) and \(-j\log g\),
and subtract the integral one-cochain supplied by the four preceding
evaluations. In frame convention \(e_j=e^{2\pi ir_{ij}}e_i\),
set \(a_i=\sum_k\rho_kr_{ki}\); then \(a_j-a_i=r_{ij}\).
The reciprocal exponentials give RFB.35 and its inverse. Remove
the remaining extending boundary windings by the actual functions
\(\theta,b,\gamma\), then extend its real logarithm with a collar
cutoff. This makes the boundary arrow literally RFB.32--34. The
cap gauge extends because \(b=0\) there. Every original annular
integer, finite character, and full nonzero point/O module is
carried throughout this construction.

\par

\medskip\noindent\textbf{RFB.5. The whole Pontryagin and Euler integrals}

\par\nobreak

The full outer stable coframe is
\((du+(\alpha_q+\phi_\kappa)dv,dv,d\gamma,d\alpha_q,d\theta,dr_q)\).
The actual source map just constructed has matrix

\[
 G_j=I+j\phi_gE_{12}-jE_{32},\qquad
 G_j^{-1}dG_j=j\,d\phi_gE_{12},\qquad
 \phi_g=\operatorname{Im}\log g/(2\pi).
 \tag{RFB.37}
\]

Here \(\phi_g\) is independent of scalar phase. The Maurer form
squares to zero. Interpolating the two actual boundary connections
by a cutoff multiple of this form gives \(\operatorname{tr}F^2=0\)
on the whole collar. No line-phase term has been added to a tangent
frame; its phase remains in RFB.33.

Choose nested tubes containing all surgery and connection changes.
Away from them take the tangent connection pulled from the actual
nodal source, with the trivial \(\gamma\) direction and the marked
flat outside collar. A regular level of the original source graph
is \(N^3\times S^1_\gamma\), on which that source \(p_1\) form
vanishes. The supported homotopy making \(b\) constant near the
larger tube gives a relative regular-level cobordism; choose a
regular value outside that constant. Its new level misses the
surgery. Stokes and the zero full outer transgression preserve the
pairing with the closed \(p_1\) form. This proves

\[
                    \langle p_1,\Xi\rangle=0.
 \tag{RFB.38}
\]

The exact boundary primitive RFB.34 contributes zero to this
pairing because RFB.37 vanishes throughout its collar. Therefore
\(\ell_{M,j}=\ell_{M,0}\) for these actual completed pieces,
with the original exterior unchanged.

For the Euler cube choose the affine interior curvature supplied
by RFB.35: \(F_j=F_0+12j\Xi\). Since \(\Xi^2=0\), its entire
interior difference is linear in \(j\), with one fixed coefficient.
Use the piece coefficient \(\gamma=\gamma_j\), so the retained
exterior coefficients are \(\gamma_{\rm ext}=\gamma+jv\) and
\(d_{\rm ext}=d_j+ju\). Its exact curvature is

\[
 F_{{\rm ext},j}=(d\alpha_q+12d\theta)\wedge(d\gamma+j\,dv),
 \qquad d\beta_j=-j\,d\alpha_q\wedge dv.
 \tag{RFB.39}
\]

Orient the outward boundary by
\(\Omega=d\alpha_q\wedge du\wedge dv\wedge d\gamma\wedge d\theta\).
The marked period cell has integral one and the exact deck maps
preserve its volume. Direct exterior algebra gives

\[
 3\int_0^1\beta_j\wedge(F_{{\rm ext},j}+a\,d\beta_j)^2\,da
                         =36j^2\Omega.
 \tag{RFB.40}
\]

The actual outward connection runs from piece to exterior, so its
cutoff runs from one to zero. Its contribution is \(-36j^2\).
The physical line is \(M=U^*\), and the K68 original-Jin parameter
arrow gives \(k_M=\int F^3\) in this positive coefficient
orientation. Both signs are the same actual Hopf/physical arrow
used in NC and MN, not a choice of positive-hyperplane generator.
Consequently, for one fixed interior coefficient \(L\),

\[
       k_{M,j}=k_{M,0}+Lj-36j^2,
       \qquad\ell_{M,j}=\ell_{M,0}.
 \tag{RFB.41}
\]

The coefficient is fixed by an actual integer comparison. In the
\(j=0\) family\textquotesingle s \(C_2\) Cox bordism the first outer fixed-branch
map is the identity, and the second is

\[
        \text{base }\kappa^{-1},\qquad
        \text{normals }\operatorname{diag}(\alpha^{-1},\kappa\alpha).
 \tag{RFB.42}
\]

These are exactly the full RFB.2 functions and are independent of
the scalar phase. The two inner charged boundaries have the same
physical-circle gauge from RFB.18--23; it extends through these
constant outer maps. Thus the line and normal-three-plane
transitions on the additional sphere boundary pull back from the
three-dimensional fixed factor. Their \(b^2,p_1(\nu)\), and tangent
\(p_1\) vanish there. The actual sphere-bundle formulas
\(4\int x^3=3b^2+p_1(\nu)\) and
\(\int p_1x=p_1(T)+p_1(\nu)\) give zero separately. This is the
full K62 SCN.9 integer map with the present \(e=4\) seam; the finite
cap, scalar phase and integral cochain have been fixed by RFB.20.
The target fan is \(T_1\): its triangle identification is the
literal \(h^{-1},B\) translation with all RFB.3 units. Its native
cap is exactly MN.12 by RFB.10 and RFB.23. Hence
\(k_{M,1}=k_{M,0}\), so \(L=36\).

Substitution gives the complete named relative formula

\[
 (k_{M,j}-k_{M,0},\ell_{M,j}-\ell_{M,0})
                         =(-36j(j-1),0).
 \tag{RFB.43}
\]

In the two requested directions the differences are

\[
 T_1\longrightarrow T_0:\ (0,0),\qquad
 T_0\longrightarrow T_{-1}:\ (-72,0).
 \tag{RFB.44}
\]

This is RFB.1. The induced whole scalar-reference index change is
\((4(-72)-0)/24=-12\). Its integer numerator was obtained from
the actual interior line, outer primitive and both cap maps. No
individual central class or free coordinate was divided to obtain it.

\par

\medskip\noindent\textbf{RFB.6. Exact scope and reproducibility}

\par\nobreak

The completed scalar boundary family with fan \(T_{-1}\) is the
genuine same-fan reference for the collision NS boundary. RFB.1
converts the \textbf{whole} RB \(T_1\) characteristic pair to that
reference. The boundary source/point/root square is still SI2,
with original root \(r_0\sigma^4\), finite character and annular
integers. The supplied correspondence, quotient diffeomorphism,
relative line morphism and cap maps have explicit inverses on
their stated domains. The reference\textquotesingle s original analytic filling
and its individual fixed-component data remain RB and MN.

The companion \path{code/part3/e6_selected_i3_m2_reference_fan_bridge.py}
checks the full character and generator reset, both Cox charts,
common-graph projections, three endpoint/cap overlaps, exact
real-period inverse and integer deck square, nilpotent tangent
collar, whole exterior-algebra cube, and the two signed pair
differences. Its check count and final source hash are recorded in
the generated JSON. These identities supplement the actual map
construction above; they do not replace its compact flow,
Cartier, cap or relative line arguments.


% End expanded verified cp36 input


% Expanded verified cp36 input: notes/part3_e6_selected_i3_m2_return_support_18.tex
\par\bigskip\noindent\textbf{P3K.111 supporting proof 18: Independent actual doubled scalar native arrows}\par\nobreak
\par

\medskip\noindent\textbf{Selected-I3 \(M_{2}\): the actual scalar-reference native products and signed flops}

\par\nobreak

Working proof, 18 November 2026. Fix the original source \(s=s_*\ne0\),
its point line \(M_2\), chronological translation \(-3G\), cubic
good cap with norm one and prescribed finite character. The scalar
root is \(r_0\sigma^4\), and \(\lambda=\lambda_0\sigma^{12}\).
All functions, inverse source charts, point/O modules, root actions,
sector integers and forward/inverse closing arrows are exactly
RFB.1--3. This is the independent scalar word \((0,4)\).

For the actual \(T_j\) family, the two single native free pairs are

\[
 (m_{C_1},n_{C_1})=(12(1-j),12(j-1)),\qquad
 (m_{C_2},n_{C_2})=(12j,-12j).
 \tag{RFS.1}
\]

These coordinates belong to K23\textquotesingle s actual relative exterior group.
The signed Cox maps give the central changes

\[
 R_{j-1}-R_j=12(j-1),\qquad R_{j+1}-R_j=-12j
                       \quad\hbox{in }\Theta_7.
 \tag{RFS.2}
\]

The first and second reverse changes \(T_1\to T_0\to T_{-1}\)
therefore contribute \(0,-12\). The whole physical integer-pair
comparison \((-72,0)\) is proved independently in RFB.4--5; no
integer Chern number is inferred from RFS.2.

\par

\medskip\noindent\textbf{RFS.1. The full single products and their literal doubles}

\par\nobreak

Keep the complete functions

\[
 \kappa=P^2/c_q,\quad\alpha=P^2,\quad g=P^{-3},\quad
 \eta=\sigma^{12}\lambda_0\eta_2^6t^{-3},\qquad
 \eta_2=a_T\Theta'_Q(1)\Theta_Q(P^2)/\Theta_Q(P)^2,
\]
\[
 \kappa_j=\kappa g^j,\quad
 \alpha_j=\alpha\eta^j\kappa^{-j}g^{-j^2},\quad
 \eta_j=\eta g^{-j}.
 \tag{RFS.3}
\]

The original selected height is \(q=t-t_c(s_*)\). The two vertex
triples satisfy \(q=abc=uvw\), with the three complete rows RFB.7.
The \(C_1\) native charts are
\((c,a,b)\) and \((\eta_jv,gw,(g\eta_j)^{-1}u)\);
the \(C_2\) native charts are \((a,b,c)\) and \((u,v,w)\).
Direct substitution gives \(z'=z^{-1},n'_i=zn_i\) in each case.

The forward endpoint is the actual \(q^4\) arrow reindexed by
\(B^{-4}\), so its bounded point map is
\((Z_j,V_j)\mapsto(g^{-4}Z_j,\eta_j^{-4}V_j)\).
Use exactly the RFB.18--23 cut maps in these two native charts.
They preserve \(q\), have reciprocal inverses evaluated at that
same \(q\), and close on full collars with all jets. The original
cubic cap retains \(\Phi_{\rho+4}q^4=h_\infty^4\Phi_\rho\), its
three chronological source/line factors, character and sector integers.

This specifies a single native product before taking a cover. Its
double is its literal pullback. On the second lap use
\(C_\theta E^{-1}\), with the original \(E\), and at the final cut
\(E^2\). Thus the capped endpoint is the actual square \(f^2\).
The root/ordinary square is pulled back to
\(\Phi_{\rho+8}q^8=h_\infty^8\Phi_\rho\), with the same two
chronological laps of the original cochain. No independent double
seam is selected.

Write the doubled scalar circle as \(\sigma=\rho e^{2i\vartheta}\),
\(0\le\vartheta\le2\pi\). The winding matrix of the full periods
in their displayed \(h_j,B\) order is

\[
 \mathsf B_j^{(2)}=2\begin{pmatrix}0&0\\12j&12\end{pmatrix}.
 \tag{RFS.4}
\]

All four remaining factors are full \(q\)-dependent units with
zero parameter winding. In particular the native closing is still
the one just specified; RFS.4 does not define a new closing.

\par

\medskip\noindent\textbf{RFS.2. The full ordinary product and its inverse}

\par\nobreak

At the second vertex use the full arrow

\[
          (u,v,w)\longmapsto(u,\eta_j(q,\vartheta)v,
                                      \kappa_j(q)w).
 \tag{RFS.5}
\]

Its target height is \(H_\vartheta(q)=q\eta_j(q,\vartheta)\kappa_j(q)\).
Recover the original \(q\) by \(H_\vartheta^{-1}\) and then divide
by these complete functions. The derivative at zero is nonzero;
compactness of the scalar circle gives one uniform injective disc.
For the normal orders \((a,c),(a,b),(b,c)\), the three whole
matrices are

\[
 D_0=\operatorname{diag}(\eta_j/\alpha_j,\kappa_j\alpha_j),\quad
 D_1=\operatorname{diag}(\kappa_j/g,g\eta_j),\quad
 D_2=\operatorname{diag}(\eta_j,\kappa_j).
 \tag{RFS.6}
\]

On the literal double their determinant winding is \(2t=24\),
and their first-entry windings are

\[
              t=12,\qquad l=(24(1-j),0,24).
 \tag{RFS.7}
\]

Divide each entry by its full reference value at the same \(q\)
and by the displayed angular factor. The quotient has a unique
based periodic logarithm on this protected disc. Retain those
entire logarithms and all their \(q\) derivatives.

For any integer \(d\), use the endpoint-flat based path

\[
 H_d(x,\vartheta)=
 (\sin\varphi(x)e^{id\vartheta}+j_{\rm quat}\cos\varphi(x))
 (\sin\varphi(x)+j_{\rm quat}\cos\varphi(x))^{-1},
 \quad\varphi(0)=0,\quad\varphi(1)=\pi/2.
 \tag{RFS.8}
\]

It equals one at \(x=0\) and both parameter ends, and equals
\(e^{id\vartheta}\) at \(x=1\). On
\(Q=n_0+j_{\rm quat}\bar n_1\) use
\(Q\mapsto H_tQH_{l_i-t}\), then compose the two full based-log
exponentials with the same radial cutoff. The radial endpoints
are exactly identity and \(D_i(q,0)^{-1}D_i(q,\vartheta)\).

For completeness, its interior inverse retains the original height.
On a bounded curve tube this map is \(F(n)=A n+G(n)\), where
\(A\) is its invertible value at \(q=0\). Since \(q=zn_0n_1\),
the complete unit dependence gives \(G=O(|n|^3)\) uniformly with
its first derivative. Choose the tube with
\(\|A^{-1}DG\|<1/2\). Iterating
\(n\mapsto A^{-1}n'-A^{-1}G(n)\) gives the unique smooth inverse,
and the source height is the actual \(zn_0n_1\) at that recovered
point. At the vertex use the \(H_\vartheta\) inverse from RFS.5.
The radial endpoint equalities glue the two constructions.

The normal rescaling \((z,n)\mapsto(z,\tau n)\), followed by the
cutoff flow of its conjugated family velocity, restores the full
nonlinear germs, with the compact inverse bound just proved. It
fixes the protected vertex collars, and the three open curve tubes
are disjoint. Carry the rounded boundary by this same isotopy.
The flow is in the closed doubled coordinates; it therefore
preserves the literal pulled-back seam and all source/cap arrows.
This constructs the actual ordinary double product, including its
inverse, before its free integers are evaluated.

\par

\medskip\noindent\textbf{RFS.3. The actual action and stable-derivative prisms}

\par\nobreak

The physical tangent weights are \(w=(0,1,-1)\) and the ordered
normal differences are \(d=(-2,-1,-1)\). Under RFS.8 the normal
action is

\[
 Q\longmapsto
 (H_t e^{-iw_i\phi/2}H_t^{-1})Q
 (H_{l_i-t}^{-1}e^{id_i\phi/2}H_{l_i-t}).
 \tag{RFS.9}
\]

The full unit maps and normal flow give its based homotopy to the
original action. The actual trace at \(q_H=1\), with
\(U=\cos x+j_{\rm quat}\sin x\), \(E_i=e^{iw_i\phi/2}\), is

\[
 E_i^{-1}UE_iH_tE_i^{-1}H_t^{-1}U^{-1}
                 H_{l_i-t}^{-1}e^{id_i\phi/2}H_{l_i-t}.
 \tag{RFS.10}
\]

Differentiating the displayed \(H_d\) gives the oriented degrees
\(-d,d\) of \(H_diH_d^{-1}\) and \(H_d^{-1}iH_d\).
The positive quaternion volume pairs angle \(c\phi/2\) and axis
degree \(e\) by \(ce\). The Maurer product boundary terms cancel
on the periodic parameter charts and have central constant radial
ends. Thus the normal-meridian and stable-derivative degrees are

\[
 I_i=tw_i+d_i(l_i-t),\qquad W_i=tw_i-d_i(l_i-t),
\]
\[
 I=(24(2j-1),24,-24),\qquad W=(24(1-2j),0,0).
 \tag{RFS.11}
\]

The stable lift is the lifted polar map to Spin: left quaternion
multiplication contributes its third degree and right multiplication
its negative. The base rotation depends only on \(\phi\), and
its endpoint Spin sign cancels the normal sign. Close the traces
by the actual vertex normal sections and link paths, capping their
parameter-independent end orbits by \(D^2_\vartheta\times S^1_\phi\).
The vertex maps factor through dimension two. The full
coordinate-to-ambient frame arrow extends over the annulus times
\(D^4\), so contributes zero on these capped three-cycles. This
gives the actual marked prism with the full RFS.6 frames.

For \(\gamma_i=\partial x_i\), the oriented positive-circle
identity \(\partial\Delta=-\Delta\partial\) yields

\[
 V(w_ie_i-w_0e_0)=I_ie_i-I_0e_0,\qquad
 w_i\delta_i-w_0\delta_0=W_i-W_0.
 \tag{RFS.12}
\]

Hence the actual two selected entries are

\[
 V_{11}=24,\quad\delta_1=24(2j-1),\qquad
 V_{22}=24,\quad\delta_2=24(1-2j).
 \tag{RFS.13}
\]

The ambient-to-local frame is literally \(c=u^{-1}q_H\).
Its derivative contains \(-u^{-1}(du)u^{-1}q_H\), removed by the
invertible triangular homotopy. Its normal-meridian third class
is \(+1\); in the K23 oriented relative disc it contributes
\(-V_{ii}\). Thus, in these same products,

\[
 2n_i=\delta_i-V_{ii},\qquad m_i=-V_{ii}-n_i.
 \tag{RFS.14}
\]

The ordinary double pairs are therefore

\[
 C_1:(-24j,24(j-1)),\qquad
 C_2:(24(j-1),-24j).
 \tag{RFS.15}
\]

No fixed \(C_0\) value has entered these equations.

\par

\medskip\noindent\textbf{RFS.4. The native correction, actual square and full flop}

\par\nobreak

The native quaternion is
\(q_H=U(z)(n_0+j_{\rm quat}\bar n_1)\), with
\(U(z)=(1+j_{\rm quat}z)/(1+|z|^2)^{1/2}\).
The actual native-to-ordinary boundary arrow of RFS.6--8 is

\[
 q_H\longmapsto(UH_tU^{-1})q_HH_{l_i-t}.
 \tag{RFS.16}
\]

For \(C_1\), its native second normals are
\((gw,(g\eta_j)^{-1}u)\), and the ordinary normals are
\((\kappa_jw,u)\). Their exact ratio is \(D_1\), so RFS.16
has retained every \(g,\kappa_j,\eta_j\) factor. The \(C_2\)
ratio is directly \(D_2\). The full based unit logarithms and
the inverse height maps are those in RFS.2.

The left factor in RFS.16 has degree \(2t=24\) on the
Hopf-oriented sphere times parameter; direct substitution of
RFS.8 in its Maurer boundary integral gives twice the angular
winding. The right factor has only radial and parameter variables
and third degree zero. Extend the inverse factors over the
exterior boundary cylinder by its explicit radial retraction;
the endpoint degrees are \((-24,0)\). They extend over
\(D^3\times D^4\), so no disc-supported central class is added.
Thus the ordinary return is the native return followed by
\(L^{-24}\), and the native doubles are

\[
 C_1:(24(1-j),24(j-1)),\qquad C_2:(24j,-24j).
 \tag{RFS.17}
\]

Use the connection defining the single native product of RFS.1.
Its literal pullback has endpoint \(f^2\), with both laps of the
same original and cap closing. The corrected ordinary connection
has the same native boundary connection. Affine interpolation
relative to that boundary gives an isotopy of relative endpoints.
The free-coordinate homomorphism to \(\mathbb Z^2\) takes
\(f^2\) to twice the coordinates of \(f\); multiplication by two
in this free group is injective. This proves RFS.1. The central
coordinate is not divided.

The full C1 opposite-ruling chart is RFB.13, with Cox lift

\[
 (a_0,a_1,c_0,c_1)\mapsto
 (a_0,Da_1,(AD)^{-1}c_0,Ac_1).
\]

It has the full reciprocal inverse at the same \(q\), retains the
positive moduli by the transported Hermitian metric, and induces
the actual target \(T_{j-1}\) through RFB.14. The C2 lift is
\((a_0,(AD)^{-1}a_1,Ac_0,Dc_1)\), and its actual target is
\(T_{j+1}\). For arbitrary four Cox angular weights the exact
link is IF3.17: it retains the full factor
\(e^{hS\vartheta/2}\), or equivalently \(e^{kS\vartheta/2}\),
and both paired central endpoints. Its exterior is precisely
RFB.17, with the written inverse and Hopf inverses.

In these actual relative products the full K23 identity is

\[
 T(L^mR^nC^a)T^{-1}=L^{-n}R^{m+2n}C^{a+n}.
 \tag{RFS.18}
\]

The fixed anti-quaternion reset restores the first-ruling complex
normal convention, swaps the two free entries and preserves the
central entry. Hence the true central change is \(n\), giving
RFS.2 in the two stated directions. For example the C1 target
free pair after that full reset is exactly the C2 pair at \(j-1\);
the reciprocal C2/C1 check is identical by the displayed equations.
All exterior source/root/cap/sector arrows are those in RFB.3,
extended by identity away from the selected tube. This proves the
signed scalar fan comparison on the actual named reference.

The companion \path{code/part3/e6_selected_i3_m2_scalar_native.py} checks
the entire trace, variation/derivative relations, ordinary/native
corrections, both target resets, and signed recurrence, including
the full coefficient squares on the common Cox chart. It locks the
complete geometric inputs and the RFB source separately.


% End expanded verified cp36 input


% Expanded verified cp36 input: notes/part3_e6_selected_i3_m2_return_support_19.tex
\par\bigskip\noindent\textbf{P3K.111 supporting proof 19: Full original attachment, parameter square and outer transfer}\par\nobreak
\par

\medskip\noindent\textbf{Selected-I3 \(M_{2}\): full compact attachment, same-fan reference and physical source return}

\par\nobreak

Working proof, 18 November 2026. We retain the original source
\[
 y^2=x^3+r_{\rm src}x^2+tx+t^2,\quad G=(0,t),\quad
 s=1/4-r_{\rm src},\quad M_2=\mathcal O([2G]-[O]),
 \tag{ASI2.1}
\]
translation by \(-3G\), the selected height \(q=t-t_c(s)\),
the good cubic cap \((m,n)=(3,1)\), and its prescribed finite
line character and full reference cochain. The positive original
\(s\)-meridian is the one in SI2.1. The actual selected native
star is the one derived from the two protected zero-end facets
in NS.1--3. Its point-coordinate matrix is
\[
 P_N=\begin{pmatrix}1&-1\\-1&2\end{pmatrix};
 \tag{ASI2.2}
\]
in the retained convention \(T_j=T(e_1+je_2,e_2)\), this is
\(T_{-1}\), by the literal triangle and period map RFB.4--7.

The auxiliary filling uses \(\lambda=\lambda_0s^{12}\) and
scalar root \(r_0s^4\), \(r_0^3=\lambda_0\), with \(\lambda_0\)
in the specified common strict contraction range. Choose a
small positive basepoint \(s_*\), and retain
\(r_{\rm base}=r_0s_*^4\), \(\lambda_{\rm base}=r_{\rm base}^3\).
The return computed below keeps this original scalar fixed.
The completed Jin physical pair and index are
\[
 (k_M,\ell_M)_{\rm source}=(71,-4),\qquad
 \frac{4k_M-\ell_M}{24}=12,\qquad
                 [\Phi_{\rm source}]=12\ \hbox{in }\mathbb Z/28.
 \tag{ASI2.3}
\]
We prove the attachment and directed reference comparison before
using the pair in ASI2.3. All auxiliary and physical lines retain
their stated Hopf orientations.

\par

\medskip\noindent\textbf{ASI2.1. The actual completed periodic piece and compactness}

\par\nobreak

The complete functions and point coefficient are
\[
 Q=P^2R,\quad R-P=su_s,\quad t=Pu_t,\quad q=Rc_q,\quad
 \eta_2=a_T\Theta'_Q(1)\frac{\Theta_Q(P^2)}{\Theta_Q(P)^2},
 \quad V_{\rm reg}=\eta_2^{-1}V_{\rm pt}.
 \tag{ASI2.4}
\]
The full Miller telescope in SI2.6--9 is
\[
 B_{2,3}=f_2(p)f_2(p-3G)f_2(p-G)^2f_2(p-2G)^2
 =\eta_2^6\frac{f_{P^2}(Z)}{f_{P^2}(P^{-3}Z)},\quad
 f_2=\frac{y-t-x/2}{x-s}.
 \tag{ASI2.5}
\]
Thus the genuine point periods are
\[
 h:(Z,V_{\rm pt})\mapsto(P^2RZ,P^2V_{\rm pt}),\quad
 B:(Z,V_{\rm pt})\mapsto
       (P^{-3}Z,\lambda q\eta_2^6t^{-3}V_{\rm pt}).
 \tag{ASI2.6}
\]
Their inverses retain the same full functions at the recovered
source point. The entire \(\Theta'_Q(1)\), with central value
\(-1\), remains in \(\eta_2\) and its reciprocal.

Use the complete new common fan and its zero, one and infinity
endpoints constructed in supporting proof 20.  Their common boundary
has 1,926 primitive regular three-cones.  The endpoint counts are
54,121, 2,022 and 2,381.  The complete original-lattice semigroup maps,
full character inverses, paired period facets and source cuts are all
retained in that construction and its exact certificates.  Each
subdivision is proper over the same original convex raw cells.
The actual base maps are
\[
       (P,R)=(P,pP),\quad(P,(1+v)P),\quad(Ru,R).
 \tag{ASI2.7}
\]
Their inverses on the overlap are \(p=R/P,v=R/P-1,u=P/R\).
The original \(s,t,q\) are the entire functions in ASI2.4.
We use this executed total-line completion with its physical
circle and derivative maps. The source/Cartier graphs below
specify its actual finite boundary attachment.

Here is an analytic compactness check with those full coefficients.
In the balanced tensor coefficient \(V_{2,T}=V_{\rm reg}/Z\),
the real logarithmic period columns in the order
\((\log|Z|,\log|V_{2,T}|)\) are
\[
 \begin{gathered}
 \binom{-2A-C}{C},\qquad\binom{3A}{-D_B},\\
 A=-\log|P|,\quad C=-\log|R|,\quad
 D_B=-\log|\lambda qK|,\quad K=(P/t)^3\eta_2^6 .
 \end{gathered}
 \tag{ASI2.8}
\]
Their determinant is
\[
                    (2A+C)D_B-3AC .
 \tag{ASI2.9}
\]
The full units in \(K,c_q,u_s\), and their reciprocals,
are bounded on one smaller closed original bidisc. The exact
equation \(R-P=su_s\) gives
\[
 D_B\geq12\min(A,C)+C+C_*
 \tag{ASI2.10}
\]
after choosing the stated strict contraction range for
\(\lambda_0\) and then shrinking the bidisc; \(C_*>0\) is
fixed. Specifically a sufficient choice of
\(-\log|\lambda_0|\) is larger than the supremum of
\(\log|c_qK|-12\log|u_s|+12\log2\), plus \(C_*\).
Consequently ASI2.9 is bounded below by
\[
 \begin{cases}
 23AC+13C^2+C_*(2A+C),&A\geq C,\\
 24A^2+11AC+C^2+C_*(2A+C),&A\leq C.
 \end{cases}
 \tag{ASI2.11}
\]
The exact inverse matrix therefore gives the marked real period
coefficients. Their values in \([0,1]^2\), together with the
compact angular variables, give a fundamental cut.

For completeness the finite-support comparison is explicit.
In the three endpoint lattices, the valuation period matrices are
\[
 M_0(r,z)=
 \begin{pmatrix}10r+z&8r+z\\3r&6r+z\end{pmatrix},\quad
 M_1(r,z)=
 \begin{pmatrix}3r&0\\r&12r+12z\end{pmatrix},\quad
 M_\infty(r,z)=
 \begin{pmatrix}23r&11r\\3r&3r+z\end{pmatrix}.
 \tag{ASI2.12}
\]
The zero support is the six-vertex polygon
\[
 (0,0),(-5r-z,0),(-10r-z,-3r),
 (-13r-2z,-6r-z),(-8r-z,-6r-z),(-3r-z,-3r-z),
 \tag{ASI2.13}
\]
the one support is the parallelogram generated by the two columns
of \(M_1\), and the infinity support is
\[
 (0,0),(-12r,0),(-23r,-3r),(-23r,-3r-z),
                         (-11r,-3r-z),(0,-z).
 \tag{ASI2.14}
\]
The matrices and polygons retain their original lattices; the
index-six and index-thirty-three inspection maps used to derive
them are inverted only on their stated image lattices.
Direct multiplication by \(M_0^{-1}\), respectively
\(M_\infty^{-1}\), puts every polygon vertex in \([-2,2]^2\).
For example the nonintegral zero coordinates have numerators
\(-30r^2-11rz-z^2,\ 15r^2+3rz,\ -21r^2-10rz-z^2,\ 6r^2+2rz\)
over \(36r^2+13rz+z^2\). At infinity they have numerators
\(-36r-12z,\ 36r,\ -23z,\ 11z\) over \(36r+23z\).
The asserted inequalities follow coefficient by coefficient
after multiplication by these positive denominators. The one
parallelogram has period coordinates exactly \([0,1]^2\).
Their period translates tile the stated native support, with the
full unit arrows recorded in the parent construction. Thus a
fixed finite set of translates, for example indices in
\([-2,3]^2\), covers the cut in valuation coordinates.

The full logarithms perturb those cuts through the exact matrix
ASI2.8. On each endpoint domain, the inverse of ASI2.12 is
uniformly bounded by a constant times the reciprocal of the
smaller retained nonzero original log-radius. Shrinking that
radius bounds the product of this inverse with the complete
unit-log perturbation by \(1/2\). The inverse
\((I+M^{-1}\Delta)^{-1}=\sum_{k\geq0}(-M^{-1}\Delta)^k\)
then gives the actual norm-coordinate comparison. Hence only
a fixed finite enlargement of those period translates meets
the exact cut. All their coordinate moduli lie in closed
polydiscs of the finitely many proper parent charts.
On the one lower rank-one face the remaining source period is
the full \(L:(X,Y)\mapsto(P^3pX,PpY)\); its source multiplier
is a strict contraction. On the infinity lower rank-one face
the remaining period is the full \(L:(X,Z)\mapsto
(R^{23}f_\infty^2X,R^3Z)\), again with a strict source
contraction. The chain labels and their source/line factors
are retained. Those contractions prove local finiteness on
these faces and give their actual completed source cycles.
Away from them the inverse real period matrix proves local
finiteness directly. The separated parent charts, their
proper subdivisions and these locally finite quotient maps
therefore give a Hausdorff compact cut.

Choose a closed original \(t\)-disc \(U\) containing \(0\)
and \(t_c(s)\), and excluding the other simple singular value
and the actual infinity cap. Its boundary misses all source
critical values for the closed small source-parameter disc.
It retains every horizontal point/O section chart. Round
its two base corners in the actual monomial coordinates
with an invariant radial function. The compact resulting
collision piece is denoted \(W_c\).

\par

\medskip\noindent\textbf{ASI2.2. Whole point/O boundary, all source discs and the reference body}

\par\nobreak

For a source band \(j\), put \(x=Q^jZ,y=Q^{1-j}/Z\).
The full graph and point/O coefficients are
\[
 x=P^2a,\ ay=R;\qquad xc=P^2,\ y=Rc,\qquad
 P=uw,\ (x,c)=(u,uw^2)\ \hbox{or }(u^2w,w),
 \tag{ASI2.15}
\]
\[
 V_{\rm reg}=-\eta_2^{-1}P^{-2j}E A
   =\eta_2^{-1}P^{-2j}c^{-1}EC,\quad C=-cA,
 \qquad
 V_{\rm reg}=\eta_2^{-1}P^{-2j-2}(x-P^2)E_OC_O.
 \tag{ASI2.16}
\]
Here \(E,E_O\) are the complete convergent products SI2.17--19.
The first graph charts have \(ds/dP\) a unit with value \(-1\);
the two second charts have \(ds/dR\) a unit with value \(+1\).
Every inverse recovers the original \(P,R\) first, and then
uses the reciprocal coefficient in the same Cartier module.
The additional \(P^{-2}\) at \(O\) remains. On the finite
original \(t\)-boundary \(P\ne0\), these are precisely the
original source and line maps through the section charts.
At a section coincidence the module numerator and denominator
are carried as a line-module family.

Let \(W_{r,1}\) be the actual RB scalar-disc reference on the
fixed source \(s_*\), with scalar \(\lambda_0\sigma^{12}\),
root \(r_0\sigma^4\), and the same finite character and
sector integers. The ordinary chain is \(U_kW_k=\sigma\).
At the selected node its full equation is
\[
             U_jW_j=\sigma^{12}q,\qquad
                       x_i\bar y_i=q,
 \tag{ASI2.17}
\]
with the ordered factors
\((\sigma,\ldots,\sigma,\bar y_i,x_i)\), twelve copies
of \(\sigma\). For factors \(f_1,\ldots,f_{14}\) its
fourteen charts are
\[
 f_r=s_rt_r,\quad
 U=s_r\prod_{\nu<r}f_\nu,\quad
 W=t_r\prod_{\nu>r}f_\nu .
 \tag{ASI2.18}
\]
On each adjacent overlap the inverse divides by the
retained nonzero factor. RB.14--17a are their literal
source-band, point-line and period maps. This includes the
full \(q\), \(\kappa,\alpha,g,\Gamma\) factors.
On a source-unit patch it reduces by those displayed
coordinate divisions to the ordinary chain. The twelve
fixed source components have \(\alpha_l=\beta_l=0\);
the overlap with the selected moving curve requires the
specified nonzero \(v_{12}\). Thus every source component
retains all three source vanishing discs. The reference\textquotesingle s
moving component is the distinct fixed sphere times its
parameter disc supplied by the last two factor charts.
Its selected fan is \(T_1\), with matrix
\(\left(\begin{smallmatrix}1&-1\\1&0\end{smallmatrix}\right)\).

The side boundary from \(W_c\) to this actual reference is
the whole IL.26--31 map. In the one characters
\(X=Z^{-1},Y=V_{\rm reg}/Z,v=p-1\), it is
\[
 \alpha_l=Yv^{-l},\quad\beta_l=v^{l+1}/Y,\qquad
 A_l=H_{+,l}\alpha_l,\quad B_l=H_{-,l}\beta_l,
\]
\[
 H_{+,l}=\eta_2u_s^lP^{-l}X^{-1},\quad
 H_{-,l}=\eta_2^{-1}XP^{l+1}u_s^{-l-1},\quad
                    H_{+,l}H_{-,l}=P/u_s.
 \tag{ASI2.19}
\]
For an arbitrary label \(12j+l\), the actual inverse
source period is
\[
 X_*=p^jX,\quad Y_*=P^{-12j}v^{-12j}c_{12}^{-j}Y.
 \tag{ASI2.20}
\]
The two target coefficients in ASI2.19, evaluated at
that recovered source point, also carry
\(\gamma^j,\gamma^{-j}\), with
\[
                 \gamma=pu_s^{12}c_{12}
                        =P^2\lambda_0q\eta_2^6/t^3.
 \tag{ASI2.21}
\]
These follow from the full \(T=hB\) arrow, including
its additional source \(h\) factor. They preserve
\(A_lB_l=s\) and have their stated reciprocal inverse.

The source transport is the actual compact relative
flow \(T_s:S_0\to S_s\) of IL.28--30, with final target
\(K_s=T_{s_*}T_s^{-1}\). Its boundary formula holds
the marked real periods fixed:
\[
 Z_s=e^{b\operatorname{Log}(Q_s/Q_0)}Z_0,\qquad
 V_{{\rm pt},s}
      =e^{b\operatorname{Log}(P_s^2/P_0^2)}V_{{\rm pt},0}.
 \tag{ASI2.22}
\]
Both logarithms have value zero at \(s=0\); under the
source deck \(b\mapsto b+1\), their targets acquire
exactly \(Q_s,P_s^2\). The local vector fields
\((\partial_Ps)^{-1}\partial_P\) and
\((\partial_Rs)^{-1}\partial_R\) on ASI2.15
project to one. A finite relative partition, agreeing
with the derivative of ASI2.22 on the boundary,
gives the complete source flow; its backward flow is
the inverse.

For genuine vector frames \(e_j=g_{ij}e_i\), its line
connection has
\[
 a_i=-\sum_k\chi_k\,Xg_{ik}/g_{ik},\quad
 a_j=a_i+Xg_{ij}/g_{ij},\quad
              X^{\rm hor}=X-a_iw_i\partial_{w_i}.
 \tag{ASI2.23}
\]
A relative affine correction makes this equal to the
full ASI2.16 boundary derivative. The line ODE and
its reciprocal dual ODE supply the entire target
\((K_sz,K_s^{\rm line}H_+\alpha,
K_s^{\rm dual}H_-\beta)\), with inverse first
\(K_s^{-1}\) and then the reciprocal coefficient.
On every period translate it is conjugated by the
full original period and its inverse. This retains
the derivative of every unit and integral source
deck. Since \(v=\alpha\beta\), its first normal
derivative on the fixed locus is zero. Circle
averaging keeps the source lift and its complex
diagonal normal derivative. This matches the entire
side-boundary tangent and line frames.

The actual three source discs give
\[
 H_1(\partial S)=\mathbb Z\langle b,\alpha_t\rangle
                 \oplus\mathbb Z/3\langle a\rangle
       \longrightarrow H_1(S)=\mathbb Z\langle b\rangle,
 \quad (a,\alpha_t,b)\mapsto(0,0,b).
 \tag{ASI2.24}
\]
The full reference cube frame is
\(e^{-2\tau_L\operatorname{Log}p}X^2V_{\rm pt}^3\),
\(\tau_L=\log|X|/\log|Q|\), with inverse divisor
\(6F-2C_1-C_2\). Its order-three flat boundary connection,
the primitive \(P/\epsilon\) fibre clutch, the complete
point/O modules and all \(\gamma^j\) are those of IL.33--38.
Thus ASI2.19 is the actual boundary map used in the
twelve closed normal doubles, not a choice of their
boundary degree.

\par

\medskip\noindent\textbf{ASI2.3. The original finite cap and full exterior cylinder}

\par\nobreak

On the infinity root keep the literal source map
\[
 q_\infty=z^3,\quad t=z^{-3},\quad
 x=z^{-2}X',\quad y=z^{-3}Y',\quad
 Y'^2=X'^3+r_{\rm src}z^2X'^2+zX'+1.
 \tag{ASI2.25}
\]
The deck is
\((z,X',Y')\mapsto(\zeta_3z,\zeta_3^{-4}X',Y')\).
The marked sections are
\(G'=(0,1)\), \(2G'=(sz^2,-1-sz^3/2)\).
The actual point correction on the two quotient graph
rays \((2/3,1/3),(1/3,2/3)\) is
\[
       \beta^*\mathcal O([2G']-[O])
                    =\pi^*M_2+E_P-E_O.
 \tag{ASI2.26}
\]
Its local monomial coefficient is
\(\min(v_w,2v_z)\) at the two marked sections,
with the specified Cartier generators and reciprocal.
The small division \(3R'_0=3G'\) has
\[
 \tau(3G')=
 z^2\frac{-2s(2+sz^3)}
 {8+6sz^3+(s^2-s^3)z^6},\qquad
 \tau(R'_0)=sz^2\ell_R,\quad\ell_R(0,0)=-1/6.
 \tag{ASI2.27}
\]
It preserves the actual graph ideals
\((w,z),(w,z^2),(z,w^2)\), with its inverse.
The normalized chord from SI2.26 has value one at
\(\tau=0\), has its stated horizontal divisor, and
commutes with the deck because the chord and
\(\tau\) weights cancel. Write this full Hom as \(F\).
The complete root coefficient is
\[
 B'_{\rm sel}=h_\infty B'_{2,3},\quad
 h_\infty=1-t_c(s)z^3,\quad
 \beta=B'_{\rm sel}/F^3,\quad
 \widetilde F=\beta^{1/3}F,\quad
 A_\lambda=r_\lambda\zeta_3^d\widetilde F D_3 .
 \tag{ASI2.28}
\]
The cube is chronological. The analytic invariant
cube root is fixed by its central value, and
\(A_\lambda^3=B_\lambda=\lambda B'_{\rm sel}\)
as a full source and line map.

On an original infinity sector set
\(\ell_\alpha=\log_\alpha q_\infty/(2\pi i)\),
\(z_\alpha=e^{2\pi i\ell_\alpha/3}\). The actual
annular arrow is
\[
 \Phi_\alpha=
 \operatorname{Exp}\{\ell_\alpha\operatorname{Log}_\alpha
       \widetilde F+
       \ell_\alpha(\log r_\lambda/(2\pi i)+d/3)\delta
       +h_\alpha\}\,I_{z_\alpha}.
 \tag{ASI2.29}
\]
It satisfies
\(\Phi_\beta=A_\lambda^{j_0}
B_\lambda^{j_{\alpha\beta}}\Phi_\alpha\) and
\(\Phi_\alpha B_\lambda=B_\lambda\Phi_\alpha\).
The map \(I_z\) carries ASI2.26; the original reference
cochain \(h_\alpha\) is the fixed-contour, reference-zero
solution SI2.28/OG.12--13 with the entire source periods
and line lattice. No integer sector term is reset.
Its inverse reverses the actual Lie, root and
Cartier arrows at the recovered source point.

On the outer scalar circle the finite cap is a compact
proper family over the source interval. Extend the
prescribed side-boundary source field through its
actual root charts, agreeing on its annulus by the
derivative of ASI2.29. The difference of two such
lifts is vertical; a relative partition and finite
deck average preserve the prescribed interval
component. The full line connection is ASI2.23 in
the actual corrected modules. Compactness gives
its interval flow and backward inverse. On the
ordinary exterior use the original smooth source
charts and the same construction. On a period
translate use the actual conjugating map and its
inverse, so the complete functions and derivatives
remain. This produces the actual exterior cylinder
between the coupled and fixed-source families.
It is constructed on the nonzero outer scalar
circle, where all original cap maps are defined.

The scalar root coordinate \(\rho\) has the exact
forward square
\[
             \Phi_{\rho+4}q^4=h_\infty^4\Phi_\rho,
 \tag{ASI2.30}
\]
with endpoint-to-initial arrows \(q^{-4},h_\infty^{-4}\).
It follows from the entire identity
\(q_\infty q=h_\infty\), including ASI2.29\textquotesingle s
integral cochain. The field just constructed
transports this square literally. At \(q=0\) the
ordinary arrow is its actual \(B^4\) reindex followed
by the full bounded source/line factor NS.24. Thus
the selected zero is retained rather than divided out.
This closes every original ordinary/root boundary
of the relative piece \(W_c-W_{r,1}\).

\par

\medskip\noindent\textbf{ASI2.4. Both native Hopf caps and the entire inner ledger}

\par\nobreak

The two actual collision native charts are
\[
 q=Za_0c_0=z a_\infty c_\infty,\quad
 z=Z^{-1},\quad(a_\infty,c_\infty)=Z(a_0,c_0),
 \quad c_0=ZV_{\rm pt},\quad a_0=q/(Z^2V_{\rm pt}).
 \tag{ASI2.31}
\]
NS.8--15 give their holomorphic extension and inverse
at both vertices, recovering the original \(s,q\)
and then every full unit. Their original volume is
\[
 \Omega_{\rm orig}=\frac2{a_T}\,
 ds\wedge dZ\wedge da_0\wedge dc_0
 =-\frac2{a_T}\,ds\wedge dz\wedge da_\infty\wedge dc_\infty .
 \tag{ASI2.32}
\]
The endpoint is
\[
 E_0=(P^{12}Z,P^{12}\mathcal A^{-1}a_0,
                         P^{-24}\mathcal A c_0),
       \qquad\mathcal A=(c_q/f_0)^4 .
 \tag{ASI2.33}
\]
It preserves the whole \(q\), and its reciprocal
inverse evaluates the coefficients at that same
recovered original \(s,q\).

The NC cut map has
\[
 C_\theta(Z,a,c)=
      (\alpha_\theta Z,(\alpha_\theta\beta_\theta)^{-1}a,
                                       \beta_\theta c),
 \quad C_1E=C_0,
 \tag{ASI2.34}
\]
where its two full logarithms are the NC.5 logarithms
of the corresponding two factors in ASI2.33, with
the specified endpoint-flat cutoff. Its velocity
annihilates \(q\). Cut off only that vertical
velocity, keeping the original parameter-horizontal
lift, and integrate on both native charts. The
backward flow is the inverse. Its support misses
the source boundary and the actual infinity cap.

Use the held Cox quaternion
\[
 T_Z=\begin{pmatrix}1&-\bar Z\\Z&1\end{pmatrix},\quad
 q_H=T_Z(a,\bar c)^t,\quad
 T'_{1/Z}\operatorname{diag}(Z,\bar Z)=T_Z.
 \tag{ASI2.35}
\]
The entire normal matrix is
\[
 G_\theta=T_{\alpha_\theta Z}
       \operatorname{diag}((\alpha_\theta\beta_\theta)^{-1},
                              \overline{\beta_\theta})T_Z^{-1}.
 \tag{ASI2.36}
\]
Its inverse is the reversed matrix product at the
recovered base point. On the Hermitian three-ball
the exact fixed-sphere map is
\[
 \Psi_\alpha(Y)=
 \frac{\operatorname{diag}(\alpha,1)Y
                    \operatorname{diag}(\alpha,1)^*}
 {\operatorname{tr}(\operatorname{diag}(\alpha,1)Y
                    \operatorname{diag}(\alpha,1)^*)}.
 \tag{ASI2.37}
\]
Its denominator is positive and its inverse is
\(\Psi_{\alpha^{-1}}\). Its rank-one boundary is
the actual \(Z\mapsto\alpha Z\) map. In ASI2.36
multiply the two full cut logarithms by a radial
function zero near the ball centre and one near
the boundary. The resulting \(\widehat G_\theta\)
is direction independent near the centre. The
map
\[
 \widehat C_\theta(Y,q_H)=
 \left(\Psi_\alpha Y,
       \widehat G_\theta(Y)q_H/
                         \|\widehat G_\theta(Y)q_H\|\right)
 \tag{ASI2.38}
\]
has inverse first \(\Psi_{\alpha^{-1}}\), then
\(\widehat G_\theta^{-1}\) at that recovered
point. Its seam is
\(\widehat E=\widehat C_1^{-1}\widehat C_0\).
For the actual resolved normal derivative \(N_\theta\)
of NS.23, the directed equality is
\[
          B_\theta=C_\theta^{-1}N_\theta,\qquad
                    \widehat C_\theta B_\theta=N_\theta.
 \tag{ASI2.39}
\]
The reciprocal equality uses the inverse maps in
the reversed order. This attaches the actual
product Hopf cap with every normal coefficient
and positive spinor frame retained.

The reference has its own regular native
coefficients \(c_R=V_{\rm pt}/Z,a_R=q/V_{\rm pt}\),
its exact second chart from RB.17a, and its own
endpoint MN.12. Applying ASI2.34--39 with those
full reference factors is exactly its specified
own cap, as proved in MN.3. We do not identify
the two native vertex charts by the dense
\(Z^2\) expression; their complete reference
continuation is constructed in ASI2.5 below.

The complete new fixed-face inventory and its exact ownership are
those of supporting proof 20 D.  It contains 35,908 compact physical
components, the twelve complete lower source owners and the single
zero moving owner.  Every endpoint chart belongs to exactly one
generic, compact-corner or separately capped piece.  The actual
physical face and period maps, all finite slice weights and the
same-atlas SHA certificate are retained.

The compact physical pair is \((895,139)\); the twelve lower
pairs sum to \((12,-24)\).  For the new moving surface the
full NC.19--20 bordism carries its actual cap to the closed
replacement with roots \(x=h+D_-\), \(y=-h-D_+\).
The thirty-one-chart derivative, complete divisor matrix and all
twenty-nine chronological source blowups give
\(c_1(E)^2=-1200\), \(c_2(E)=-294\), \(p_1(TF)=-87\).
The original RB reference roots remain \(h-12u,-h-12u\),
and its own cap still has \(c_1^2=c_2=p_1=0\).
From the actual physical Euler denominator,
writing \(x=-c_1(N_-),y=c_1(N_+)\), the
physical surface contributions are
\[
 k_M=c_1(E)^2-c_2(E),\qquad
 \ell_M=p_1(TF)+c_1(E)^2-4c_2(E).
 \tag{ASI2.41}
\]
These are IL.41 after the literal negative-line
conjugation. Hence the actual new moving-minus-own-reference pair is
\((-906,-111)\).  The complete new inner ledger with the same
RB \(T_1\) reference is
\[
 (895,139)+(12,-24)+(-906,-111)=(1,4).
\]
This uses the new collision atlas for every collision contribution
and includes the actual original fixed-source reference exactly once.

\par

\medskip\noindent\textbf{ASI2.5. Positive outer Stokes, Jin reversal and the actual reference change}

\par\nobreak

On the free quotient use the physical nonzero-vector
line \(M\). On a finite stabilizer chart use its
actual covering slice and reciprocal covering degree.
Average the tangent and line connections over the
physical circle and finite deck group. Across each
identified boundary transport those connections
through the full derivative and line arrow above;
place the identical interpolation on the opposite
collar with its interval reversed. Its complete
transgression cancels pointwise. The stable tangent
splitting into the quotient tangent and the
physical orbit direction is the actual one in
these slice charts. Stokes therefore applies to
\(c_1(M)^3\) and \(p_1c_1(M)\), including every
finite covering factor in the new complete fixed-face certificate and IL.

The orientation is fixed by the original complex
source and the physical circle:
\[
 \operatorname{or}(W)=d\phi_\delta\wedge
           \operatorname{or}(Y)\wedge dr_s\wedge d\theta.
 \tag{ASI2.43}
\]
The outward parameter boundary in the circle-first
quotient is \(-\operatorname{or}(Y)\wedge d\theta\),
as the radial normal crosses five directions.
The removed fixed tubes have the two opposite
reversals of the actual normal-Hopf boundary map
ASI2.35. These give exactly the complex
fixed-surface orientations used in supporting proof 20, IL and
the retained native/cap maps. Equations ASI2.32 and NC.20 fix the native
and cap signs, including the original nonzero
volume unit. Consequently, with \(O\) denoting
the completed positive-parameter physical pair,
\[
                 O_c-O_{r,1}=(1,4).
 \tag{ASI2.44}
\]

For the same actual lifted relative return \(f\),
the map
\[
 \begin{gathered}
 Q_f^+=(Y\times[0,1])/((y,1)\sim(f(y),0))\\
 \longrightarrow M_f=(Y\times[0,1])/((y,0)\sim(f(y),1)),\\
 [y,t]\longmapsto[y,1-t].
 \end{gathered}
 \tag{ASI2.45}
\]
has line map \([v,t]\mapsto[v,1-t]\), with the same
full endpoint lift. It sends the two displayed
endpoint relations to one another. Its inverse
is the identical formula. It preserves the
physical line and reverses the parameter
orientation. Write \(J=-O\) for the resulting
Jin-oriented pair.

The whole RFB construction has the precise type
\[
                    J_{r,-1}-J_{r,1}=(-72,0).
 \tag{ASI2.46}
\]
Indeed RFB.33--40 computes with \(U=M^*\);
its physical dual and original-to-Jin map give
\(k_M=+\int F_U^3\) in that coefficient
orientation, explicitly as in RFB.40--43.
Thus in the positive outer orientation
\[
                    O_{r,-1}-O_{r,1}=(72,0).
 \tag{ASI2.47}
\]
These are directed equalities on the entire
fixed-source scalar reference word \((0,4)\).

We recall its actual morphisms to specify its
use here. For \(j=1,0,-1\),
\[
 Z_j=Z,\quad V_j=Z^{-j}V_{\rm pt},\quad h_j=hB^j,\quad
 \kappa_j=\kappa g^j,\quad
 \alpha_j=\alpha\eta^j\kappa^{-j}g^{-j^2},\quad
 \eta_j=\eta g^{-j}.
 \tag{ASI2.48}
\]
Here \(\eta=\sigma^{12}\lambda_0\eta_2^6t^{-3}\),
\(\kappa=P^2/c_q,\alpha=P^2,g=P^{-3}\), all at
the fixed original source. Every inverse is
the displayed reciprocal character and period
substitution. Each reverse \(C_1\) change has
the full Cox map
\[
 (a',b',c')=(a/b,b,bc),\quad
 (a,b,c)=(a'b',b',c'/b'),
 \tag{ASI2.49}
\]
on its dense overlap, and the actual two
projections of \(\operatorname{Tot}\mathcal O(-1,-1)\)
\[
 p_+(A,C,\tau)=[A_0,A_1,\tau C_0,\tau C_1],\quad
 p_-(A,C,\tau)=[\tau A_0,\tau A_1,C_0,C_1].
 \tag{ASI2.50}
\]
They pull \(q\) back to
\(\tau^2A_0A_1C_0C_1\); their inverses off the
zero sections are the specified projective
direction and nonzero fibre coefficient.
The entire bounded Cox lift and its inverse
are RFB.17. The six chart rows RFB.7 retain
every \(\kappa_j,\alpha_j,g,\eta_j\).
The common link/exterior map is
\((y,q_H)\mapsto(q_H^{-1}yq_H,q_H)\), with
its conjugation inverse and unrestricted
four-weight Hopf lift.

Their exact cut square is
\(C_jS_j=S_jC_0\), where
\[
 S_j=(Z,Z^ja,Z^{-j}c),\quad
 \mathcal D_j=\mathcal D\mathcal A^{-j}.
 \tag{ASI2.51}
\]
The full inverse square and the same
Hermitian-ball construction ASI2.37--39
carry it through both Hopf caps. The
ordinary cubic cap and ASI2.30 are fixed.
RFB.24--38 construct the completed quotient
map and inverse by the actual log-norm
period coefficients, source/graph flows,
Cartier spinor arrow and relative Cech
line map. The latter is
\(U_j\to I_j^*U_0\otimes\mathcal T^{12j}\)
with the exact endpoint-to-initial clutch
\(w\mapsto e^{2\pi ib}w\), all outer
primitive \(-j(du+\alpha_qdv)\) retained.
The full nilpotent tangent collar has
zero \(p_1\) transgression. Its relative
regular-level cobordism and the entire
Euler collar give RFB.43, and the actual
integer \(0\to1\) Cox comparison fixes
its linear coefficient. These maps prove
ASI2.46 for this reference and its own
three native caps. They are the required
finite continuation through \(T_1,T_0,T_{-1}\).

At \(j=-1\), \(a_j=q/(Z^2V_{\rm pt})\),
\(c_j=ZV_{\rm pt}\). Hence this final
reference has exactly ASI2.31\textquotesingle s native
chart, point marking, full endpoint
and cap product. Combining the directed
whole-reference equality with ASI2.44 gives
\[
 O_c-O_{r,-1}=(-71,4),\qquad
                    J_c-J_{r,-1}=(71,-4).
 \tag{ASI2.52}
\]
The correction is made once on the whole
reference after all its components and
its own cap have been included.

\par

\medskip\noindent\textbf{ASI2.6. All four parameter edges, retained winding and the source word}

\par\nobreak

Use the independent source cut \(\tau\) and
relative scalar-root coordinate \(\rho\),
based at the actual \(r_{\rm base}\).
The source, point and root family has
\(A_\rho=e^{2\pi i\rho}A_{\rm base}\),
\(B_\rho=e^{6\pi i\rho}B_{\rm base}\).
On its original sector, the complete
cochain is the same ASI2.29 with the
corresponding \(\log r_\lambda\), and
its retained integer shifts. Equivalently
its varying part is \(\rho\ell_\infty\delta\).
For every integer root turn the full
ordinary/cap forward arrows are
\(q^j,h_\infty^j\), with reciprocal
endpoint-to-initial arrows and the
actual \(B^j\) reindex at the selected zero.
At four this is exactly ASI2.30.

Put
\[
 G=\operatorname{Log}(P^{-3}),\quad
 J=\operatorname{Log}(\lambda_{\rm base}\eta_2^6t^{-3}),
 \quad
 A_\rho=e^{\rho G},\quad
 D_\rho=e^{\rho J+3\pi i\rho(\rho-1)}.
 \tag{ASI2.53}
\]
Here \(A_\rho,D_\rho\) denote the two point-coordinate
factors in this display; the root deck action
above remains the specified three-dimensional
cap map. Their types distinguish the two
uses. The native current-to-reference arrow is
\[
 C_\rho^N:(Z,a,c)\mapsto
      (A_\rho Z,(A_\rho^2D_\rho)^{-1}a,
                                  A_\rho D_\rho c).
 \tag{ASI2.54}
\]
It preserves the entire \(q\); its inverse
recovers that \(q\) and divides by the same
full functions. With
\(\eta_\rho=e^{6\pi i\rho}e^J\), subtraction
of the complete logarithms gives
\[
 C_{\rho+1}^N\,S(g^{-1},\eta_\rho^{-1})=C_\rho^N,
 \quad C_0^N=1,\quad C_4^N=S(g^4,e^{4J}),
 \tag{ASI2.55}
\]
where \(S(a,d)=(aZ,(a^2d)^{-1}a_N,ad\,c_N)\)
with \(a_N,c_N\) its independent normal
coordinates. This is the inverse of the
actual forward endpoint at the same
scalar-root endpoint. At the source-cut
edges, both \(G,J\) increase by \(-6\pi i\).
The native mismatch factors are therefore
\[
          (A,\beta)=(e^{-6\pi i\rho},e^{-12\pi i\rho}).
 \tag{ASI2.56}
\]
They are retained for every intermediate
integer corner, including the full
constant and based-unit values at the
initial source edge.

Writing \(w=e^{\pi i\rho}\), the full
held-Cox matrix identity is
\[
 T_{w^{-6}Z}\operatorname{diag}(w^{18},w^{12})
       =\operatorname{diag}(w^{18},w^{12})T_Z.
 \tag{ASI2.57}
\]
In its right-complex quaternion coefficient
this is exactly
\[
                         q_H\mapsto w^3q_Hw^{15}.
 \tag{ASI2.58}
\]
Both half-phase endpoints have sign
\((-1)^\rho\) at every integer corner.
The left adjoint is the actual rotation
on the fixed sphere and normal orbit
sphere, and extends linearly over their
specified three-balls. The full right
factor is its principal Hopf-line lift.
Thus this is an actual cap boundary map,
with closed paired Spin endpoints.
For arbitrary complex moduli the complete
matrix ASI2.36 and its positive Hermitian
ball map provide the same full extension,
without replacing the moduli by phases.

On the coupled diagonal \(\rho=4\tau\),
the original DZ product and ASI2.54
differ by the directed physical scalar
from NS.37:
\[
 F(\tau)=e^{3\pi i(16\tau^2+4\tau)}
       =e^{48\pi i\tau(\tau-1)}e^{60\pi i\tau},
 \quad(a,c)\mapsto(F^{-1}a,Fc).
 \tag{ASI2.59}
\]
The first displayed logarithm vanishes at
both endpoints and contracts by scaling
that logarithm. The second is the actual
positive physical loop of winding thirty.
Its held-Hopf factor is the inverse right
phase; it acts by the same global physical
circle on every original point, root,
Cartier and cap chart. Its degree-three
pullback is zero because it depends only
on \(\tau\). It commutes with the full
source and finite-deck maps, extends on
the product Hopf cap, and has paired
Spin endpoint \(+1\). The nonzero winding
is retained in this map. Replacing the
unflattened \(\tau\) by the NC endpoint-flat
cutoff adds only the displayed full
zero-endpoint logarithm differences;
the cutoff vertical flow transports all
normal jets and leaves ASI2.30 fixed.

The whole scalar reference at \(j=-1\)
uses the same native conjugation ASI2.54:
in RFB.19 its \(D_j=D A^{-j}\) is \(DA\).
Its source-cut quaternion is therefore
exactly ASI2.58. This is the actual
same-fan product obtained after RFB,
including its cap and original point-line
coefficient. All four edges of the
source/root parameter rectangle are now
specified by ASI2.29--30 and ASI2.53--59.

The coupled word is \((1,4)\), the
fixed-source scalar word is \((0,4)\),
and the fixed-scalar source word is
\((1,0)\). Pull back the displayed
family by the endpoint-fixed polygonal
homotopy from the diagonal to the source
edge followed by the scalar edge.
On each corner the full functional
equation ASI2.55 and the finite-cap
square ASI2.30 agree. The Hermitian
ball and full normal matrices extend
those maps over the selected caps,
with the retained winding ASI2.59.
Thickening these intervals gives the
actual pair-of-pants bordism of the
completed six-manifolds and their
physical line. Connections extend by
the same boundary derivatives.
Stokes on this bordism is additivity
of the two integer characteristic
numbers. Consequently ASI2.52 is
exactly the Jin physical pair of the
fixed-scalar source word, proving
the first equality in ASI2.3.

\par

\medskip\noindent\textbf{ASI2.7. The oriented sphere, native pair and marked generator}

\par\nobreak

At $r_{\rm src}=0$ the source is literally the E6-B source of
\cref{night:P3F.41}.  Use its sector $(a,b,m,n)=(2,3,3,1)$,
with the present finite character and entire annular reference, and
instantiate its selected Hom section by the exact coefficient EX.2.
Equivalently, in \cref{night:P3F.46} take
$(m_6,n_6;m_2,n_2)=(3,1;1,0)$, whose determinant is one.
At the second site F46.5 is the identity root
$(z,X,Y)=(t,x,y)$, its norm-zero map is the identity, and
F46.9 gives $H_2=B_\lambda$.  Thus this is the original ordinary
$III$ piece.  These source, line and cap maps give the initial
oriented sphere.

The $(2,3;3)$ row of \cref{night:P3K.96} supplies the full family
through that initial unselected $III$ value.  Its EX.1 is exactly
the six-factor multiplier ASI2.5, and its EX.2 is the same
$qB_{2,3}/t^3$.  Its cubic map is
$q_\infty=z^3$, $x=z^{-2}X$, $y=z^{-3}Y$, $t=z^{-3}$;
each original $f_2$ has weight $z^{-1}$, so the six factors
and selected-height factor give exactly
$B_{\rm sel}=(1-t_cz^3)B'_{2,3}$.  EX.14--19 retain the point/O
Cartier correction, finite character and entire reference cochain.
The shared source and translation charts used there are K94 EU.0,
with the identity on $(x,y,t,r_{\rm src})$; the present line maps
are the displayed $M_2$ maps of EX.1--2 and EX.14--19.

Continue over the actual real path
$0\le r_{\rm src}\le1/4-s_*$, with $0<s_*<1/12$, using
\[
 \rho=-\sqrt{1-3r_{\rm src}},\qquad
 t_c=\frac{(\rho-1)^2(2\rho+1)}{27},\qquad
 t_o=\frac{(\rho+1)^2(1-2\rho)}{27}.
\]
Here $t_c<0$, $t_o>0$ away from the initial endpoint, and
$t_c-t_o=4\rho^3/27\ne0$.  Thus the only source collision on
this path is the initial one just treated.  The height remains
$q=t-t_c$, the original sections remain marked, and the complete
SI2 source, Miller, toric and cubic-cap maps continue along it.

The fixed scalar satisfies a uniform bound on the whole path.
Choose a fixed small $\delta>0$ in the selected-collision chart
range.  The compact initial portion
$[0,1/4-\delta]$ has a uniform scalar bound $\varepsilon>0$;
choose $s_*$ sufficiently small that
$|\lambda_0|s_*^{12}<\varepsilon$.  On the remaining portion put
$\sigma=1/4-r_{\rm src}\ge s_*$.  Then
$|\lambda_{\rm base}|/\sigma^{12}
=|\lambda_0|(s_*/\sigma)^{12}\le|\lambda_0|$.
Decreasing the scalar increases $D_B$ in ASI2.8, so the exact
ASI2.10--11 bounds still hold.  The cubic root remains the fixed
$r_0s_*^4$, with its prescribed character and reference cochain.

Use the actual local source/cap parameter lifts and a relative
partition agreeing with the full derivative on protected collars;
average over the physical circle and finite deck group.  Their
flow gives the source map.  In the original Cartier frames its
line lift is the IL.29 linear ODE
$\dot w_i=-a_iw_i$, with
$a_i=-\sum_k\chi_kX(g_{ik})/g_{ik}$.
The backward source flow and inverse linear ODE give its inverse.
The original Hopf tube and the primitive F41/F35 boundary maps
travel through this same source/line/cap family.  Their Gysin and
van Kampen maps therefore retain the oriented sphere topology on
the stated selected fan.

The physical line and held quaternion
are the ones in ASI2.35--39. Thus
K68\textquotesingle s explicit physical-line inverse
clutch identifies the class of the
return with
\[
        (4k_M-\ell_M)/24=(284+4)/24=12.
 \tag{ASI2.60}
\]
The actual fixed-\(C_0\) free-coordinate
map in that same product is
\(k_M=\mu+\nu,\ell_M=4\mu\).
It is obtained by the same directed
Hopf transfer and negative-normal
conjugation, with both current cap
products fixed in ASI2.39. Hence
the source native pair is
\[
                         (\mu,\nu)=(-1,72).
 \tag{ASI2.61}
\]
This division by four is the stated
integral Hopf coordinate formula;
no division of a homotopy-sphere
class is used.

Finally the original marked-node
cover has \(v=w^2-1\),
\(D=v-2\), \(s=v(v+4)/(4D^2)\).
At \(w=-1\), \(Z_T=Z_w\), and
\(ds/dw\ne0\); the positive small
\(w=-1\) circle is exactly one
positive original \(s\)-meridian.
The full point/source map at the
other mark first sends
\((Z,V_{\rm pt})\mapsto(Z^{-1},P^2V_{\rm pt})\),
then uses the actual target period
\((Q,P^{-2})\), giving
\((Q/Z,V_{\rm pt})\).
Its reciprocal and all \(r_2\)
Cartier factors are SI2.21.
Thus the result transports to the
positive \(w=+1\) loop on the
literally reflected fan \(JT_{-1}\).
The original source and good cubic
cap depend on \(w^2\), and their
finite character and annular
cochain are unchanged by that
specified source identity. The
fan marking in ASI2.2 is retained
through the entire statement.


% End expanded verified cp36 input

\par\bigskip\noindent\textbf{P3K.111 supporting proof 20: An explicit auxiliary atlas and its complete physical contribution}\par\nobreak

The following construction uses the original source, point line
\(M_2=\mathcal O([2G]-[O])\), selected height, translation by \(-3G\),
scalar \(\lambda_0s^{12}\), root \(r_0s^4\), and the full finite-cap and
annular reference maps of SI2.  It gives another complete toric subdivision
of the three original periodic cell models, with an explicit identification
of its outer source, physical line and caps.  Its complete inner contribution
is computed below.

\medskip\noindent\textbf{A. The original cells and the actual common maps.}
Use \((X_c,Z_c,P)=(V_{\rm reg}/Z,Z,P)\).  Its full periods are
\[
 L_c=(P^{23}p^{-1}f_0^2X_c,P^3p^3Z_c,P),\qquad
 T_c=(P^{11}p^{-1}f_0X_c,P^3p^2Z_c,P).
\]
For a word \((i,j)\), put \(A=23i+11j\), \(B=3i+3j\),
\(u=2i+j\).  The valuation shift is \((Ah,Bh,0)\); the full p and
\(f_0\) columns are \((-i-j,3i+2j,0)^t\) and \((u,0,0)^t\).
In this lattice the three original height-one cells are
\[
\begin{split}
H&=((0,0),(-12,0),(-23,-3),(-11,-3)),\\
H_0&=((0,0),(5,0),(13,-3),(19,-6),(14,-6),(6,-3)),\\
H_1&=((0,0),(1,-3),(13,-3),(12,0)).
\end{split}
\]
They have area thirty-six.  The second and third are the literal images of
the zero hexagon in SI2.44 and the one parallelogram with period columns
\((3,1)^t,(0,12)^t\) in SI2.46 under
\((X_c,Z_c)=(1/(X_0Y_0),Y_0)\) and
\((X_c,Z_c)=(Y_1,1/X_1)\), respectively.

Intersect all translates of \(H_0,H_1\) meeting H.  The two coordinates of
\(M^{-1}(x,y)\), with
\(M=\left(\begin{smallmatrix}23&11\\3&3\end{smallmatrix}\right)\),
give exact finite translate bounds
by their vertex extrema.  Rational half-plane clipping gives six polygons
of disjoint interiors and total area thirty-six.  Insert their vertices and
all period copies in every incident edge, and use ordered ears.  The eight
initial simplicial cones are then regularized by integral nonzero quotient
representatives in their half-open parallelepipeds.  Every primitive ray is
inserted in its entire incident star and every period copy in H.  Each
positive barycentric coefficient is strictly less than one, so the
determinant decreases on every affected cone.  The full execution has 959
stellar orbits and 1,926 regular cones, of maximum primitive height nineteen.
Its exact area is thirty-six.  Its complete boundary edge sets agree under
\[
(a,b,h)\mapsto(a+11h,b+3h,h),\qquad
(a,b,h)\mapsto(a+12h,b,h),
\]
with the full \(T_c\) and \(L_cT_c^{-1}\) coefficients above.

Each child cone \(\tau\subset\sigma\) has the actual affine parent map
\[
 \mathbb C[\sigma^\vee\cap M]\longrightarrow
 \mathbb C[\tau^\vee\cap M],\qquad\chi^m\longmapsto\chi^m.
\]
All exponents are integral and nonnegative.  Equal character sums preserve
every binomial relation; character localization gives the inverse on the
original dense torus.  The saved owner word supplies all p and \(f_0\)
powers, with reciprocal inverse.  The complete subdivision covers each
parent support, so the valuative criterion gives its proper toric map.
These are maps to the original convex periodic cell models given by the
full SI2.6 periods and the raw supports ASI2.12--14.

\medskip\noindent\textbf{B. The three complete endpoint atlases.}
At zero the full native-to-common arrow is
\[
 X_c=P^Ap^{-i-j}f_0^u/(X_0Y_0),\qquad
 Z_c=P^Bp^{3i+2j}Y_0,
\]
with inverse
\[
 X_0=P^{A+B}p^{2i+j}f_0^u/(X_cZ_c),\qquad
 Y_0=P^{-B}p^{-(3i+2j)}Z_c.
\]
Join the complete native upper fan to the original lower rays
\(A,B_1,B_2,B_3\) of NS.3.  Join A to every upper cone.  On the four
remaining upper boundary sides, given in order by
\(3x-5y+15r=0\), \(x-y+7r=0\), \(y+6r=0\) and
\(3x-5y-6r=0\), join each complete interval to \((A,B_1)\) on the
first two and to \((A,B_3)\) on the last two.  Add
\((A,B_1,U_3,B_2)\) and \((A,B_3,U_3,B_2)\), with
\(U_3=(-13,-6,1,0)\).  These are the cones over the facets of the
convex raw slice not containing A.  Resolve every singular quotient
in the original lattice.  Subdivide every edge crossing
\(y+3r+z=0\) by its ordinary two-coordinate blowup, always in the current
whole fan.  The 8,394 source-cut stars give 48,700 regular cones.  Completing
all period ray orbits by ordinary face blowups takes 138 additional stars;
the complete period-facet comparison then has no unmatched triangle.
There are 54,121 final regular four-cones.  Their generic boundary is
exactly the 1,926-cone common fan, and their two original lower facets
remain.  All 302 paired period facets carry the complete maps
\[
L_0=(P^{10}pf_0X_0,P^3Y_0,P,p),\qquad
T_0=(P^8pf_0X_0,P^6pY_0,P,p),
\]
their actual composite \(L_0T_0^{-1}\), and all full reciprocals.
Each cone retains one source band, with both nonnegative rows
\(y+3jr+jz\), \(-y+3(1-j)r+(1-j)z\).  Their sum is \(3r+z\).
The exact whole slice volume is \(29/2\).

At one the complete arrow is
\[
 X_c=P^Av^{12u}p^{-i-j}c_{12}^uY_1,\qquad
 Z_c=P^Bp^{3i+2j}/X_1,
\]
with inverse obtained by the displayed reciprocal factors.  The native
top-edge chain is exactly
\(U(0,1),U(1,1)\), where
\(U(a,b)=(3a,a+12b,b,0)\).  Join every upper cone to \(E_0\), and this
edge to every \(E_l,E_{l+1}\), \(0\leq l<12\).  Resolving all twelve
determinant-three lower cones gives 2,022 regular four-cones.  Every full
L and T facet and both p and \(c_{12}\) columns agree with their
reciprocals.  The complete boundary has 1,926 generic facets and twenty-five
period pairs; all other facets are paired internally.  Its original source
rows are \(3r-x,x\), with the full p coefficient in \(pP^3/X_1\),
and its whole slice volume is eighteen.

At infinity use the literal map
\[
(X_c,Z_c,P,p)=(X,Z,Ru,u^{-1}),\qquad f_0=u^{-12}f_\infty.
\]
Join the upper cones to \(A_0\), and the twelve intervals on each of the
two specified lower boundary edges of ASI2.14 to \(A_0,A_1\).
The full period arrows are
\[
L_\infty=(R^{23}f_\infty^2X,R^3Z),\qquad
T_\infty=(R^{11}f_\infty X,R^3uZ),
\]
and their complete composite and inverse arrows.  The raw determinants
are one and eleven.  The 120 complete quotient stars give 2,381 regular
four-cones, of whole slice volume \(95/6\).  Every original L, T and L/T
facet is paired, with the full \(f_\infty\) factor and inverse; there are
twenty-five period pairs.  Both source rows \(y+3r+2z,-y\) are
nonnegative, and their sum is \(3r+2z\).

All additional corner rays in these constructions have both base heights
positive.  At zero the entire P-invertible subfan is the two original lower
triangles; at one it is exactly the twelve-edge chain on
\(E_0,\ldots,E_{12}\); at infinity the R-invertible subfan is exactly
\(\langle A_0,A_1\rangle\).  The identity-character maps therefore induce
literal isomorphisms on these localized coordinate rings, with the same
full period functions.  Since \(t=Pu_t\), \(s=P(p-1)/u_s\), and at
infinity \(t=Ru\,u_t\), \(s=R(1-u)/u_s\), every added divisor maps to
the original collision \((s,t)=(0,0)\).  Thus the new whole model has the
same original punctured source family and outer character maps.

\medskip\noindent\textbf{C. The unchanged exterior modules and the second action.}
The actual common-band source and point arrow is
\[
x=pP^3Z_c,\quad y=Z_c^{-1},\quad xy=P^2R,
\qquad v_2=\eta P^2X_cZ_c/f_{P^2}(x).
\]
All full endpoint substitutions above are used in this equation.  The
other retained source band has the full formula
\(x=Q^jZ_c\), \(y=Q^{1-j}/Z_c\),
\(v_2=\eta P^{2j}X_cZ_c/f_{P^2}(x)\); its actual integer j and all
owner period coefficients remain in the following maps.  The
proper point graph has the literal ideals \((x,P^2)\) and charts
\[
x=P^2a,\quad ay=R;\qquad xc=P^2,\quad y=Rc.
\]
Its second chart has the two resolutions I2P.2.  The full endpoint fans
do not all map to this proper graph: the rows
\(\operatorname{ord}x-2\operatorname{ord}P\) have both signs in
8,723 zero cones, 229 one cones and 184 infinity cones.  We do not pull
back and resolve that graph inside the counted atlas.  Instead the
regular periodic atlas constructed in A--B is used as the filling itself.
Its whole circle action and its regular tangent bundle are already
defined there.  Its graph identification with the original family is
required on the unchanged exterior and the specified lower source stars.

Every pure-lower face of the localized exterior maps to the resolved
graph.  Where \(\operatorname{ord}x\geq2\operatorname{ord}P\), the
map is \(a=x/P^2\).  Where
\(\operatorname{ord}x\leq\operatorname{ord}P\), its resolved coordinates
are \((u,w)=(x,P/x)\); where
\(\operatorname{ord}P\leq\operatorname{ord}x\leq2\operatorname{ord}P\),
they are \((u,w)=(x/P,P^2/x)\).  Their original source maps are respectively
\[
(P,x,c,y)=(uw,u,uw^2,Ruw^2),\qquad
(P,x,c,y)=(uw,u^2w,w,Rw).
\]
All these rows are nonnegative on the indicated actual faces.  At
infinity the entire localized lower cone uses the second resolved chart;
its full coordinates are \(u=R^2u_{\rm base}Z_c\),
\(w=R^{-1}Z_c^{-1}\), with \(P=Ru_{\rm base}\).  All twelve one-end
lower stars also satisfy these three chart inequalities on each of their
three complete charts.  The saved reciprocal maps preserve the original
point graph and source functions on the entire regions where they are
used below.

Glue this filling to the original exterior by the identity-character
isomorphisms of B and these graph arrows.  On their overlap every original
I2P coefficient and its inverse is the same full function.  The point/O
Cartier modules are consequently identified by the identity of their
specified generators and coefficients.  This construction supplies a
smooth filling of the original marked exterior without imposing a
source-graph projection on its new compact interior.

The commuting compact torus action, for \((\xi,\zeta)\in S^1\times S^1\),
\((X_c,Z_c)\mapsto(\xi X_c,\zeta Z_c)\), with P,p fixed, preserves the
ideal by its invertible generator matrix \(\operatorname{diag}(\zeta,1)\).
Its source maps are \(x'=\zeta x,y'=\zeta^{-1}y\).  On both resolved
node charts they are \(u'=\zeta u,w'=\zeta^{-1}w\).
Put \(b=P^2\), and write E and \(E'\) for the entire I2P unit at
\((x,y)\) and \((\zeta x,\zeta^{-1}y)\).  The exact point-module arrows are
\[
a'=\zeta a,\quad A'=\xi\zeta(E/E')A;\qquad
c'=\zeta^{-1}c,\quad C'=\xi(E/E')C.
\]
They preserve \(C=-cA\).  On the O module they are
\[
C'_O=\xi\zeta\frac{x-b}{\zeta x-b}\frac{E_O}{E'_O}C_O.
\]
Use the invariant small node domains and an invariant O annulus
\(1/2<|x|<2\), after shrinking the original base so \(|P^2|<1/4\).
Both \(x-P^2\) and \(\zeta x-P^2\) are nonzero throughout that annulus;
the complete \(E_O\) products and reciprocals are units there.  The
node \(E\) products remain units on the invariant node domains.
Thus these are invertible module arrows on the entire stated domains,
including the original point and O sections.  They follow directly from
I2P's factorizations of \(f_b\), and retain the additional \(P^{-2}\)
at O.  They extend through \(P=0\); the displayed E ratios are units.
All inverse arrows use the reciprocal at the recovered source point.
Multiplication cancels the middle E factor, proving the full action law.
In the genuine coefficient the same action is
\[
v'_2=\xi\zeta\frac{f_b(x)}{f_b(\zeta x)}v_2.
\]
For \(\zeta=1\) it is exactly \(v'_2=\xi v_2\), so this comparison
retains the original physical nonzero-vector line.

The native physical and second cocharacters are respectively
\[
\begin{array}{c|cc}
&\delta&a\\\hline
0&(-1,0,0,0)&(-1,1,0,0)\\
1&(0,1,0,0)&(-1,0,0,0)\\
\infty&(1,0,0,0)&(0,1,0,0).
\end{array}
\]
Every full period matrix fixes both.  The full units depend on the fixed
base parameters, so the two actions commute with all period maps.  This
proves the torus action on the actual regular atlas; on its unchanged
exterior and twelve lower source stars it is the displayed action on the
original point modules.  The action on the compact interior is the
literal toric action, so no unperformed graph resolution enters its
fixed-face or tangent-bundle calculation.

\medskip\noindent\textbf{D. Complete physical ownership and compact localization.}
For every full cone first compute \(B^{-1}\delta\).  Its nonzero entries
give the physical normal face.  Identify these faces through the complete
internal and period maps before using the second action.  The resulting
generic physical components number 1,664.  Their three endpoint stars are
the actual full discs or saturated half-plane fans between the two marked
section rays.  The compact corner inventories have 33,847 components at
zero, sixty at one and 337 at infinity.  Each point has its full chart,
each curve its opposite primitive tangent rays, and each surface its full
angularly ordered complete quotient fan.  These give 35,908 compact
physical components.  The remaining thirteen physical owners are exactly
the single moving orbit identified by T0 and the twelve named one-end
lower faces.  Their source stabilizers and full unit maps are retained.

Each compact physical component F is stable under the second circle:
the commuting connected torus acts continuously on the finite component
set, hence fixes every component.  Retain its actual physical normal
bundle E.  Every weight \(w_i\) on E is nonzero.  Therefore
\[
\int_F\frac{u^3}{e_{T^2}(E)},\qquad
\int_F\frac{u[p_1(TF)+p_{1,T^2}(E)]}{e_{T^2}(E)}
\]
are regular at the second parameter \(v=0\), after u is inverted.
Localize the second action inside this F.  The actual tangent and normal
sequences give
\(N_{H/X}=N_{H/F}+E|_H\) at each second-fixed piece H.  Thus its full
Euler denominator is the product of the second tangent Euler factor and
the retained physical normal factor.  All v poles cancel in the sum for
this particular F.  The original outward-meridian-one class is
\(u=c_1(U)\), with \(U=M^{-1}\); K68's actual Hermitian dual keeps
the quotient orientation and gives \(c_1(M)=-u\).  The physical pair
is consequently minus the \(u^{-1}\) coefficient of the displayed
two integrals at \(v=0\).  This fixes the sign from the original line
and circle maps before any residue is evaluated.

All common upper rays have positive P height, and every endpoint ray has
positive total base height.  Two such rays cannot span the torus plane
of zero base heights.  The second-fixed pieces consequently have complex
dimension at most one.  For a point put
\(W_i=w_i+\varepsilon a_i\).  Its pair is
\[
-1/\prod_iW_i,\qquad-\sum_iW_i^2/\prod_iW_i.
\]
For a curve with its actual normal degrees \(d_i\), put
\(W=\prod_iW_i\), \(S=\sum_i d_i/W_i\),
\(Q=\sum_iW_i^2\), \(T=\sum_iW_id_i\).  Its pair is
\[
S/W,\qquad(QS-2T)/W.
\]
These are direct expansions of the preceding Euler factors.  Every
piece is assigned to its physical owner first.  The 67 second-fixed
pieces in the thirteen prescribed noncompact owners remain in their
separate framed calculations below.

There is one generic second-fixed curve for each ordered common cone N.
Its unique endpoint chart has transverse ray \(l_e\) of base multiplicity
one.  The full coefficient clutch is
\[
c_e=N^{-1}F_el_e,\qquad d=-c_0-c_1-c_\infty,
\]
where the full endpoint character matrices obtained from B are
\[
F_0=\begin{pmatrix}-1&-1&23i+11j&-i-j\\0&1&3i+3j&3i+2j\\0&0&1&0\end{pmatrix},
\]
\[
F_1=\begin{pmatrix}0&1&23i+11j&12(2i+j)\\-1&0&3i+3j&0\\0&0&1&0\end{pmatrix},
\qquad F_\infty=\begin{pmatrix}1&0&0&0\\0&1&0&0\\0&0&1&1\end{pmatrix}.
\]
The zero and one owner words determine their respective i,j.
Every generic column is checked as its actual
permutation column; all p, \(f_0,c_{12},\eta\) coefficients and their
inverses are retained.  Each corner curve uses its entire two-chart wall
relation.  The whole computation includes all 1,926 generic curves and
52,546 compact corner second-fixed pieces.  It verifies the cancellation
of every negative \(\varepsilon\) coefficient separately for every
compact physical owner.  A second evaluation at \(u=1+\varepsilon\),
\(v=2\varepsilon\), retaining the numerator \(u^3\) or \(up_1\),
gives the same value for every owner.  The complete compact sum is
\[
                         (k_M,\ell_M)_{\rm compact}=(895,139).
\]

\medskip\noindent\textbf{E. The actual moving cap and twelve source doubles.}
The moving physical owner has thirty-one full charts.  Align its second
native face by the entire T0 map.  Evaluating the actual NS.23 coefficient
characters in both native frames gives thirty compact source rays, the
two horizontal rays, and both full normal divisors.  Every adjacent
derivative agrees.  The complete source intersection matrix and a second
calculation by twenty-nine chronological ordinary source blowdowns give
\[
D_-^2=-309,\quad D_-D_+=293,\quad D_+^2=-305,
\qquad hD_-=-4,\quad hD_+=5.
\]
Every contracted ray is the sum of its current neighbors.  Its maps are
\((u,w)\mapsto(u,uw)\) and \((u,w)\mapsto(uw,w)\), with the actual
reciprocal localizations and normal coefficient powers.  The marked
endpoint is the native sphere times the original P disc.

Attach the original NC cap through the new full derivative
\(N^+_\theta=G_0\operatorname{diag}(\overline{e_-},e_+)\).
The two native chart matrices obey NC.18.  With
\(B_\theta=C_\theta^{-1}N_\theta\), the actual cap map satisfies
\(\widehat C_\theta B_\theta=N_\theta\) and the reversed inverse
identity.  NC.20's rounded \(D^3\times D^2\) bordism therefore carries
this specified cap-completed bundle to the closed replacement with the
same full derivative clutch.  Its ordered positive line classes are
\(x=h+D_-\), \(y=-h-D_+\).  The actual twenty-nine blowups give
\[
x^2=-317,\quad xy=-294,\quad y^2=-295,
\qquad(c_1^2,c_2,p_1(TF))=(-1200,-294,-87).
\]
All full base units retain their nonzero constants and the reciprocal
NC.14 based-logarithm extension.  The reference uses its original own
cap and has \(c_1^2=c_2=p_1=0\).  The physical Euler denominator in
ASI2.41 gives the actual moving-minus-own-reference pair
\[
                              (-906,-111).
\]

For each of the twelve lower labels the complete new source star is the
three-chart minimal I3 star, with actual degree vectors
\((-1,0,1)\), \((-1,1,0)\).  Its source and line maps are the original
IL.20--30: in particular
\[
H_{+,l}=\eta_2u_s^lP^{-l}X^{-1},\qquad
H_{-,l}=\eta_2^{-1}XP^{l+1}u_s^{-l-1},\qquad
H_{+,l}H_{-,l}=P/u_s.
\]
Their period factors retain
\(\gamma=pu_s^{12}c_{12}=P^2\lambda_0q\eta_2^6/t^3\), including every
\(\gamma^j\) of a source-cell representative.  The actual source flow,
line ODE, inverse and cap square of IL.28--31 transport the entire marked
source body to its fixed reference.  The collision relative pairings are
\((-2/3,-1/3,-2/3)\).  The reference cube frame IL.33--37 and its
reciprocal give \((-2/3,+2/3,-2/3)\).  The boundary torsion characters
agree by \(I(-1,-1)^t=(1,1)^t\), with
\(I=\left(\begin{smallmatrix}-2&1\\1&-2\end{smallmatrix}\right)\);
the retained fibre gauge has zero pairing
against both normal lines.  Thus every actual closed double has normal
pairings \((0,-1,0)\) and tangent \(p_1=0\), with physical pair
\((1,-2)\).  All twelve give \((12,-24)\).

\medskip\noindent\textbf{F. The outer physical bundle, finite square and directed conclusion.}
The localized identity-character maps in B identify the new outer family
with the original source and line family on the entire punctured base.
The graph maps in C identify that entire exterior in the actual point/O
modules.  The moving cap uses the full outer normal derivative of NS.23a;
its twenty-six ambient source-graph crossings introduce no additional
resolution in the filling.  Its interior normal bundles are the actual
divisor bundles computed in E, and its entire cap attachment is the
displayed NC derivative map.  The twelve lower source bodies use the
complete resolved graph charts explicitly checked in C and their full
IL source flows and line maps from E.
They preserve the whole analytic functions P,R, \(s,t,q=t-t_c(s)\),
\(u_s,u_t,c_q,\eta\), and hence the selected root \(r_0s^4\).
The ordinary infinity cap and its original graph ideals lie in this
unchanged exterior.  Its full invariant cube root, finite deck character,
integral annular cochain and every sector integer are therefore the
literal ASI2.26--29 maps.  Their exact forward square remains
\[
                   \Phi_{\rho+4}q^4=h_\infty^4\Phi_\rho,
                   \qquad h_\infty=1-t_c(s)z^3,
\]
with the same reciprocal arrows and actual selected-zero reindex.
The fixed-source RB reference is unchanged on \(s=s_*\ne0\), including
its fourteen-factor source maps and its own caps.

On the physical quotient transport tangent and line connections through
these actual derivative and coefficient maps.  The identical
interpolation on each opposite collar has its interval reversed, so
its complete transgression cancels pointwise.  All primitive physical
weights and finite slice covering factors occur in the preceding
Euler denominators.  The full fixed-face inventory exhausts the compact
owners and thirteen completed noncompact owners.  The original complex
source volumes and the held Cox cap maps retain ASI2.32 and ASI2.43's
circle-first orientation.  Stokes therefore gives
\[
 O_c-O_{r,1}=(895,139)+(12,-24)+(-906,-111)=(1,4).
\]
This equality uses the new atlas for every collision-side contribution
and the exact original fixed-source reference for every comparison.

The whole reference change RFB.30--44 has the same source and line maps:
\(J_j:U_j\to I_j^*U_0\otimes\mathcal T^{12j}\), with the prescribed
outer primitive \(-j(du+\alpha_qdv)\), actual cap trivialization and
endpoint-to-initial clutch.  Its nilpotent tangent collar has zero
\(p_1\) transgression; its full Euler collar is \(-36j^2\), and the
actual integer \(0\to1\) Cox comparison fixes the linear coefficient
thirty-six.  Thus the same entire positive reference change is
\(O_{r,-1}-O_{r,1}=(72,0)\).  The literal map
\([y,t]\mapsto[y,1-t]\), with its same physical-line lift, reverses
the parameter orientation.  Hence
\[
 O_c-O_{r,-1}=(-71,4),\qquad J_c-J_{r,-1}=(71,-4).
\]
The full source/root rectangle, its four finite-cap squares, the native
right phase fifteen and retained physical winding thirty are the
unchanged ASI2.53--59 maps.  Their actual pair-of-pants bordism identifies
this relative pair with the fixed-scalar source word.  In the held Cox
product, \(\ell_M=4\mu\), \(k_M=\mu+\nu\), giving
\((\mu,\nu)=(-1,72)\).  The original physical inverse clutch gives
\((4k_M-\ell_M)/24=12\).

\medskip\noindent\textbf{G. Exact construction records.}
The complete common and endpoint matrices, all stellar morphisms and
period units are in the corresponding
\path{code/current_audits/audit_sign/K111_new_regular_common_fan.json},
\path{code/current_audits/audit_sign/K111_new_zero_endpoint_v2.json},
\path{code/current_audits/audit_sign/K111_new_one_endpoint.json} and
\path{code/current_audits/audit_sign/K111_new_infinity_endpoint.json}.
The compact physical-owner and three-end clutch certificate is
\path{code/current_audits/audit_sign/K111_T2_compact_ledger.json}; it
contains the exact ordered common-cone indices, full endpoint matrices,
weights, curve degrees, owner assignments and every per-owner result.
Its program performs 983,447 exact checks.  The native moving certificate
\path{code/current_audits/audit_sign/K111_new_moving_cap.json} retains both
intersection routes and all thirty-one actual normal charts.  The full
point/O and punctured-boundary comparison is
\path{code/current_audits/audit_sign/K111_actual_boundary_comparison.json},
with 476,963 checks.  Their generating programs are in the adjacent
\path{code/} directory.  Each computation starts from the original cells,
characters and full period functions; the complete geometric comparison
is the one proved in A--F.
The separate program
\path{code/K111_compact_owner_independent.py} recomputes all compact
owner Laurent coefficients and both substitutions, including every
ordered three-end clutch, in 484,310 exact checks.  Its result is
\path{code/outputs/K111_compact_owner_independent.json}.
The separate native-divisor and chronological-blowup calculation
\path{code/K111_moving_cap_independent.py} performs 378 checks and
gives the same cap roots and moving pair; its result is
\path{code/outputs/K111_moving_cap_independent.json}.
These are independent implementations within the same model family.
The exact graph-lift scope is
\path{code/current_audits/audit_sign/K111_point_graph_lift_scope.json}.
Its 20,711 checks retain every failed ambient graph-containment test,
every actual exterior face map, and all thirty-six resolved lower-source
chart maps.  It records the filling construction's precise use of the
original graph without changing the counted atlas.
The complete same-atlas certificate
\path{code/current_audits/audit_sign/K111_same_atlas_link.json} has
286,678 checks.  It assigns every final endpoint cone exactly once to
the generic, compact-corner or separately capped ledger; recomputes
its actual physical-owner map and corner weights; matches all moving
ordered cones and full T0 alignments; and records the SHA256 identities
of every final fan and both independently recomputed output inputs.

\end{proof}
